
\Data\ID{2000017201}
\Updated{2024-03-03 09:03:26}
\Level{A}
\SubArea{Matrices and determinants}
\LANGen{
\Question{
Which of the given matrices is of order \(3\) and has an entry of \(2\) at the position \((3,2)\)?}
{\[ 
\left (\MATRIX{ 
1& 2 & 3\cr 
 4 & 3 & 4
\cr 
3 & 2 & 1 } \right )
\]}{\[ 
\left (\MATRIX{ 
1& 2 & 3\cr 
 4 & 3 & 2
\cr 
3 & 4 & 1 } \right )
\]}{\[ 
\left (\MATRIX{ 
1& 2 & 3\cr 
 4 & 3 & 2
\cr 
3 & 2 & 2 \cr
1 & 2 & 3} \right )
\]}{\[ 
\left (\MATRIX{ 
1& 2 \cr 
 3 & 4 
\cr 
3 & 2  } \right )
\]}{}{}}
\LANGpl{
\Question{
Która z podanych macierzy jest rzędu \(3\) i ma wartość \(2\) na pozycji \((3,2)\)?}
{\[ 
\left (\MATRIX{ 
1& 2 & 3\cr 
 4 & 3 & 4
\cr 
3 & 2 & 1 } \right )
\]}{\[ 
\left (\MATRIX{ 
1& 2 & 3\cr 
 4 & 3 & 2
\cr 
3 & 4 & 1 } \right )
\]}{\[ 
\left (\MATRIX{ 
1& 2 & 3\cr 
 4 & 3 & 2
\cr 
3 & 2 & 2 \cr
1 & 2 & 3} \right )
\]}{\[ 
\left (\MATRIX{ 
1& 2 \cr 
 3 & 4 
\cr 
3 & 2  } \right )
\]}{}{}}
\LANGcs{
\Question{
Která z uvedených matic je řádu \(3\) a současně má na pozici \((3,2)\) prvek \(2\)?}
{\[ 
\left (\MATRIX{ 
1& 2 & 3\cr 
 4 & 3 & 4
\cr 
3 & 2 & 1 } \right )
\]}{\[ 
\left (\MATRIX{ 
1& 2 & 3\cr 
 4 & 3 & 2
\cr 
3 & 4 & 1 } \right )
\]}{\[ 
\left (\MATRIX{ 
1& 2 & 3\cr 
 4 & 3 & 2
\cr 
3 & 2 & 2 \cr
1 & 2 & 3} \right )
\]}{\[ 
\left (\MATRIX{ 
1& 2 \cr 
 3 & 4 
\cr 
3 & 2  } \right )
\]}{}{}}
\LANGsk{
\Question{
Ktorá z uvedených matíc je rádu 3 a súčasne má na pozícii (3,2) prvok 2?}
{\[ 
\left (\MATRIX{ 
1& 2 & 3\cr 
 4 & 3 & 4
\cr 
3 & 2 & 1 } \right )
\]}{\[ 
\left (\MATRIX{ 
1& 2 & 3\cr 
 4 & 3 & 2
\cr 
3 & 4 & 1 } \right )
\]}{\[ 
\left (\MATRIX{ 
1& 2 & 3\cr 
 4 & 3 & 2
\cr 
3 & 2 & 2 \cr
1 & 2 & 3} \right )
\]}{\[ 
\left (\MATRIX{ 
1& 2 \cr 
 3 & 4 
\cr 
3 & 2  } \right )
\]}{}{}}
\LANGes{
\Question{
¿Cuál de las matrices dadas es de orden \(3\) y tiene la entrada \(2\) en la posición \((3,2)\)?}
{\[ 
\left (\MATRIX{ 
1& 2 & 3\cr 
 4 & 3 & 4
\cr 
3 & 2 & 1 } \right )
\]}{\[ 
\left (\MATRIX{ 
1& 2 & 3\cr 
 4 & 3 & 2
\cr 
3 & 4 & 1 } \right )
\]}{\[ 
\left (\MATRIX{ 
1& 2 & 3\cr 
 4 & 3 & 2
\cr 
3 & 2 & 2 \cr
1 & 2 & 3} \right )
\]}{\[ 
\left (\MATRIX{ 
1& 2 \cr 
 3 & 4 
\cr 
3 & 2  } \right )
\]}{}{}}
\End

\Data\ID{2000017202}
\Updated{2024-03-03 09:03:26}
\Level{A}
\SubArea{Matrices and determinants}
\LANGen{
\Question{
Which of the given matrices \(K\), \(L\), \(M\), and \(N\) have the same main diagonal?
\[ 
K=\left (\MATRIX{ 
1& 7 & 8\cr 
 4 & 2 & 9
\cr 
5 & 6 & 3 } \right ),\quad
L=\left (\MATRIX{ 
3& 9 & 8\cr 
 4 & 2 & 7
\cr 
5 & 6 & 1 } \right ),
\]
\[ 
M=\left (\MATRIX{ 
1& 7 & 9\cr 
 5 & 2 & 8
\cr 
4 & 6 & 3 } \right ),
\quad
N=\left (\MATRIX{ 
2& 7 & 8\cr 
 4 & 2 & 6
\cr 
5 & 7 & 5 } \right )
\]}
{\(K\) and \(M\)}{\(K\), \(L\) and \(M\)}{\(K\), \(L\) and \(N\)}{The main diagonals of any two matrices differ.}{}{}}
\LANGpl{
\Question{
Które z podanych macierzy \(K\), \(L\), \(M\) i \(N\) mają tę samą główną przekątną?
\[ 
K=\left (\MATRIX{ 
1& 7 & 8\cr 
 4 & 2 & 9
\cr 
5 & 6 & 3 } \right ),
\quad
L=\left (\MATRIX{ 
3& 9 & 8\cr 
 4 & 2 & 7
\cr 
5 & 6 & 1 } \right ),
\]
\[ 
M=\left (\MATRIX{ 
1& 7 & 9\cr 
 5 & 2 & 8
\cr 
4 & 6 & 3 } \right ),
\quad
N=\left (\MATRIX{ 
2& 7 & 8\cr 
 4 & 2 & 6
\cr 
5 & 7 & 5 } \right )
\]}
{\(K\) i \(M\)}{\(K\), \(L\) i \(M\)}{\(K\), \(L\) i \(N\)}{Główne przekątne dowolnych dwóch macierzy różnią się.}{}{}}
\LANGcs{
\Question{
Které z uvedených matic \(K\), \(L\), \(M\), a \(N\) mají tutéž hlavní diagonálu?
\[ 
K=\left (\MATRIX{ 
1& 7 & 8\cr 
 4 & 2 & 9
\cr 
5 & 6 & 3 } \right ),\quad
L=\left (\MATRIX{ 
3& 9 & 8\cr 
 4 & 2 & 7
\cr 
5 & 6 & 1 } \right ),
\]
\[ 
M=\left (\MATRIX{ 
1& 7 & 9\cr 
 5 & 2 & 8
\cr 
4 & 6 & 3 } \right ),
\quad
N=\left (\MATRIX{ 
2& 7 & 8\cr 
 4 & 2 & 6
\cr 
5 & 7 & 5 } \right )
\]}
{\(K\) a \(M\)}{\(K\), \(L\) a \(M\)}{\(K\), \(L\) a \(N\)}{Žádné dvě matice nemají stejnou hlavní diagonálu.}{}{}}
\LANGsk{
\Question{
Ktoré z uvedených matíc K, L, M, a N majú rovnakú hlavnú diagonálu?
\[ 
K=\left (\MATRIX{ 
1& 7 & 8\cr 
 4 & 2 & 9
\cr 
5 & 6 & 3 } \right ),
\quad
L=\left (\MATRIX{ 
3& 9 & 8\cr 
 4 & 2 & 7
\cr 
5 & 6 & 1 } \right ),
\]
\[ 
M=\left (\MATRIX{ 
1& 7 & 9\cr 
 5 & 2 & 8
\cr 
4 & 6 & 3 } \right ),
\quad
N=\left (\MATRIX{ 
2& 7 & 8\cr 
 4 & 2 & 6
\cr 
5 & 7 & 5 } \right )
\]}
{\(K\) a \(M\)}{\(K\), \(L\) a \(M\)}{\(K\), \(L\) a \(N\)}{Žiadne dve matice nemajú rovnakú hlavnú diagonálu.}{}{}}
\LANGes{
\Question{
¿Cuáles de las matrices dadas \(K\), \(L\), \(M\), y \(N\) tienen la misma diagonal principal?
\[ 
K=\left (\MATRIX{ 
1& 7 & 8\cr 
 4 & 2 & 9
\cr 
5 & 6 & 3 } \right ),
\quad
L=\left (\MATRIX{ 
3& 9 & 8\cr 
 4 & 2 & 7
\cr 
5 & 6 & 1 } \right ),
\]
\[ 
M=\left (\MATRIX{ 
1& 7 & 9\cr 
 5 & 2 & 8
\cr 
4 & 6 & 3 } \right ),
\quad
N=\left (\MATRIX{ 
2& 7 & 8\cr 
 4 & 2 & 6
\cr 
5 & 7 & 5 } \right )
\]}
{\(K\) y \(M\)}{\(K\), \(L\) y \(M\)}{\(K\), \(L\) y \(N\)}{Las diagonales principales de dos matrices cualesquiera son diferentes.}{}{}}
\End

\Data\ID{2000017203}
\Updated{2024-03-03 09:03:25}
\Level{A}
\SubArea{Matrices and determinants}
\LANGen{
\Question{
Which of the given matrices have the same entry at the position \( (1,2)\)?
\[ 
K=\left (\MATRIX{ 
1& \sqrt2 & 3 & \sqrt5\cr 
\sqrt3& 2 & 1 & 5\cr 
4& 1 & 1& 0\cr 
} \right ),
\quad
L=\left (\MATRIX{ 
1& \sqrt2 & 3\cr 
 \sqrt3 & 2 & 1\cr 
4 & 1& 1\cr
\sqrt5 & 5& 0
 } \right ),
\]
\[ 
M=\left (\MATRIX{ 
\sqrt3& 2 & 1\cr 
 \sqrt3 & 4 & 0
} \right ),
\quad
N=\left (\MATRIX{ 
1& \sqrt3 & 4\cr 
 \sqrt3 & 2 & 1
\cr 
3 & 1 & 1 \cr
\sqrt5 & 5 & 0
 } \right )
\]}
{\(K\) and \(L\)}{\(K\), \(L\) and \(N\)}{\(K\), \(L\), \(M\), and \(N\)}{\(L\) and \(N\)}{}{}}
\LANGpl{
\Question{
Która z podanych macierzy ma tą samą wartość na pozycji  \( (1,2)\)?
\[ 
K=\left (\MATRIX{ 
1& \sqrt2 & 3 & \sqrt5\cr 
\sqrt3& 2 & 1 & 5\cr 
4& 1 & 1& 0\cr 
} \right ),
\quad
L=\left (\MATRIX{ 
1& \sqrt2 & 3\cr 
 \sqrt3 & 2 & 1\cr 
4 & 1& 1\cr
\sqrt5 & 5& 0
 } \right ),
\]
\[ 
M=\left (\MATRIX{ 
\sqrt3& 2 & 1\cr 
 \sqrt3 & 4 & 0
} \right ),
\quad
N=\left (\MATRIX{ 
1& \sqrt3 & 4\cr 
 \sqrt3 & 2 & 1
\cr 
3 & 1 & 1 \cr
\sqrt5 & 5 & 0
 } \right )
\]}
{\(K\) i \(L\)}{\(K\), \(L\) i \(N\)}{\(K\), \(L\), \(M\), i \(N\)}{\(L\) i \(N\)}{}{}}
\LANGcs{
\Question{
Ve kterých maticích je na pozici  \( (1,2)\) tentýž prvek?
\[ 
K=\left (\MATRIX{ 
1& \sqrt2 & 3 & \sqrt5\cr 
\sqrt3& 2 & 1 & 5\cr 
4& 1 & 1& 0\cr 
} \right ),
\quad
L=\left (\MATRIX{ 
1& \sqrt2 & 3\cr 
 \sqrt3 & 2 & 1\cr 
4 & 1& 1\cr
\sqrt5 & 5& 0
 } \right ),
\]
\[ 
M=\left (\MATRIX{ 
\sqrt3& 2 & 1\cr 
 \sqrt3 & 4 & 0
} \right ),
\quad
N=\left (\MATRIX{ 
1& \sqrt3 & 4\cr 
 \sqrt3 & 2 & 1
\cr 
3 & 1 & 1 \cr
\sqrt5 & 5 & 0
 } \right )
\]}
{\(K\) a \(L\)}{\(K\), \(L\) a \(N\)}{\(K\), \(L\), \(M\), a \(N\)}{\(L\) a \(N\)}{}{}}
\LANGsk{
\Question{
V ktorých maticiach je na pozícii (1,2) rovnaký prvok?
\[ 
K=\left (\MATRIX{ 
1& \sqrt2 & 3 & \sqrt5\cr 
\sqrt3& 2 & 1 & 5\cr 
4& 1 & 1& 0\cr 
} \right ),
\quad
L=\left (\MATRIX{ 
1& \sqrt2 & 3\cr 
 \sqrt3 & 2 & 1\cr 
4 & 1& 1\cr
\sqrt5 & 5& 0
 } \right ),
\]
\[ 
M=\left (\MATRIX{ 
\sqrt3& 2 & 1\cr 
 \sqrt3 & 4 & 0
} \right ),
\quad
N=\left (\MATRIX{ 
1& \sqrt3 & 4\cr 
 \sqrt3 & 2 & 1
\cr 
3 & 1 & 1 \cr
\sqrt5 & 5 & 0
 } \right )
\]}
{\(K\) a \(L\)}{\(K\), \(L\) a \(N\)}{\(K\), \(L\), \(M\), a \(N\)}{\(L\) a \(N\)}{}{}}
\LANGes{
\Question{
¿Cuáles de las matrices dadas tienen la misma entrada en la posición \( (1,2)\)?
\[ 
K=\left (\MATRIX{ 
1& \sqrt2 & 3 & \sqrt5\cr 
\sqrt3& 2 & 1 & 5\cr 
4& 1 & 1& 0\cr 
} \right ),
\quad
L=\left (\MATRIX{ 
1& \sqrt2 & 3\cr 
 \sqrt3 & 2 & 1\cr 
4 & 1& 1\cr
\sqrt5 & 5& 0
 } \right ),
\]
\[ 
M=\left (\MATRIX{ 
\sqrt3& 2 & 1\cr 
 \sqrt3 & 4 & 0
} \right ),
\quad
N=\left (\MATRIX{ 
1& \sqrt3 & 4\cr 
 \sqrt3 & 2 & 1
\cr 
3 & 1 & 1 \cr
\sqrt5 & 5 & 0
 } \right )
\]}
{\(K\) y \(L\)}{\(K\), \(L\) y \(N\)}{\(K\), \(L\), \(M\), y \(N\)}{\(L\) y \(N\)}{}{}}
\End

\Data\ID{2000017204}
\Updated{2024-03-03 09:03:25}
\Level{A}
\SubArea{Matrices and determinants}
\LANGen{
\Question{
Which of the given matrices is the matrix \((m_{i,j})\), where \(i=1, \dots, 3\) and \(j=1,~2\)?}
{\(
\left (\MATRIX{ 
8& 7\cr 
6 & 5\cr 
4 & 3\cr
 } \right )
\)}{\(
\left (\MATRIX{ 
8& 7 & 6\cr 
5 & 4 & 3\cr 
 } \right )
\)}{\(
\left (\MATRIX{ 
8& 7 & 6\cr 
5 & 4 & 3\cr 
2 & 1 & 0
 } \right )
\)}{\(
\left (\MATRIX{ 
8& 7 \cr 
6 & 5 \cr 
 } \right )
\)}{}{}}
\LANGpl{
\Question{
Która z podanych macierzy jest macierzą \((m_{i,j})\), gdzie \(i=1, \dots, 3\) oraz \(j=1,~2\)?}
{\(
\left (\MATRIX{ 
8& 7\cr 
6 & 5\cr 
4 & 3\cr
 } \right )
\)}{\(
\left (\MATRIX{ 
8& 7 & 6\cr 
5 & 4 & 3\cr 
 } \right )
\)}{\(
\left (\MATRIX{ 
8& 7 & 6\cr 
5 & 4 & 3\cr 
2 & 1 & 0
 } \right )
\)}{\(
\left (\MATRIX{ 
8& 7 \cr 
6 & 5 \cr 
 } \right )
\)}{}{}}
\LANGcs{
\Question{
Která z uvedených matic je maticí \((m_{i,j})\), kde \(i=1, \dots, 3\) a \(j=1,~2\)?}
{\(
\left (\MATRIX{ 
8& 7\cr 
6 & 5\cr 
4 & 3\cr
 } \right )
\)}{\(
\left (\MATRIX{ 
8& 7 & 6\cr 
5 & 4 & 3\cr 
 } \right )
\)}{\(
\left (\MATRIX{ 
8& 7 & 6\cr 
5 & 4 & 3\cr 
2 & 1 & 0
 } \right )
\)}{\(
\left (\MATRIX{ 
8& 7 \cr 
6 & 5 \cr 
 } \right )
\)}{}{}}
\LANGsk{
\Question{
Ktorá z uvedených matíc je matica \((m_{i,j})\), kde \(i=1, \dots, 3\) a \(j=1,~2\)?}
{\(
\left (\MATRIX{ 
8& 7\cr 
6 & 5\cr 
4 & 3\cr
 } \right )
\)}{\(
\left (\MATRIX{ 
8& 7 & 6\cr 
5 & 4 & 3\cr 
 } \right )
\)}{\(
\left (\MATRIX{ 
8& 7 & 6\cr 
5 & 4 & 3\cr 
2 & 1 & 0
 } \right )
\)}{\(
\left (\MATRIX{ 
8& 7 \cr 
6 & 5 \cr 
 } \right )
\)}{}{}}
\LANGes{
\Question{
¿Cuál de las matrices dadas es la matriz \((m_{i,j})\), donde \(i=1, \dots, 3\) y \(j=1,~2\)?}
{\(
\left (\MATRIX{ 
8& 7\cr 
6 & 5\cr 
4 & 3\cr
 } \right )
\)}{\(
\left (\MATRIX{ 
8& 7 & 6\cr 
5 & 4 & 3\cr 
 } \right )
\)}{\(
\left (\MATRIX{ 
8& 7 & 6\cr 
5 & 4 & 3\cr 
2 & 1 & 0
 } \right )
\)}{\(
\left (\MATRIX{ 
8& 7 \cr 
6 & 5 \cr 
 } \right )
\)}{}{}}
\End

\Data\ID{2000017205}
\Updated{2024-03-03 09:03:25}
\Level{A}
\SubArea{Matrices and determinants}
\LANGen{
\Question{
Which of the given matrices is the matrix \((m_{i,j})\), where \(i=1, \dots, 3\) and \(j=1,\dots,3\) with \(m_{i,j}=i+j+1\)?}
{\(
\left (\MATRIX{ 
3& 4 & 5\cr 
4 & 5 & 6\cr 
5 & 6 & 7
 } \right )
\)}{\(
\left (\MATRIX{ 
2& 3 & 4\cr 
3 & 4 & 5\cr 
4 & 5 & 6
 } \right )
\)}{\(
\left (\MATRIX{ 
3& 4 & 5\cr 
3 & 4 & 5\cr 
4 & 5 & 6
 } \right )
\)}{\(
\left (\MATRIX{ 
2& 3 & 4\cr 
2 & 3 & 4\cr 
4 & 5 & 6
 } \right )
\)}{}{}}
\LANGpl{
\Question{
Która z podanych macierzy jest macierzą \((m_{i,j})\), gdzie \(i=1, \dots, 3\) oraz \(j=1,\dots,3\) z \(m_{i,j}=i+j+1\)?}
{\(
\left (\MATRIX{ 
3& 4 & 5\cr 
4 & 5 & 6\cr 
5 & 6 & 7
 } \right )
\)}{\(
\left (\MATRIX{ 
2& 3 & 4\cr 
3 & 4 & 5\cr 
4 & 5 & 6
 } \right )
\)}{\(
\left (\MATRIX{ 
3& 4 & 5\cr 
3 & 4 & 5\cr 
4 & 5 & 6
 } \right )
\)}{\(
\left (\MATRIX{ 
2& 3 & 4\cr 
2 & 3 & 4\cr 
4 & 5 & 6
 } \right )
\)}{}{}}
\LANGcs{
\Question{
Která z uvedených matic je maticí \((m_{i,j})\), kde \(i=1, \dots, 3\) a \(j=1,\dots,3\), je-li \(m_{i,j}=i+j+1\)?}
{\(
\left (\MATRIX{ 
3& 4 & 5\cr 
4 & 5 & 6\cr 
5 & 6 & 7
 } \right )
\)}{\(
\left (\MATRIX{ 
2& 3 & 4\cr 
3 & 4 & 5\cr 
4 & 5 & 6
 } \right )
\)}{\(
\left (\MATRIX{ 
3& 4 & 5\cr 
3 & 4 & 5\cr 
4 & 5 & 6
 } \right )
\)}{\(
\left (\MATRIX{ 
2& 3 & 4\cr 
2 & 3 & 4\cr 
4 & 5 & 6
 } \right )
\)}{}{}}
\LANGsk{
\Question{
Ktorá z uvedených matíc je matica \((m_{i,j})\), kde \(i=1, \dots, 3\) a \(j=1,\dots,3\), ak \(m_{i,j}=i+j+1\)?}
{\(
\left (\MATRIX{ 
3& 4 & 5\cr 
4 & 5 & 6\cr 
5 & 6 & 7
 } \right )
\)}{\(
\left (\MATRIX{ 
2& 3 & 4\cr 
3 & 4 & 5\cr 
4 & 5 & 6
 } \right )
\)}{\(
\left (\MATRIX{ 
3& 4 & 5\cr 
3 & 4 & 5\cr 
4 & 5 & 6
 } \right )
\)}{\(
\left (\MATRIX{ 
2& 3 & 4\cr 
2 & 3 & 4\cr 
4 & 5 & 6
 } \right )
\)}{}{}}
\LANGes{
\Question{
¿Cuál de las matrices dadas es la matriz \((m_{i,j})\), donde \(i=1, \dots, 3\) y \(j=1,\dots,3\) con \(m_{i,j}=i+j+1\)?}
{\(
\left (\MATRIX{ 
3& 4 & 5\cr 
4 & 5 & 6\cr 
5 & 6 & 7
 } \right )
\)}{\(
\left (\MATRIX{ 
2& 3 & 4\cr 
3 & 4 & 5\cr 
4 & 5 & 6
 } \right )
\)}{\(
\left (\MATRIX{ 
3& 4 & 5\cr 
3 & 4 & 5\cr 
4 & 5 & 6
 } \right )
\)}{\(
\left (\MATRIX{ 
2& 3 & 4\cr 
2 & 3 & 4\cr 
4 & 5 & 6
 } \right )
\)}{}{}}
\End

\Data\ID{2000017207}
\Updated{2024-03-03 09:03:25}
\Level{A}
\SubArea{Matrices and determinants}
\LANGen{
\Question{
For what values of \(a\) and \(b\) is the matrix
\[
\left (\MATRIX{ 
a+5& a+3 & 0\cr 
0 & a-2 & 0\cr 
-a-b & a+b & a+1
 } \right )
\]
a diagonal matrix?}
{\(a=-3\), \(b=3\)}{\(a=2\), \(b=-2\)}{\(a=3\), \(b=-3\)}{\(a=-3\), \(b=-3\)}{}{}}
\LANGpl{
\Question{
Dla jakich wartości \(a\) i \(b\) macierz
\[
\left (\MATRIX{ 
a+5& a+3 & 0\cr 
0 & a-2 & 0\cr 
-a-b & a+b & a+1
 } \right )
\]
to macierz diagonalna?}
{\(a=-3\), \(b=3\)}{\(a=2\), \(b=-2\)}{\(a=3\), \(b=-3\)}{\(a=-3\), \(b=-3\)}{}{}}
\LANGcs{
\Question{
Pro jaké hodnoty \(a\) a \(b\) je matice
\[
\left (\MATRIX{ 
a+5& a+3 & 0\cr 
0 & a-2 & 0\cr 
-a-b & a+b & a+1
 } \right )
\]
diagonální maticí?}
{\(a=-3\), \(b=3\)}{\(a=2\), \(b=-2\)}{\(a=3\), \(b=-3\)}{\(a=-3\), \(b=-3\)}{}{}}
\LANGsk{
\Question{
Pre aké hodnoty \(a\) a \(b\) je matica
\[
\left (\MATRIX{ 
a+5& a+3 & 0\cr 
0 & a-2 & 0\cr 
-a-b & a+b & a+1
 } \right )
\]
diagonálnou maticou?}
{\(a=-3\), \(b=3\)}{\(a=2\), \(b=-2\)}{\(a=3\), \(b=-3\)}{\(a=-3\), \(b=-3\)}{}{}}
\LANGes{
\Question{
Halla los valores de \(a\) y \(b\) para que la matriz sea una matriz diagonal.
\[
\left (\MATRIX{ 
a+5& a+3 & 0\cr 
0 & a-2 & 0\cr 
-a-b & a+b & a+1
 } \right )
\]}
{\(a=-3\), \(b=3\)}{\(a=2\), \(b=-2\)}{\(a=3\), \(b=-3\)}{\(a=-3\), \(b=-3\)}{}{}}
\End

\Data\ID{2000017401}
\Updated{2024-03-03 09:03:25}
\Level{A}
\SubArea{Matrices and determinants}
\LANGen{
\Question{
Find the product of the following matrices:
\[ 
A=\left (\MATRIX{ 
\frac12  & \frac34\cr 
-1 & \frac32 } \right ),~
B=\left (\MATRIX{ 
\frac32 & -\frac14 \cr 
0& \frac12 } \right )
\]}
{\( AB=\left (\MATRIX{ 
\frac34 & \frac14 \cr 
-\frac32& 1 } \right )
\)}{\( AB=\left (\MATRIX{ 
\frac34 & -\frac3{16} \cr 
0 & \frac34} \right )
\)}{\( AB=\left (\MATRIX{ 
\frac34 & \frac12 \cr 
-\frac32 & 1 } \right )
\)}{\( AB=\left (\MATRIX{ 
2 & \frac12 \cr 
-1 & 2} \right )
\)}{}{}}
\LANGpl{
\Question{
Oblicz iloczyn następujących macierzy:
\[ 
A=\left (\MATRIX{ 
\frac12  & \frac34\cr 
-1 & \frac32 } \right ),~
B=\left (\MATRIX{ 
\frac32 & -\frac14 \cr 
0& \frac12 } \right )
\]}
{\( AB=\left (\MATRIX{ 
\frac34 & \frac14 \cr 
-\frac32& 1 } \right )
\)}{\( AB=\left (\MATRIX{ 
\frac34 & -\frac3{16} \cr 
0 & \frac34} \right )
\)}{\( AB=\left (\MATRIX{ 
\frac34 & \frac12 \cr 
-\frac32 & 1 } \right )
\)}{\( AB=\left (\MATRIX{ 
2 & \frac12 \cr 
-1 & 2} \right )
\)}{}{}}
\LANGcs{
\Question{
Vypočítejte součin následujících matic:
\[ 
A=\left (\MATRIX{ 
\frac12  & \frac34\cr 
-1 & \frac32 } \right ),~
B=\left (\MATRIX{ 
\frac32 & -\frac14 \cr 
0& \frac12 } \right )
\]}
{\( AB=\left (\MATRIX{ 
\frac34 & \frac14 \cr 
-\frac32& 1 } \right )
\)}{\( AB=\left (\MATRIX{ 
\frac34 & -\frac3{16} \cr 
0 & \frac34} \right )
\)}{\( AB=\left (\MATRIX{ 
\frac34 & \frac12 \cr 
-\frac32 & 1 } \right )
\)}{\( AB=\left (\MATRIX{ 
2 & \frac12 \cr 
-1 & 2} \right )
\)}{}{}}
\LANGsk{
\Question{
Vypočítajte súčin nasledujúcich matíc:
\[ 
A=\left (\MATRIX{ 
\frac12  & \frac34\cr 
-1 & \frac32 } \right ),~
B=\left (\MATRIX{ 
\frac32 & -\frac14 \cr 
0& \frac12 } \right )
\]}
{\( AB=\left (\MATRIX{ 
\frac34 & \frac14 \cr 
-\frac32& 1 } \right )
\)}{\( AB=\left (\MATRIX{ 
\frac34 & -\frac3{16} \cr 
0 & \frac34} \right )
\)}{\( AB=\left (\MATRIX{ 
\frac34 & \frac12 \cr 
-\frac32 & 1 } \right )
\)}{\( AB=\left (\MATRIX{ 
2 & \frac12 \cr 
-1 & 2} \right )
\)}{}{}}
\LANGes{
\Question{
Calcula el producto de las siguientes matrices:
\[ 
A=\left (\MATRIX{ 
\frac12  & \frac34\cr 
-1 & \frac32 } \right ),~
B=\left (\MATRIX{ 
\frac32 & -\frac14 \cr 
0& \frac12 } \right )
\]}
{\( AB=\left (\MATRIX{ 
\frac34 & \frac14 \cr 
-\frac32& 1 } \right )
\)}{\( AB=\left (\MATRIX{ 
\frac34 & -\frac3{16} \cr 
0 & \frac34} \right )
\)}{\( AB=\left (\MATRIX{ 
\frac34 & \frac12 \cr 
-\frac32 & 1 } \right )
\)}{\( AB=\left (\MATRIX{ 
2 & \frac12 \cr 
-1 & 2} \right )
\)}{}{}}
\End

\Data\ID{2000017402}
\Updated{2024-03-03 09:03:25}
\Level{A}
\SubArea{Matrices and determinants}
\LANGen{
\Question{
Find the product:
\[ 
\left (\MATRIX{ 
1.2& 1 & 0\cr 
 0.3 & 0.1 & 1 \cr 
2 & 0 & 4} \right ) \cdot
\left (\MATRIX{ 
1& 1 & 0\cr 
 0 & 1 & 1 \cr 
1 & 1 & 1 } \right )
\]}
{\(
\left (\MATRIX{ 
1.2& 2.2 & 1\cr 
 1.3 & 1.4& 1.1 \cr 
6 & 6& 4 } \right )
\)}{\(
\left (\MATRIX{ 
1.2& 1 & 0\cr 
0& 0.1 & 1\cr 
2 & 0& 4 } \right )
\)}{\(
\left (\MATRIX{ 
3.2& 3.2 & 1\cr 
 1.4 & 1.4& 1.4 \cr 
6 & 7& 4 } \right )
\)}{\(
\left (\MATRIX{ 
1.2& 2.2 & 1\cr 
 1.1 & 1.4& 0 \cr 
2 & 4& 6 } \right )
\)}{}{}}
\LANGpl{
\Question{
Oblicz:
\[ 
\left (\MATRIX{ 
1{,}2& 1 & 0\cr 
 0{,}3 & 0{,}1 & 1 \cr 
2 & 0 & 4} \right ) \cdot
\left (\MATRIX{ 
1& 1 & 0\cr 
 0 & 1 & 1 \cr 
1 & 1 & 1 } \right )
\]}
{\(
\left (\MATRIX{ 
1{,}2& 2{,}2 & 1\cr 
 1{,}3 & 1{,}4& 1{,}1 \cr 
6 & 6& 4 } \right )
\)}{\(
\left (\MATRIX{ 
1{,}2& 1 & 0\cr 
0& 0{,}1 & 1\cr 
2 & 0& 4 } \right )
\)}{\(
\left (\MATRIX{ 
3{,}2& 3{,}2 & 1\cr 
 1{,}4 & 1{,}4& 1{,}4 \cr 
6 & 7& 4 } \right )
\)}{\(
\left (\MATRIX{ 
1{,}2& 2{,}2 & 1\cr 
 1{,}1 & 1{,}4& 0 \cr 
2 & 4& 6 } \right )
\)}{}{}}
\LANGcs{
\Question{
Vypočítejte součin:
\[ 
\left (\MATRIX{ 
1{,}2& 1 & 0\cr 
 0{,}3 & 0{,}1 & 1 \cr 
2 & 0 & 4} \right ) \cdot
\left (\MATRIX{ 
1& 1 & 0\cr 
 0 & 1 & 1 \cr 
1 & 1 & 1 } \right )
\]}
{\(
\left (\MATRIX{ 
1{,}2& 2{,}2 & 1\cr 
 1{,}3 & 1{,}4& 1{,}1 \cr 
6 & 6& 4 } \right )
\)}{\(
\left (\MATRIX{ 
1{,}2& 1 & 0\cr 
0& 0{,}1 & 1\cr 
2 & 0& 4 } \right )
\)}{\(
\left (\MATRIX{ 
3{,}2& 3{,}2 & 1\cr 
 1{,}4 & 1{,}4& 1{,}4 \cr 
6 & 7& 4 } \right )
\)}{\(
\left (\MATRIX{ 
1{,}2& 2{,}2 & 1\cr 
 1{,}1 & 1{,}4& 0 \cr 
2 & 4& 6 } \right )
\)}{}{}}
\LANGsk{
\Question{
Vypočítajte súčin:
\[ 
\left (\MATRIX{ 
1{,}2& 1 & 0\cr 
 0{,}3 & 0{,}1 & 1 \cr 
2 & 0 & 4} \right ) \cdot
\left (\MATRIX{ 
1& 1 & 0\cr 
 0 & 1 & 1 \cr 
1 & 1 & 1 } \right )
\]}
{\(
\left (\MATRIX{ 
1{,}2& 2{,}2 & 1\cr 
 1{,}3 & 1{,}4& 1{,}1 \cr 
6 & 6& 4 } \right )
\)}{\(
\left (\MATRIX{ 
1{,}2& 1 & 0\cr 
0& 0{,}1 & 1\cr 
2 & 0& 4 } \right )
\)}{\(
\left (\MATRIX{ 
3{,}2& 3{,}2 & 1\cr 
 1{,}4 & 1{,}4& 1{,}4 \cr 
6 & 7& 4 } \right )
\)}{\(
\left (\MATRIX{ 
1{,}2& 2{,}2 & 1\cr 
 1{,}1 & 1{,}4& 0 \cr 
2 & 4& 6 } \right )
\)}{}{}}
\LANGes{
\Question{
Calcula el producto:
\[ 
\left (\MATRIX{ 
1.2& 1 & 0\cr 
 0.3 & 0.1 & 1 \cr 
2 & 0 & 4} \right ) \cdot
\left (\MATRIX{ 
1& 1 & 0\cr 
 0 & 1 & 1 \cr 
1 & 1 & 1 } \right )
\]}
{\(
\left (\MATRIX{ 
1.2& 2.2 & 1\cr 
 1.3 & 1.4& 1.1 \cr 
6 & 6& 4 } \right )
\)}{\(
\left (\MATRIX{ 
1.2& 1 & 0\cr 
0& 0.1 & 1\cr 
2 & 0& 4 } \right )
\)}{\(
\left (\MATRIX{ 
3.2& 3.2 & 1\cr 
 1.4 & 1.4& 1.4 \cr 
6 & 7& 4 } \right )
\)}{\(
\left (\MATRIX{ 
1.2& 2.2 & 1\cr 
 1.1 & 1.4& 0 \cr 
2 & 4& 6 } \right )
\)}{}{}}
\End

\Data\ID{2000017403}
\Updated{2024-03-03 09:03:25}
\Level{A}
\SubArea{Matrices and determinants}
\LANGen{
\Question{
For which real numbers \(a\), \(b\), \(c\), \(d\) is the following equality true?
\[ 
2 \cdot \left (\MATRIX{ 
a& c\cr 
 b & d \cr 
} \right ) 
-
\left (\MATRIX{ 
2& 1 \cr 
3 & 5 \cr 
} \right ) 
=
\left (\MATRIX{ 
4& 7\cr 
-5 & -5\cr 
} \right ) 
\]}
{\(a=3\), \(b=-1\), \(c=4\), \(d=0\)}{\(a=1\), \(b=1\), \(c=4\), \(d=1\)}{\(a=3\), \(b=-4\), \(c=4\), \(d=0\)}{\(a=3\), \(b=-1\), \(c=4\), \(d=1\)}{}{}}
\LANGpl{
\Question{
Dla jakich liczb rzeczywistych \(a\), \(b\), \(c\), \(d\) podana równość jest prawdziwa?
\[ 
2 \cdot \left (\MATRIX{ 
a& c\cr 
 b & d \cr 
} \right ) 
-
\left (\MATRIX{ 
2& 1 \cr 
3 & 5 \cr 
} \right ) 
=
\left (\MATRIX{ 
4& 7\cr 
-5 & -5\cr 
} \right ) 
\]}
{\(a=3\), \(b=-1\), \(c=4\), \(d=0\)}{\(a=1\), \(b=1\), \(c=4\), \(d=1\)}{\(a=3\), \(b=-4\), \(c=4\), \(d=0\)}{\(a=3\), \(b=-1\), \(c=4\), \(d=1\)}{}{}}
\LANGcs{
\Question{
Pro která reálná čísla \(a\), \(b\), \(c\), \(d\) platí rovnost?
\[ 
2 \cdot \left (\MATRIX{ 
a& c\cr 
 b & d \cr 
} \right ) 
-
\left (\MATRIX{ 
2& 1 \cr 
3 & 5 \cr 
} \right ) 
=
\left (\MATRIX{ 
4& 7\cr 
-5 & -5\cr 
} \right ) 
\]}
{\(a=3\), \(b=-1\), \(c=4\), \(d=0\)}{\(a=1\), \(b=1\), \(c=4\), \(d=1\)}{\(a=3\), \(b=-4\), \(c=4\), \(d=0\)}{\(a=3\), \(b=-1\), \(c=4\), \(d=1\)}{}{}}
\LANGsk{
\Question{
Pre ktoré reálne čísla \(a\), \(b\), \(c\), \(d\) platí rovnosť?
\[ 
2 \cdot \left (\MATRIX{ 
a& c\cr 
 b & d \cr 
} \right ) 
-
\left (\MATRIX{ 
2& 1 \cr 
3 & 5 \cr 
} \right ) 
=
\left (\MATRIX{ 
4& 7\cr 
-5 & -5\cr 
} \right ) 
\]}
{\(a=3\), \(b=-1\), \(c=4\), \(d=0\)}{\(a=1\), \(b=1\), \(c=4\), \(d=1\)}{\(a=3\), \(b=-4\), \(c=4\), \(d=0\)}{\(a=3\), \(b=-1\), \(c=4\), \(d=1\)}{}{}}
\LANGes{
\Question{
¿Para qué números reales \(a\), \(b\), \(c\), \(d\) es la siguiente ecuación verdadera?
\[ 
2 \cdot \left (\MATRIX{ 
a& c\cr 
 b & d \cr 
} \right ) 
-
\left (\MATRIX{ 
2& 1 \cr 
3 & 5 \cr 
} \right ) 
=
\left (\MATRIX{ 
4& 7\cr 
-5 & -5\cr 
} \right ) 
\]}
{\(a=3\), \(b=-1\), \(c=4\), \(d=0\)}{\(a=1\), \(b=1\), \(c=4\), \(d=1\)}{\(a=3\), \(b=-4\), \(c=4\), \(d=0\)}{\(a=3\), \(b=-1\), \(c=4\), \(d=1\)}{}{}}
\End

\Data\ID{2000017404}
\Updated{2024-03-03 09:03:25}
\Level{A}
\SubArea{Matrices and determinants}
\LANGen{
\Question{
Consider the matrices:
\[ 
A=\left (\MATRIX{ 
1& -1 & 2\cr 
3 & 4 & 1 \cr 
0 & 1 & 0} \right ),  
\ B=\left (\MATRIX{ 
2& 3 & -1\cr 
-2 & 1 & 1 \cr 
1 & 4 & -1 } \right ).
\]
Find the transpose of the matrix \(A+B\).}
{\(
\left (\MATRIX{ 
3& 1 & 1\cr 
 2 & 5& 5 \cr 
1 & 2& -1} \right )
\)}{\(
\left (\MATRIX{ 
3& 2 & 1\cr 
1& 5 & 2\cr 
1 & 5& -1 } \right )
\)}{\(
\left (\MATRIX{ 
3& 1 & 1\cr 
2 & 5&5 \cr 
1 & 2& 1 } \right )
\)}{\(
\left (\MATRIX{ 
3& 2 & 1\cr 
 2 &5& 2 \cr 
1&2& -1 } \right )
\)}{}{}}
\LANGpl{
\Question{
Dane są macierze:
\[ 
A=\left (\MATRIX{ 
1& -1 & 2\cr 
3 & 4 & 1 \cr 
0 & 1 & 0} \right ),
\ B=\left (\MATRIX{ 
2& 3 & -1\cr 
-2 & 1 & 1 \cr 
1 & 4 & -1 } \right ).
\]
Znajdź macierz transponowaną do \(A+B\).}
{\(
\left (\MATRIX{ 
3& 1 & 1\cr 
 2 & 5& 5 \cr 
1 & 2& -1} \right )
\)}{\(
\left (\MATRIX{ 
3& 2 & 1\cr 
1& 5 & 2\cr 
1 & 5& -1 } \right )
\)}{\(
\left (\MATRIX{ 
3& 1 & 1\cr 
2 & 5&5 \cr 
1 & 2& 1 } \right )
\)}{\(
\left (\MATRIX{ 
3& 2 & 1\cr 
 2 &5& 2 \cr 
1&2& -1 } \right )
\)}{}{}}
\LANGcs{
\Question{
Jsou dány matice:
\[ 
A=\left (\MATRIX{ 
1& -1 & 2\cr 
3 & 4 & 1 \cr 
0 & 1 & 0} \right ),
\ B=\left (\MATRIX{ 
2& 3 & -1\cr 
-2 & 1 & 1 \cr 
1 & 4 & -1 } \right ).
\]
Určete matici transponovanou k matici \(A+B\).}
{\(
\left (\MATRIX{ 
3& 1 & 1\cr 
 2 & 5& 5 \cr 
1 & 2& -1} \right )
\)}{\(
\left (\MATRIX{ 
3& 2 & 1\cr 
1& 5 & 2\cr 
1 & 5& -1 } \right )
\)}{\(
\left (\MATRIX{ 
3& 1 & 1\cr 
2 & 5&5 \cr 
1 & 2& 1 } \right )
\)}{\(
\left (\MATRIX{ 
3& 2 & 1\cr 
 2 &5& 2 \cr 
1&2& -1 } \right )
\)}{}{}}
\LANGsk{
\Question{
Sú dané matice:
\[ 
A=\left (\MATRIX{ 
1& -1 & 2\cr 
3 & 4 & 1 \cr 
0 & 1 & 0} \right ),
\ B=\left (\MATRIX{ 
2& 3 & -1\cr 
-2 & 1 & 1 \cr 
1 & 4 & -1 } \right ).
\]
Určte transponovanú maticu k matici \(A+B\).}
{\(
\left (\MATRIX{ 
3& 1 & 1\cr 
 2 & 5& 5 \cr 
1 & 2& -1} \right )
\)}{\(
\left (\MATRIX{ 
3& 2 & 1\cr 
1& 5 & 2\cr 
1 & 5& -1 } \right )
\)}{\(
\left (\MATRIX{ 
3& 1 & 1\cr 
2 & 5&5 \cr 
1 & 2& 1 } \right )
\)}{\(
\left (\MATRIX{ 
3& 2 & 1\cr 
 2 &5& 2 \cr 
1&2& -1 } \right )
\)}{}{}}
\LANGes{
\Question{
Dadas las matrices:
\[ 
A=\left (\MATRIX{ 
1& -1 & 2\cr 
3 & 4 & 1 \cr 
0 & 1 & 0} \right ),
\ B=\left (\MATRIX{ 
2& 3 & -1\cr 
-2 & 1 & 1 \cr 
1 & 4 & -1 } \right ).
\]
Calcula la traspuesta \(A+B\).}
{\(
\left (\MATRIX{ 
3& 1 & 1\cr 
 2 & 5& 5 \cr 
1 & 2& -1} \right )
\)}{\(
\left (\MATRIX{ 
3& 2 & 1\cr 
1& 5 & 2\cr 
1 & 5& -1 } \right )
\)}{\(
\left (\MATRIX{ 
3& 1 & 1\cr 
2 & 5&5 \cr 
1 & 2& 1 } \right )
\)}{\(
\left (\MATRIX{ 
3& 2 & 1\cr 
 2 &5& 2 \cr 
1&2& -1 } \right )
\)}{}{}}
\End

\Data\ID{2000017405}
\Updated{2024-03-03 09:03:25}
\Level{A}
\SubArea{Matrices and determinants}
\LANGen{
\Question{
For what value of \(a\) does the product given below result in a diagonal matrix?
\[ 
\left (\MATRIX{ 
3& -2\cr 
 a & -2 \cr 
} \right ) 
\cdot
\left (\MATRIX{ 
4& 4 \cr 
10 & 6 \cr 
} \right ) 
\]}
{\(5\)}{\(-5\)}{\(4\)}{\(-4\)}{}{}}
\LANGpl{
\Question{
Dla jakiej wartości \(a\) podany poniżej iloczyn daje macierz diagonalną?
\[ 
\left (\MATRIX{ 
3& -2\cr 
 a & -2 \cr 
} \right ) 
\cdot
\left (\MATRIX{ 
4& 4 \cr 
10 & 6 \cr 
} \right ) 
\]}
{\(5\)}{\(-5\)}{\(4\)}{\(-4\)}{}{}}
\LANGcs{
\Question{
Pro kterou hodnotu \(a\) tvoří uvedený součin matic diagonální matici?
\[ 
\left (\MATRIX{ 
3& -2\cr 
 a & -2 \cr 
} \right ) 
\cdot
\left (\MATRIX{ 
4& 4 \cr 
10 & 6 \cr 
} \right ) 
\]}
{\(5\)}{\(-5\)}{\(4\)}{\(-4\)}{}{}}
\LANGsk{
\Question{
Pre ktorú hodnotu \(a\) tvorí uvedený súčin matíc diagonálnu maticu?
\[ 
\left (\MATRIX{ 
3& -2\cr 
 a & -2 \cr 
} \right ) 
\cdot
\left (\MATRIX{ 
4& 4 \cr 
10 & 6 \cr 
} \right ) 
\]}
{\(5\)}{\(-5\)}{\(4\)}{\(-4\)}{}{}}
\LANGes{
\Question{
¿Para qué valor de \(a\) el producto dado da como resultado una matriz diagonal?
\[ 
\left (\MATRIX{ 
3& -2\cr 
 a & -2 \cr 
} \right ) 
\cdot
\left (\MATRIX{ 
4& 4 \cr 
10 & 6 \cr 
} \right ) 
\]}
{\(5\)}{\(-5\)}{\(4\)}{\(-4\)}{}{}}
\End

\Data\ID{2000017406}
\Updated{2024-03-03 09:03:25}
\Level{A}
\SubArea{Matrices and determinants}
\LANGen{
\Question{
For which numbers \(x\) and \(y\) is the following equality true?
\[ 
\left (\MATRIX{ 
2& 3\cr 
 1 & -5 \cr 
} \right ) 
\cdot
\left (\MATRIX{ 
x\cr 
 y  \cr 
} \right ) 
=
\left (\MATRIX{ 
1.5& x \cr 
y & 5 \cr 
} \right ) 
\cdot
\left (\MATRIX{ 
-2\cr 
 -1  \cr 
} \right ) 
\]}
{\(x=-2\), \(y=1\)}{\(x=2\), \(y=1\)}{\(x=-5\), \(y=4\)}{\(x=-2\), \(y=-1\)}{}{}}
\LANGpl{
\Question{
Dla jakich liczb \(x\) i \(y\) podana równość jest prawdziwa?
\[ 
\left (\MATRIX{ 
2& 3\cr 
 1 & -5 \cr 
} \right ) 
\cdot
\left (\MATRIX{ 
x\cr 
 y  \cr 
} \right ) 
=
\left (\MATRIX{ 
1{,}5& x \cr 
y & 5 \cr 
} \right ) 
\cdot
\left (\MATRIX{ 
-2\cr 
 -1  \cr 
} \right ) 
\]}
{\(x=-2\), \(y=1\)}{\(x=2\), \(y=1\)}{\(x=-5\), \(y=4\)}{\(x=-2\), \(y=-1\)}{}{}}
\LANGcs{
\Question{
Pro která čísla \(x\) a \(y\) platí následující rovnost?
\[ 
\left (\MATRIX{ 
2& 3\cr 
 1 & -5 \cr 
} \right ) 
\cdot
\left (\MATRIX{ 
x\cr 
 y  \cr 
} \right ) 
=
\left (\MATRIX{ 
1{,}5& x \cr 
y & 5 \cr 
} \right ) 
\cdot
\left (\MATRIX{ 
-2\cr 
 -1  \cr 
} \right ) 
\]}
{\(x=-2\), \(y=1\)}{\(x=2\), \(y=1\)}{\(x=-5\), \(y=4\)}{\(x=-2\), \(y=-1\)}{}{}}
\LANGsk{
\Question{
Pre ktoré čísla \(x\) a \(y\) platí naslednujúca rovnosť?
\[ 
\left (\MATRIX{ 
2& 3\cr 
 1 & -5 \cr 
} \right ) 
\cdot
\left (\MATRIX{ 
x\cr 
 y  \cr 
} \right ) 
=
\left (\MATRIX{ 
1{,}5& x \cr 
y & 5 \cr 
} \right ) 
\cdot
\left (\MATRIX{ 
-2\cr 
 -1  \cr 
} \right ) 
\]}
{\(x=-2\), \(y=1\)}{\(x=2\), \(y=1\)}{\(x=-5\), \(y=4\)}{\(x=-2\), \(y=-1\)}{}{}}
\LANGes{
\Question{
¿Para qué números \(x\) y \(y\) es la siguiente ecuación verdadera?
\[ 
\left (\MATRIX{ 
2& 3\cr 
 1 & -5 \cr 
} \right ) 
\cdot
\left (\MATRIX{ 
x\cr 
 y  \cr 
} \right ) 
=
\left (\MATRIX{ 
1.5& x \cr 
y & 5 \cr 
} \right ) 
\cdot
\left (\MATRIX{ 
-2\cr 
 -1  \cr 
} \right ) 
\]}
{\(x=-2\), \(y=1\)}{\(x=2\), \(y=1\)}{\(x=-5\), \(y=4\)}{\(x=-2\), \(y=-1\)}{}{}}
\End

\Data\ID{2000017407}
\Updated{2024-03-03 09:03:25}
\Level{A}
\SubArea{Matrices and determinants}
\LANGen{
\Question{
Find \(A^2-B^2\), if:
\[ 
A=\left (\MATRIX{ 
2& 1\cr 
 3 & 0 \cr 
} \right ) ,
B=\left (\MATRIX{ 
1& 4 \cr 
2 & -1 \cr 
} \right ) 
\]}
{\( \left (\MATRIX{ 
-2& 2 \cr 
6 & -6\cr 
} \right )  \)}{\( \left (\MATRIX{ 
-2& 2 \cr 
6 & 6\cr 
} \right )  \)}{\( \left (\MATRIX{ 
-2& 2 \cr 
6 & -4\cr 
} \right )  \)}{\( \left (\MATRIX{ 
2& 2 \cr 
6 & 6\cr 
} \right )  \)}{}{}}
\LANGpl{
\Question{
Znajdź \(A^2-B^2\), jeśli:
\[ 
A=\left (\MATRIX{ 
2& 1\cr 
 3 & 0 \cr 
} \right ) ,
B=\left (\MATRIX{ 
1& 4 \cr 
2 & -1 \cr 
} \right ) 
\]}
{\( \left (\MATRIX{ 
-2& 2 \cr 
6 & -6\cr 
} \right )  \)}{\( \left (\MATRIX{ 
-2& 2 \cr 
6 & 6\cr 
} \right )  \)}{\( \left (\MATRIX{ 
-2& 2 \cr 
6 & -4\cr 
} \right )  \)}{\( \left (\MATRIX{ 
2& 2 \cr 
6 & 6\cr 
} \right )  \)}{}{}}
\LANGcs{
\Question{
Určete \(A^2-B^2\), jestliže:
\[ 
A=\left (\MATRIX{ 
2& 1\cr 
 3 & 0 \cr 
} \right ) ,
B=\left (\MATRIX{ 
1& 4 \cr 
2 & -1 \cr 
} \right ) 
\]}
{\( \left (\MATRIX{ 
-2& 2 \cr 
6 & -6\cr 
} \right )  \)}{\( \left (\MATRIX{ 
-2& 2 \cr 
6 & 6\cr 
} \right )  \)}{\( \left (\MATRIX{ 
-2& 2 \cr 
6 & -4\cr 
} \right )  \)}{\( \left (\MATRIX{ 
2& 2 \cr 
6 & 6\cr 
} \right )  \)}{}{}}
\LANGsk{
\Question{
Určte \(A^2-B^2\), ak:
\[ 
A=\left (\MATRIX{ 
2& 1\cr 
 3 & 0 \cr 
} \right ) ,
B=\left (\MATRIX{ 
1& 4 \cr 
2 & -1 \cr 
} \right ) 
\]}
{\( \left (\MATRIX{ 
-2& 2 \cr 
6 & -6\cr 
} \right )  \)}{\( \left (\MATRIX{ 
-2& 2 \cr 
6 & 6\cr 
} \right )  \)}{\( \left (\MATRIX{ 
-2& 2 \cr 
6 & -4\cr 
} \right )  \)}{\( \left (\MATRIX{ 
2& 2 \cr 
6 & 6\cr 
} \right )  \)}{}{}}
\LANGes{
\Question{
Halla \(A^2-B^2\), suponiendo que:
\[ 
A=\left (\MATRIX{ 
2& 1\cr 
 3 & 0 \cr 
} \right ) ,
B=\left (\MATRIX{ 
1& 4 \cr 
2 & -1 \cr 
} \right ) 
\]}
{\( \left (\MATRIX{ 
-2& 2 \cr 
6 & -6\cr 
} \right )  \)}{\( \left (\MATRIX{ 
-2& 2 \cr 
6 & 6\cr 
} \right )  \)}{\( \left (\MATRIX{ 
-2& 2 \cr 
6 & -4\cr 
} \right )  \)}{\( \left (\MATRIX{ 
2& 2 \cr 
6 & 6\cr 
} \right )  \)}{}{}}
\End

\Data\ID{2000017408}
\Updated{2024-03-03 09:03:25}
\Level{A}
\SubArea{Matrices and determinants}
\LANGen{
\Question{
Find \(\frac{K \cdot (-L)}2\), if:
\[ 
K=\left (\MATRIX{ 
3& 0 & 1\cr 
 2 & 3 & 4 \cr 
1& -1 & 1} \right ),~
L=\left (\MATRIX{ 
2& 3 & 0\cr 
1 & 1 & -1 \cr 
2 &0& 1 } \right )
\]}
{\(
\left (\MATRIX{ 
-4& -4.5 & -0.5\cr 
 -7.5 & -4.5& -0.5\cr 
-1.5 & -1& -1 } \right )
\)}{\(
\left (\MATRIX{ 
-4& 4.5 & 0.5\cr 
 -7.5 & -4.5& -0.5\cr 
-1.5 & -1& -1 } \right )
\)}{\(
\left (\MATRIX{ 
-4& -4.5 & -0.5\cr 
 7.5 & -4.5& -0.5\cr 
-1.5 & -1& -1 } \right )
\)}{\(
\left (\MATRIX{ 
-4& -4.5 & -0.5\cr 
 -7.5 & 4.5& -0.5\cr 
-1.5 & -1& 1 } \right )
\)}{}{}}
\LANGpl{
\Question{
Znajdź \(\frac{K \cdot (-L)}2\), jeśli:
\[ 
K=\left (\MATRIX{ 
3& 0 & 1\cr 
 2 & 3 & 4 \cr 
1& -1 & 1} \right ),~
L=\left (\MATRIX{ 
2& 3 & 0\cr 
1 & 1 & -1 \cr 
2 &0& 1 } \right )
\]}
{\(
\left (\MATRIX{ 
-4& -4{,}5 & -0{,}5\cr 
 -7{,}5 & -4{,}5& -0{,}5\cr 
-1{,}5 & -1& -1 } \right )
\)}{\(
\left (\MATRIX{ 
-4& 4{,}5 & 0{,}5\cr 
 -7{,}5 & -4{,}5& -0{,}5\cr 
-1{,}5 & -1& -1 } \right )
\)}{\(
\left (\MATRIX{ 
-4& -4{,}5 & -0{,}5\cr 
 7{,}5 & -4{,}5& -0{,}5\cr 
-1{,}5 & -1& -1 } \right )
\)}{\(
\left (\MATRIX{ 
-4& -4{,}5 & -0{,}5\cr 
 -7{,}5 & 4{,}5& -0{,}5\cr 
-1{,}5 & -1& 1 } \right )
\)}{}{}}
\LANGcs{
\Question{
Určete \(\frac{K \cdot (-L)}2\), jestliže:
\[ 
K=\left (\MATRIX{ 
3& 0 & 1\cr 
 2 & 3 & 4 \cr 
1& -1 & 1} \right ),~
L=\left (\MATRIX{ 
2& 3 & 0\cr 
1 & 1 & -1 \cr 
2 &0& 1 } \right )
\]}
{\(
\left (\MATRIX{ 
-4& -4{,}5 & -0{,}5\cr 
 -7{,}5 & -4{,}5& -0{,}5\cr 
-1{,}5 & -1& -1 } \right )
\)}{\(
\left (\MATRIX{ 
-4& 4{,}5 & 0{,}5\cr 
 -7{,}5 & -4{,}5& -0{,}5\cr 
-1{,}5 & -1& -1 } \right )
\)}{\(
\left (\MATRIX{ 
-4& -4{,}5 & -0{,}5\cr 
 7{,}5 & -4{,}5& -0{,}5\cr 
-1{,}5 & -1& -1 } \right )
\)}{\(
\left (\MATRIX{ 
-4& -4{,}5 & -0{,}5\cr 
 -7{,}5 & 4{,}5& -0{,}5\cr 
-1{,}5 & -1& 1 } \right )
\)}{}{}}
\LANGsk{
\Question{
Nájdi \(\frac{K \cdot (-L)}2\), ak:
\[ 
K=\left (\MATRIX{ 
3& 0 & 1\cr 
 2 & 3 & 4 \cr 
1& -1 & 1} \right ),~
L=\left (\MATRIX{ 
2& 3 & 0\cr 
1 & 1 & -1 \cr 
2 &0& 1 } \right )
\]}
{\(
\left (\MATRIX{ 
-4& -4{,}5 & -0{,}5\cr 
 -7{,}5 & -4{,}5& -0{,}5\cr 
-1{,}5 & -1& -1 } \right )
\)}{\(
\left (\MATRIX{ 
-4& 4{,}5 & 0{,}5\cr 
 -7{,}5 & -4{,}5& -0{,}5\cr 
-1{,}5 & -1& -1 } \right )
\)}{\(
\left (\MATRIX{ 
-4& -4{,}5 & -0{,}5\cr 
 7{,}5 & -4{,}5& -0{,}5\cr 
-1{,}5 & -1& -1 } \right )
\)}{\(
\left (\MATRIX{ 
-4& -4{,}5 & -0{,}5\cr 
 -7{,}5 & 4{,}5& -0{,}5\cr 
-1{,}5 & -1& 1 } \right )
\)}{}{}}
\LANGes{
\Question{
Halla \(\frac{K \cdot (-L)}2\), suponiendo que:
\[ 
K=\left (\MATRIX{ 
3& 0 & 1\cr 
 2 & 3 & 4 \cr 
1& -1 & 1} \right ),~
L=\left (\MATRIX{ 
2& 3 & 0\cr 
1 & 1 & -1 \cr 
2 &0& 1 } \right )
\]}
{\(
\left (\MATRIX{ 
-4& -4.5 & -0.5\cr 
 -7.5 & -4.5& -0.5\cr 
-1.5 & -1& -1 } \right )
\)}{\(
\left (\MATRIX{ 
-4& 4.5 & 0.5\cr 
 -7.5 & -4.5& -0.5\cr 
-1.5 & -1& -1 } \right )
\)}{\(
\left (\MATRIX{ 
-4& -4.5 & -0.5\cr 
 7.5 & -4.5& -0.5\cr 
-1.5 & -1& -1 } \right )
\)}{\(
\left (\MATRIX{ 
-4& -4.5 & -0.5\cr 
 -7.5 & 4.5& -0.5\cr 
-1.5 & -1& 1 } \right )
\)}{}{}}
\End

\Data\ID{2000017409}
\Updated{2024-03-03 09:03:25}
\Level{A}
\SubArea{Matrices and determinants}
\LANGen{
\Question{
Let
\[ 
A=\left (\MATRIX{ 
1& 2\cr 
 a & 1 \cr 
} \right ) ,
\ B=\left (\MATRIX{ 
7& 4 \cr 
6 & 7 \cr 
} \right ),
\]
where \(a\) is a real number. For which \(a\) the equality \(A \cdot B = B \cdot A\) holds?}
{\( 3\)}{\( -3\)}{\( 3.5\)}{\( -3.5\)}{}{}}
\LANGpl{
\Question{
Dane są
\[ 
A=\left (\MATRIX{ 
1& 2\cr 
 a & 1 \cr 
} \right ) ,
\ B=\left (\MATRIX{ 
7& 4 \cr 
6 & 7 \cr 
} \right ),
\]
gdzie \(a\) to liczba rzeczywista. Dla jakiej wartości \(a\) zachodzi równość\(A \cdot B = B \cdot A\)?}
{\( 3\)}{\( -3\)}{\( 3{,}5\)}{\( -3{,}5\)}{}{}}
\LANGcs{
\Question{
Nechť
\[ 
A=\left (\MATRIX{ 
1& 2\cr 
 a & 1 \cr 
} \right ) ,
\ B=\left (\MATRIX{ 
7& 4 \cr 
6 & 7 \cr 
} \right ),
\]
kde \(a\) je reálné číslo. Pro jakou hodnotu \(a\) platí rovnost \(A \cdot B = B \cdot A\)?}
{\( 3\)}{\( -3\)}{\( 3{,}5\)}{\( -3{,}5\)}{}{}}
\LANGsk{
\Question{
Nech
\[ 
A=\left (\MATRIX{ 
1& 2\cr 
 a & 1 \cr 
} \right ) ,
\ B=\left (\MATRIX{ 
7& 4 \cr 
6 & 7 \cr 
} \right ),
\]
kde \(a\) je reálne číslo. Pre ktoré \(a\) platí rovnosť \(A \cdot B = B \cdot A\)?}
{\( 3\)}{\( -3\)}{\( 3{,}5\)}{\( -3{,}5\)}{}{}}
\LANGes{
\Question{
Suponiendo que
\[ 
A=\left (\MATRIX{ 
1& 2\cr 
 a & 1 \cr 
} \right ) ,
\ B=\left (\MATRIX{ 
7& 4 \cr 
6 & 7 \cr 
} \right ),
\]
donde \(a\) es un número real. ¿Para qué \(a\) se cumple la ecuación \(A \cdot B = B \cdot A\)?}
{\( 3\)}{\( -3\)}{\( 3.5\)}{\( -3.5\)}{}{}}
\End

\Data\ID{2000018301}
\Updated{2024-03-03 09:03:23}
\Level{A}
\SubArea{Matrices and determinants}
\LANGen{
\Question{
Find the matrix \(B\), the solution to the equation given below.
\[ 
\left (\MATRIX{ 
3&-1 &5\cr 
 1 &0&3
} \right ) 
+
B
=
\left (\MATRIX{ 
5 & 0 & 4 \cr 
3 & 2 & 1\cr 
} \right ) 
\]}
{\[ 
B=
\left (\MATRIX{ 
2 & 1 & -1\cr 
2 & 2 & -2
} \right ) 
\]}{\[ 
B=
\left (\MATRIX{ 
2 & -1 & -1\cr 
2 & 2 & -2
} \right ) 
\]}{\[ 
B=
\left (\MATRIX{ 
2 & 1 & -1\cr 
2 & -2 & -2
} \right ) 
\]}{\[ 
B=
\left (\MATRIX{ 
2 & 1 & -1\cr 
2 & 2 & 2
} \right ) 
\]}{}{}}
\LANGpl{
\Question{
Znajdź macierz \(B\), rozwiązanie równania podanego poniżej.
\[ 
\left (\MATRIX{ 
3&-1 &5\cr 
 1 &0&3
} \right ) 
+
B
=
\left (\MATRIX{ 
5 & 0 & 4 \cr 
3 & 2 & 1\cr 
} \right ) 
\]}
{\[ 
B=
\left (\MATRIX{ 
2 & 1 & -1\cr 
2 & 2 & -2
} \right ) 
\]}{\[ 
B=
\left (\MATRIX{ 
2 & -1 & -1\cr 
2 & 2 & -2
} \right ) 
\]}{\[ 
B=
\left (\MATRIX{ 
2 & 1 & -1\cr 
2 & -2 & -2
} \right ) 
\]}{\[ 
B=
\left (\MATRIX{ 
2 & 1 & -1\cr 
2 & 2 & 2
} \right ) 
\]}{}{}}
\LANGcs{
\Question{
Nalezněte matici \(B\), která je řešením níže uvedené rovnice.
\[ 
\left (\MATRIX{ 
3&-1 &5\cr 
 1 &0&3
} \right ) 
+
B
=
\left (\MATRIX{ 
5 & 0 & 4 \cr 
3 & 2 & 1\cr 
} \right ) 
\]}
{\[ 
B=
\left (\MATRIX{ 
2 & 1 & -1\cr 
2 & 2 & -2
} \right ) 
\]}{\[ 
B=
\left (\MATRIX{ 
2 & -1 & -1\cr 
2 & 2 & -2
} \right ) 
\]}{\[ 
B=
\left (\MATRIX{ 
2 & 1 & -1\cr 
2 & -2 & -2
} \right ) 
\]}{\[ 
B=
\left (\MATRIX{ 
2 & 1 & -1\cr 
2 & 2 & 2
} \right ) 
\]}{}{}}
\LANGsk{
\Question{
Nájdite maticu \(B\), ktorá je riešením nižšie uvedenej rovnice.
\[ 
\left (\MATRIX{ 
3&-1 &5\cr 
 1 &0&3
} \right ) 
+
B
=
\left (\MATRIX{ 
5 & 0 & 4 \cr 
3 & 2 & 1\cr 
} \right ) 
\]}
{\[ 
B=
\left (\MATRIX{ 
2 & 1 & -1\cr 
2 & 2 & -2
} \right ) 
\]}{\[ 
B=
\left (\MATRIX{ 
2 & -1 & -1\cr 
2 & 2 & -2
} \right ) 
\]}{\[ 
B=
\left (\MATRIX{ 
2 & 1 & -1\cr 
2 & -2 & -2
} \right ) 
\]}{\[ 
B=
\left (\MATRIX{ 
2 & 1 & -1\cr 
2 & 2 & 2
} \right ) 
\]}{}{}}
\LANGes{
\Question{
Halla la matriz \(B\):
\[ 
\left (\MATRIX{ 
3&-1 &5\cr 
 1 &0&3
} \right ) 
+
B
=
\left (\MATRIX{ 
5 & 0 & 4 \cr 
3 & 2 & 1\cr 
} \right ) 
\]}
{\[ 
B=
\left (\MATRIX{ 
2 & 1 & -1\cr 
2 & 2 & -2
} \right ) 
\]}{\[ 
B=
\left (\MATRIX{ 
2 & -1 & -1\cr 
2 & 2 & -2
} \right ) 
\]}{\[ 
B=
\left (\MATRIX{ 
2 & 1 & -1\cr 
2 & -2 & -2
} \right ) 
\]}{\[ 
B=
\left (\MATRIX{ 
2 & 1 & -1\cr 
2 & 2 & 2
} \right ) 
\]}{}{}}
\End

\Data\ID{2000018302}
\Updated{2024-03-03 09:03:23}
\Level{A}
\SubArea{Matrices and determinants}
\LANGen{
\Question{
Find the matrix \(M\) so that the equality given below is true.
\[ 
2 \cdot \left (\MATRIX{ 
-1&4\cr 
3&-5\cr 
} \right ) 
-
M
=
\left (\MATRIX{ 
-3 &6\cr 
9 & -14\cr 
} \right ) 
\]}
{\[ 
M=\left (\MATRIX{ 
1 &2\cr 
-3 & 4\cr 
} \right ) 
\]}{\[ 
M=\left (\MATRIX{ 
-1 &2\cr 
-3 & 4\cr 
} \right ) 
\]}{\[ 
M=\left (\MATRIX{ 
-1 &-2\cr 
3 & -4\cr 
} \right ) 
\]}{\[ 
M=\left (\MATRIX{ 
1 &2\cr 
3 & -4\cr 
} \right ) 
\]}{}{}}
\LANGpl{
\Question{
Znajdź macierz \(M\) tak, aby podana poniżej równość była prawdziwa.
\[ 
2 \cdot \left (\MATRIX{ 
-1&4\cr 
3&-5\cr 
} \right ) 
-
M
=
\left (\MATRIX{ 
-3 &6\cr 
9 & -14\cr 
} \right ) 
\]}
{\[ 
M=\left (\MATRIX{ 
1 &2\cr 
-3 & 4\cr 
} \right ) 
\]}{\[ 
M=\left (\MATRIX{ 
-1 &2\cr 
-3 & 4\cr 
} \right ) 
\]}{\[ 
M=\left (\MATRIX{ 
-1 &-2\cr 
3 & -4\cr 
} \right ) 
\]}{\[ 
M=\left (\MATRIX{ 
1 &2\cr 
3 & -4\cr 
} \right ) 
\]}{}{}}
\LANGcs{
\Question{
Pro kterou matici \(M\) platí následující rovnost?
\[ 
2 \cdot \left (\MATRIX{ 
-1&4\cr 
3&-5\cr 
} \right ) 
-
M
=
\left (\MATRIX{ 
-3 &6\cr 
9 & -14\cr 
} \right ) 
\]}
{\[ 
M=\left (\MATRIX{ 
1 &2\cr 
-3 & 4\cr 
} \right ) 
\]}{\[ 
M=\left (\MATRIX{ 
-1 &2\cr 
-3 & 4\cr 
} \right ) 
\]}{\[ 
M=\left (\MATRIX{ 
-1 &-2\cr 
3 & -4\cr 
} \right ) 
\]}{\[ 
M=\left (\MATRIX{ 
1 &2\cr 
3 & -4\cr 
} \right ) 
\]}{}{}}
\LANGsk{
\Question{
Pre ktorú maticu \(M\) platí nasledujúca rovnosť
\[ 
2 \cdot \left (\MATRIX{ 
-1&4\cr 
3&-5\cr 
} \right ) 
-
M
=
\left (\MATRIX{ 
-3 &6\cr 
9 & -14\cr 
} \right ) 
\]}
{\[ 
M=\left (\MATRIX{ 
1 &2\cr 
-3 & 4\cr 
} \right ) 
\]}{\[ 
M=\left (\MATRIX{ 
-1 &2\cr 
-3 & 4\cr 
} \right ) 
\]}{\[ 
M=\left (\MATRIX{ 
-1 &-2\cr 
3 & -4\cr 
} \right ) 
\]}{\[ 
M=\left (\MATRIX{ 
1 &2\cr 
3 & -4\cr 
} \right ) 
\]}{}{}}
\LANGes{
\Question{
Halla la matriz \(M\):
\[ 
2 \cdot \left (\MATRIX{ 
-1&4\cr 
3&-5\cr 
} \right ) 
-
M
=
\left (\MATRIX{ 
-3 &6\cr 
9 & -14\cr 
} \right ) 
\]}
{\[ 
M=\left (\MATRIX{ 
1 &2\cr 
-3 & 4\cr 
} \right ) 
\]}{\[ 
M=\left (\MATRIX{ 
-1 &2\cr 
-3 & 4\cr 
} \right ) 
\]}{\[ 
M=\left (\MATRIX{ 
-1 &-2\cr 
3 & -4\cr 
} \right ) 
\]}{\[ 
M=\left (\MATRIX{ 
1 &2\cr 
3 & -4\cr 
} \right ) 
\]}{}{}}
\End

\Data\ID{2000018303}
\Updated{2024-03-03 09:03:23}
\Level{A}
\SubArea{Matrices and determinants}
\LANGen{
\Question{
Let \(E\) denote an identity matrix of order \(2\) and let matrix
\[ 
M
=
\left (\MATRIX{ 
m &0\cr 
0 & 2\cr 
} \right ) .
\]
Find all the values of \(m\) so that the equality below holds. 
\[
M^2-\frac52M+E=0
\]}
{\(m=2\) or \(m=\frac12\)}{\(m=\frac12\)}{\(m=2\)}{\(m=2\) or \(m=-\frac12\)}{}{}}
\LANGpl{
\Question{
\(E\) oznacza macierz jednostkową rzędu \(2\) oraz 
\[ 
M
=
\left (\MATRIX{ 
m &0\cr 
0 & 2\cr 
} \right ) .
\]
Znajdź wszystkie wartości \(m\),  dla których zachodzi równość. 
\[
M^2-\frac52M+E=0
\]}
{\(m=2\) lub \(m=\frac12\)}{\(m=\frac12\)}{\(m=2\)}{\(m=2\) lub \(m=-\frac12\)}{}{}}
\LANGcs{
\Question{
Nechť \(E\) označuje jednotkovou matici řádu \(2\) a matice
\[ 
M
=
\left (\MATRIX{ 
m &0\cr 
0 & 2\cr 
} \right ) .
\]
Pro která reálná čísla \(m\) platí následující rovnost?
\[
M^2-\frac52M+E=0
\]}
{pro \(m=2\) nebo pro \(m=\frac12\)}{jen pro \(m=\frac12\)}{jen pro \(m=2\)}{pro \(m=2\) nebo pro \(m=-\frac12\)}{}{}}
\LANGsk{
\Question{
Nech \(E\) označuje jednotkovú maticu rádu \(2\) a maticu
\[ 
M
=
\left (\MATRIX{ 
m &0\cr 
0 & 2\cr 
} \right ) .
\]
Pre ktoré reálne čísla \(m\) platí nasledujúca rovnosť?
\[
M^2-\frac52M+E=0
\]}
{\(m=2\) alebo \(m=\frac12\)}{\(m=\frac12\)}{\(m=2\)}{\(m=2\) alebo \(m=-\frac12\)}{}{}}
\LANGes{
\Question{
Suponiendo que \(E\) es una matriz identidad de orden \(2\) y la matriz:
\[ 
M
=
\left (\MATRIX{ 
m &0\cr 
0 & 2\cr 
} \right ) .
\]
Halla todos los valores \(m\) para que se cumpla la siguiente ecuación:
\[
M^2-\frac52M+E=0.
\]}
{\(m=2\) o \(m=\frac12\)}{\(m=\frac12\)}{\(m=2\)}{\(m=2\) o \(m=-\frac12\)}{}{}}
\End

\Data\ID{2000018305}
\Updated{2024-03-03 09:03:23}
\Level{A}
\SubArea{Matrices and determinants}
\LANGen{
\Question{
Consider three matrices
\[
A
=
\left (\MATRIX{ 
3 &4\cr 
1 & 2\cr 
} \right ),~
B
=
\left (\MATRIX{ 
1 &1\cr 
0&1\cr 
} \right ),~
C
=
\left (\MATRIX{ 
1 &0\cr 
1&1\cr 
} \right ).
\]
Let \(E\) denote the identity matrix of order \(2\). Then, find \(X\), which is the solution to the following equation.
\[
C \cdot (A+X)\cdot B=E\]}
{\(
X
=
\left (\MATRIX{ 
-2 &-5\cr 
-2& 0\cr 
} \right ) 
\)}{\(
X
=
\left (\MATRIX{ 
-2 &-5\cr 
2& 0\cr 
} \right ) 
\)}{\(
X
=
\left (\MATRIX{ 
-2 &5\cr 
-2& 0\cr 
} \right ) 
\)}{\(
X
=
\left (\MATRIX{ 
-2 &5\cr 
2& 0\cr 
} \right ) 
\)}{}{}}
\LANGpl{
\Question{
Dane są trzy macierze
\[
A
=
\left (\MATRIX{ 
3 &4\cr 
1 & 2\cr 
} \right ),~
B
=
\left (\MATRIX{ 
1 &1\cr 
0&1\cr 
} \right ),~
C
=
\left (\MATRIX{ 
1 &0\cr 
1&1\cr 
} \right ).
\]
Niech \(E\) oznacza macierz jednostkową rzędu \(2\). Znajdź \(X\), które jest rozwiązaniem następującego równania.
\[
C \cdot (A+X)\cdot B=E\]}
{\(
X
=
\left (\MATRIX{ 
-2 &-5\cr 
-2& 0\cr 
} \right ) 
\)}{\(
X
=
\left (\MATRIX{ 
-2 &-5\cr 
2& 0\cr 
} \right ) 
\)}{\(
X
=
\left (\MATRIX{ 
-2 &5\cr 
-2& 0\cr 
} \right ) 
\)}{\(
X
=
\left (\MATRIX{ 
-2 &5\cr 
2& 0\cr 
} \right ) 
\)}{}{}}
\LANGcs{
\Question{
Jsou dány matice
\[
A
=
\left (\MATRIX{ 
3 &4\cr 
1 & 2\cr 
} \right ),~
B
=
\left (\MATRIX{ 
1 &1\cr 
0&1\cr 
} \right ),~
C
=
\left (\MATRIX{ 
1 &0\cr 
1&1\cr 
} \right ).
\]
Nechť \(E\) je jednotková matice řádu \(2\). Určete matici \(X\), která je řešením následující rovnice.
\[
C \cdot (A+X)\cdot B=E\]}
{\(
X
=
\left (\MATRIX{ 
-2 &-5\cr 
-2& 0\cr 
} \right ) 
\)}{\(
X
=
\left (\MATRIX{ 
-2 &-5\cr 
2& 0\cr 
} \right ) 
\)}{\(
X
=
\left (\MATRIX{ 
-2 &5\cr 
-2& 0\cr 
} \right ) 
\)}{\(
X
=
\left (\MATRIX{ 
-2 &5\cr 
2& 0\cr 
} \right ) 
\)}{}{}}
\LANGsk{
\Question{
Sú dané matice
\[
A
=
\left (\MATRIX{ 
3 &4\cr 
1 & 2\cr 
} \right ),~
B
=
\left (\MATRIX{ 
1 &1\cr 
0&1\cr 
} \right ),~
C
=
\left (\MATRIX{ 
1 &0\cr 
1&1\cr 
} \right ).
\]
Nech \(E\) je jednotková matica rádu \(2\). Určte maticu \(X\), ktorá je riešením nasledujúcej rovnice.
\[
C \cdot (A+X)\cdot B=E\]}
{\(
X
=
\left (\MATRIX{ 
-2 &-5\cr 
-2& 0\cr 
} \right ) 
\)}{\(
X
=
\left (\MATRIX{ 
-2 &-5\cr 
2& 0\cr 
} \right ) 
\)}{\(
X
=
\left (\MATRIX{ 
-2 &5\cr 
-2& 0\cr 
} \right ) 
\)}{\(
X
=
\left (\MATRIX{ 
-2 &5\cr 
2& 0\cr 
} \right ) 
\)}{}{}}
\LANGes{
\Question{
Dadas las tres matrices:
\[
A
=
\left (\MATRIX{ 
3 &4\cr 
1 & 2\cr 
} \right ),~
B
=
\left (\MATRIX{ 
1 &1\cr 
0&1\cr 
} \right ),~
C
=
\left (\MATRIX{ 
1 &0\cr 
1&1\cr 
} \right ).
\]
Suponiendo que \(E\) es una matriz identidad de orden \(2\). Halla \(X\) que es la solución de la siguiente ecuación.
\[
C \cdot (A+X)\cdot B=E\]}
{\(
X
=
\left (\MATRIX{ 
-2 &-5\cr 
-2& 0\cr 
} \right ) 
\)}{\(
X
=
\left (\MATRIX{ 
-2 &-5\cr 
2& 0\cr 
} \right ) 
\)}{\(
X
=
\left (\MATRIX{ 
-2 &5\cr 
-2& 0\cr 
} \right ) 
\)}{\(
X
=
\left (\MATRIX{ 
-2 &5\cr 
2& 0\cr 
} \right ) 
\)}{}{}}
\End

\Data\ID{2000019301}
\Updated{2024-03-03 09:03:23}
\Level{A}
\SubArea{Matrices and determinants}
\LANGen{
\Question{
Three ice cream stands of ICE company reported their July sales of four ice cream flavors in number of portions sold. All data can be seen in the table below.
\[ \begin{array}{|c|c|c|c|c|} \hline  &\text{vanilla} & \text{chocolate} & \text{nut} & \text{strawberry} \\\hline
\text{Stand 1}& 720 & 800 & 1\,200&360 \\\hline 
\text{Stand 2} & 550 & 434 & 900 & 300 \\\hline
\text{Stand 3} &610 &300 & 200 & 750 \\\hline \end{array}\]
The data provided by the stands for August sales are reported in the matrix \(A\).
\[ A=
\left (\MATRIX{ 
650& 470 & 890 & 410\cr 
500& 505 & 890 & 300\cr 
380& 520 & 350 & 800\cr } \right )
\]
If the July sales are rewritten to the matrix \(J\), by what matrix the sales of ice cream for both summer months are described?}
{matrix \(J+A\)}{matrix \(J-A\)}{matrix \(J \cdot A\)}{matrix \(2J+2A\)}{}{}}
\LANGpl{
\Question{
Trzy stoiska z lodami firmy ICE odnotowały lipcową sprzedaż czterech smaków lodów w ilości sprzedanych porcji. Wszystkie dane przedstawia poniższa tabela.
\[ \begin{array}{|c|c|c|c|c|} \hline  &\text{wanilia} & \text{czekolada} & \text{orzech} & \text{truskawka} \\\hline
\text{Stoisko 1}& 720 & 800 & 1\,200&360 \\\hline 
\text{Stoisko 2} & 550 & 434 & 900 & 300 \\\hline
\text{Stoisko 3} &610 &300 & 200 & 750 \\\hline \end{array}\]
Dane ze stoisk dotyczące sprzedaży sierpniowej są ujmowane w matrycy \(A\).
\[ A=
\left (\MATRIX{ 
650& 470 & 890 & 410\cr 
500& 505 & 890 & 300\cr 
380& 520 & 350 & 800\cr } \right )
\]
 Jeśli lipcowe wyprzedaże zostaną przepisane do matrycy \(J\), to jaką matrycą opisuje się sprzedaż lodów za oba letnie miesiące?}
{matryca \(J+A\)}{matryca \(J-A\)}{matryca \(J \cdot A\)}{matryca \(2J+2A\)}{}{}}
\LANGcs{
\Question{
Tři zmrzlinové stánky firmy ICE hlásily za červenec prodej porcí jednotlivých druhů zmrzliny. Údaje jsou uvedeny v následující tabulce.
\[ \begin{array}{|c|c|c|c|c|} \hline  &\text{vanilka} & \text{čokoláda} & \text{oříšek} & \text{jahoda} \\\hline
\text{Stánek 1}& 720 & 800 & 1\,200&360 \\\hline 
\text{Stánek 2} & 550 & 434 & 900 & 300 \\\hline
\text{Stánek 3} &610 &300 & 200 & 750 \\\hline \end{array}\]
Údaje, které stánky hlásily za srpen, jsou již stručně uvedeny v odpovídající matici  \(A\).
\[ A=
\left (\MATRIX{ 
650& 470 & 890 & 410\cr 
500& 505 & 890 & 300\cr 
380& 520 & 350 & 800\cr } \right )
\]
Pokud příslušnou matici prodeje za červenec označíme písmenem \(J\), jaká matice pak popisuje prodej jednotlivých druhů zmrzliny za oba letní měsíce?}
{matice \(J+A\)}{matice \(J-A\)}{matice \(J \cdot A\)}{matice \(2J+2A\)}{}{}}
\LANGsk{
\Question{
Tri zmrzlinové stánky firmy ICE hlásili za júl predaj porcií jednotlivých druhov zmrzliny. Údaje sú uvedené v nasledujúcej tabuľke.
\[ \begin{array}{|c|c|c|c|c|} \hline &\text{vanilka} & \text{čokoláda} & \text{oriešok} & \text{jahoda} \\\hline
\text{Stánok 1}& 720 & 800 & 1\,200&360 \\\hline
\text{Stánok 2} & 550 & 434 & 900 & 300 \\\hline
\text{Stánok 3} &610 &300 & 200 & 750 \\\hline \end{array}\]
Údaje, ktoré stánky hlásili za august, sú už stručne uvedené v zodpovedajúcej matici \(A\).
\[ A=
\left (\MATRIX{
650& 470 & 890 & 410\cr
500& 505 & 890 & 300\cr
380& 520 & 350 & 800\cr } \right )
\]
Ak príslušnú maticu predaja za júl označíme písmenom \(J\), aká matica potom popisuje predaj jednotlivých druhov zmrzliny za oba letné mesiace?}
{matica \(J+A\)}{matica \(J-A\)}{matica \(J \cdot A\)}{matica \(2J+2A\)}{}{}}
\LANGes{
\Question{
Tres puestos de helados de la empresa ICE informaron de sus ventas de julio de cuatro sabores de helado en número de porciones vendidas. Todos los datos en coronas checas pueden verse en la siguiente tabla.
\[ \begin{array}{|c|c|c|c|c|} \hline  &\text{vainilla} & \text{chocolate} & \text{nuez} & \text{fresa} \\\hline
\text{Puesto 1}& 720 & 800 & 1\,200&360 \\\hline 
\text{Puesto 2} & 550 & 434 & 900 & 300 \\\hline
\text{Puesto 3} &610 &300 & 200 & 750 \\\hline \end{array}\]
Los datos facilitados por los puestos para las ventas de agosto se recogen en la matriz \(A\).
\[ A=
\left (\MATRIX{ 
650& 470 & 890 & 410\cr 
500& 505 & 890 & 300\cr 
380& 520 & 350 & 800\cr } \right )
\]
Si las ventas del mes de julio se reescriben en la matriz \(J\), ¿con qué matriz se describen las ventas de helados de los dos meses de verano?}
{matriz \(J+A\)}{matriz \(J-A\)}{matriz \(J \cdot A\)}{matriz \(2J+2A\)}{}{}}
\End

\Data\ID{2000019302}
\Updated{2024-03-03 09:03:56}
\Level{A}
\SubArea{Matrices and determinants}
\LANGen{
\Question{
Three ice cream stands of ICE company reported their July sales of four ice cream flavors in number of portions sold. All data can be seen in the table below:
\[ \begin{array}{|c|c|c|c|c|} \hline  &\text{vanilla} & \text{chocolate} & \text{nut} & \text{strawberry} \\\hline
\text{Stand 1}& 720 & 800 & 1\,200&360 \\\hline 
\text{Stand 2} & 550 & 434 & 900 & 300 \\\hline
\text{Stand 3} &610 &300 & 200 & 750 \\\hline \end{array}\]

The profits from the sales of each specific flavor are expressed by the matrix 
\( P=
\left (\MATRIX{ 
1\cr 
1\cr 
3\cr 
2\cr } \right )
\).

Let the July sales be rewritten to the matrix \(J\). By what matrix the profits of individual ice cream stands in July are described?}
{\(J\cdot P\)}{\(P \cdot J\)}{\(J +P\)}{Profits cannot be determined by any of the matrix operations.}{}{}}
\LANGpl{
\Question{
Trzy stoiska z lodami firmy ICE odnotowały lipcową sprzedaż czterech smaków lodów w ilości sprzedanych porcji. Wszystkie dane są widoczne w poniższej tabeli:
\[ \begin{array}{|c|c|c|c|c|} \hline  &\text{wanilia} & \text{czekolada} & \text{orzech} & \text{truskawka} \\\hline
\text{Stoisko 1}& 720 & 800 & 1\,200&360 \\\hline 
\text{Stoisko 2} & 550 & 434 & 900 & 300 \\\hline
\text{Stoisko 3} &610 &300 & 200 & 750 \\\hline \end{array}\]

Zyski ze sprzedaży poszczególnych smaków są wyrażone za pomocą matrycy 
\( P=
\left (\MATRIX{ 
1\cr 
1\cr 
3\cr 
2\cr } \right )
\).

 Jeśli lipcowe wyprzedaże zostaną przepisane do matrycy \(J\), to jaką matrycą opiszemy zyski poszczególnych stoisk w lipcu?}
{\(J\cdot P\)}{\(P \cdot J\)}{\(J +P\)}{Zysków nie można określić przez żadną z operacji macierzowych.}{}{}}
\LANGcs{
\Question{
Tři zmrzlinové stánky firmy ICE hlásily za červenec prodej porcí jednotlivých druhů zmrzliny. Údaje jsou uvedeny v následující tabulce:
\[ \begin{array}{|c|c|c|c|c|} \hline  &\text{vanilka} & \text{čokoláda} & \text{oříšek} & \text{jahoda} \\\hline
\text{Stánek 1}& 720 & 800 & 1\,200&360 \\\hline 
\text{Stánek 2} & 550 & 434 & 900 & 300 \\\hline
\text{Stánek 3} &610 &300 & 200 & 750 \\\hline \end{array}\]

Zisk z prodeje jednoho kopečku zmrzliny v závislosti na jejím druhu (v korunách) vyjadřuje matice 
\( P=
\left (\MATRIX{ 
1\cr 
1\cr 
3\cr 
2\cr } \right )
\).

Matici příslušnou prodeji za červenec označme písmenem \(J\). Která matice vyjadřuje celkový zisk jednotlivých stánků firmy ICE za červenec?}
{\(J\cdot P\)}{\(P \cdot J\)}{\(J +P\)}{Zisk nelze určit žádnou operací s maticemi.}{}{}}
\LANGsk{
\Question{
Tri zmrzlinové stánky firmy ICE hlásili za júl predaj porcií jednotlivých druhov zmrzliny. Údaje sú uvedené v nasledujúcej tabuľke:
\[ \begin{array}{|c|c|c|c|c|} \hline &\text{vanilka} & \text{čokoláda} & \text{oriešok} & \text{jahoda} \\\hline
\text{Stánok 1}& 720 & 800 & 1\,200&360 \\\hline
\text{Stánok 2} & 550 & 434 & 900 & 300 \\\hline
\text{Stánok 3} &610 &300 & 200 & 750 \\\hline \end{array}\]

Zisk z predaja jedného kopčeka zmrzliny v závislosti od jej druhu (v korunách) vyjadruje matice
\( P=
\left (\MATRIX{
1\cr
1\cr
3\cr
2\cr } \right )
\).

Maticu príslušnú predaji za júl označme písmenom \(J\). Ktorá matica vyjadruje celkový zisk jednotlivých stánkov firmy ICE za júl?}
{\(J\cdot P\)}{\(P \cdot J\)}{\(J +P\)}{Zisky nemožno určiť žiadnou z maticových operácií.}{}{}}
\LANGes{
\Question{
Tres puestos de helados de la empresa ICE informaron de sus ventas de julio de cuatro sabores de helado en número de porciones vendidas. Todos los datos en coronas checas pueden verse en la siguiente tabla:
\[ \begin{array}{|c|c|c|c|c|} \hline  &\text{vainilla} & \text{chocolate} & \text{nuez} & \text{fresa} \\\hline
\text{Puesto 1}& 720 & 800 & 1\,200&360 \\\hline 
\text{Puesto 2} & 550 & 434 & 900 & 300 \\\hline
\text{Puesto 3} &610 &300 & 200 & 750 \\\hline \end{array}\]

Los beneficios de las ventas de cada sabor específico se expresan mediante la matriz 
\( P=
\left (\MATRIX{ 
1\cr 
1\cr 
3\cr 
2\cr } \right )
\).

Si se reescriben las ventas de julio en la matriz \(J\). ¿Con qué matriz se describen los beneficios de los puestos de helados individuales en julio?}
{\(J\cdot P\)}{\(P \cdot J\)}{\(J +P\)}{Los beneficios no pueden ser determinados por ninguna de las operaciones con matrices.}{}{}}
\End

\Data\ID{2000019303}
\Updated{2024-03-03 09:03:56}
\Level{A}
\SubArea{Matrices and determinants}
\LANGen{
\Question{
Three ice cream stands of ICE company reported their July sales of four ice cream flavors in number of portions sold. All data can be seen in the table below:
\[ \begin{array}{|c|c|c|c|c|} \hline  &\text{vanilla} & \text{chocolate} & \text{nut} & \text{strawberry} \\\hline
\text{Stand 1}& 720 & 800 & 1\,200&360 \\\hline 
\text{Stand 2} & 550 & 434 & 900 & 300 \\\hline
\text{Stand 3} &610 &300 & 200 & 750 \\\hline \end{array}\]

The profits from the sales of each specific flavor are expressed by the matrix 
\( P=
\left (\MATRIX{ 
1\cr 
1\cr 
3\cr 
2\cr } \right )
\).
Estimate what was the total profit of ICE company in July from all three stands.}
{more than \(12\,000\,\mathrm{CZK}\)}{from \(9\,000\,\mathrm{CZK}\) to \(12\,000\,\mathrm{CZK}\)}{from \(6\,000\,\mathrm{CZK}\) to \(9\,000\,\mathrm{CZK}\)}{less than \(6\,000\,\mathrm{CZK}\)}{}{}}
\LANGpl{
\Question{
Trzy stoiska z lodami firmy ICE odnotoway lipcową sprzedaż czterech smaków lodów w ilości sprzedanych porcji. Wszystkie dane przedstawione są w poniższej tabeli:
\[ \begin{array}{|c|c|c|c|c|} \hline  &\text{wanilia} & \text{czekolada} & \text{orzech} & \text{truskawka} \\\hline
\text{Stoisko 1}& 720 & 800 & 1\,200&360 \\\hline 
\text{Stoisko 2} & 550 & 434 & 900 & 300 \\\hline
\text{Stoisko 3} &610 &300 & 200 & 750 \\\hline \end{array}\]

Zyski ze sprzedaży każdego konkretnego smaku wyrażone są przez matrycę 
\( P=
\left (\MATRIX{ 
1\cr 
1\cr 
3\cr 
2\cr } \right )
\).
Oblicz łączny zysk firmy ICE w lipcu ze wszystkich trzech stoisk.}
{więcej niż \(12\,000\,\mathrm{CZK}\)}{od \(9\,000\,\mathrm{CZK}\) do \(12\,000\,\mathrm{CZK}\)}{od \(6\,000\,\mathrm{CZK}\) do \(9\,000\,\mathrm{CZK}\)}{mniej niż \(6\,000\,\mathrm{CZK}\)}{}{}}
\LANGcs{
\Question{
Tři zmrzlinové stánky firmy ICE hlásily za červenec prodej porcí jednotlivých druhů zmrzliny. Údaje jsou uvedeny v následující tabulce:
\[ \begin{array}{|c|c|c|c|c|} \hline  &\text{vanilka} & \text{čokoláda} & \text{oříšek} & \text{jahoda} \\\hline
\text{Stánek 1}& 720 & 800 & 1\,200&360 \\\hline 
\text{Stánek 2} & 550 & 434 & 900 & 300 \\\hline
\text{Stánek 3} &610 &300 & 200 & 750 \\\hline \end{array}\]

Zisk z prodeje jednoho kopečku zmrzliny v závislosti na jejím druhu (v korunách) vyjadřuje matice
\( P=
\left (\MATRIX{ 
1\cr 
1\cr 
3\cr 
2\cr } \right )
\).
Odhadněte, jaký zisk měla firma ICE za červenec ze všech tří stánků.}
{více než \(12\,000\,\mathrm{CZK}\)}{mezi \(9\,000\,\mathrm{CZK}\) a \(12\,000\,\mathrm{CZK}\)}{mezi \(6\,000\,\mathrm{CZK}\) a \(9\,000\,\mathrm{CZK}\)}{méně než \(6\,000\,\mathrm{CZK}\)}{}{}}
\LANGsk{
\Question{
Tri zmrzlinové stánky firmy ICE hlásili za júl predaj porcií jednotlivých druhov zmrzliny. Údaje sú uvedené v nasledujúcej tabuľke:
\[ \begin{array}{|c|c|c|c|c|} \hline &\text{vanilka} & \text{čokoláda} & \text{oriešok} & \text{jahoda} \\\hline
\text{Stánok 1}& 720 & 800 & 1\,200&360 \\\hline
\text{Stánok 2} & 550 & 434 & 900 & 300 \\\hline
\text{Stánok 3} &610 &300 & 200 & 750 \\\hline \end{array}\]

Zisk z predaja jedného kopčeka zmrzliny v závislosti od jej druhu (v korunách) vyjadruje matica
\( P=
\left (\MATRIX{
1\cr
1\cr
3\cr
2\cr } \right )
\).
Odhadnite, aký zisk mala firma ICE za júl zo všetkých troch stánkov.}
{viac ako \(12\,000\,\mathrm{CZK}\)}{od \(9\,000\,\mathrm{CZK}\) do \(12\,000\,\mathrm{CZK}\)}{od \(6\,000\,\mathrm{CZK}\) do \(9\,000\,\mathrm{CZK}\)}{menej ako \(6\,000\,\mathrm{CZK}\)}{}{}}
\LANGes{
\Question{
Tres puestos de helados de la empresa ICE informaron de sus ventas de julio de cuatro sabores de helado en número de porciones vendidas. Todos los datos en coronas checas pueden verse en la siguiente tabla:
\[ \begin{array}{|c|c|c|c|c|} \hline  &\text{vainilla} & \text{chocolate} & \text{nuez} & \text{fresa} \\\hline
\text{Puesto 1}& 720 & 800 & 1\,200&360 \\\hline 
\text{Puesto 2} & 550 & 434 & 900 & 300 \\\hline
\text{Puesto 3} &610 &300 & 200 & 750 \\\hline \end{array}\]

Los beneficios de las ventas de cada sabor específico se expresan mediante la matriz 
\( P=
\left (\MATRIX{ 
1\cr 
1\cr 
3\cr 
2\cr } \right )
\).
Calcula cuál fue el beneficio total juntando los tres puestos de la empresa ICE en julio.}
{más de \(12\,000\,\mathrm{CZK}\)}{de \(9\,000\,\mathrm{CZK}\) a \(12\,000\,\mathrm{CZK}\)}{de \(6\,000\,\mathrm{CZK}\) a \(9\,000\,\mathrm{CZK}\)}{menos de \(6\,000\,\mathrm{CZK}\)}{}{}}
\End

\Data\ID{2000019304}
\Updated{2024-03-03 09:03:20}
\Level{A}
\SubArea{Matrices and determinants}
\LANGen{
\Question{
Adam’s mom and Peter’s mom compete to see who can buy groceries cheaper. For summer weekend, they both want to buy butter, sugar, flour, and vanilla sugar and there are two stores available Aha and Praha.
\[~\]
The following tables show the quantity of items they both plan to buy and their prices in stores Aha and Praha.
\[ \begin{array}{|c|c|c|c|c|}
\hline  
 &\begin{gathered}\text{Butter}\\\text{(kg)}\end{gathered}
 &\begin{gathered}\text{Sugar}\\\text{(kg)}\end{gathered}
 &\begin{gathered}\text{Flour}\\\text{(kg)}\end{gathered}
 &\begin{gathered}\text{Vanila sugar}\\\text{(number of pieces)}\end{gathered}\\\hline
\text{Adam's mom}& 0.5&  2&1&8 \\\hline
\text{Peter's mom}& 0.25 &2&2&5  \\\hline
\end{array}\]

\[ \begin{array}{|c|c|c|} \hline  &\mathrm{Aha} & \mathrm{Praha}\\\hline
\text{Butter (per }1\,\mathrm{kg)}& 119.6\,\mathrm{CZK}&  159.6\,\mathrm{CZK}  \\\hline 
\text{Sugar (per }1\,\mathrm{kg)} & 12.5\,\mathrm{CZK}& 9.9\,\mathrm{CZK} \\\hline 
\text{Flour (per }1\,\mathrm{kg)} & 10.9\,\mathrm{CZK} & 9.5\,\mathrm{CZK} \\\hline
\text{Vanilla sugar (per } \mathrm{piece)} & 5.9\,\mathrm{CZK} & 5.1\,\mathrm{CZK} \\\hline
\end{array}\]
What follows from the next product of the two matrices?
\[ 
\left (\MATRIX{ 
0.5& 2 & 1 & 8\cr 
0.25& 2 & 2 & 5\cr }
\right )
\cdot
\left (\MATRIX{ 
119.6& 159.6 \cr 
12.5& 9.9 \cr 
10.9& 9.5 \cr 
5.9& 5.1 \cr }
\right )
=
\left (\MATRIX{ 
142.9& 149.9\cr 
106.2& 104.2 \cr }
\right )
\]}
{For Adam’s mom, it is more convenient to shop at Aha, while for Peter’s mom, it is more convenient to shop at Praha.}{For Adam’s mom, it is more convenient to shop at Praha, while for Peter’s mom, it is more convenient to shop at Aha.}{Shopping at Aha is more convenient for both moms.}{Shopping at Praha is more convenient for both moms.}{}{}}
\LANGpl{
\Question{
Mama Adama i mama Piotra rywalizują o to, kto może kupić taniej artykuły spożywcze. Na letni weekend oboje chcą kupić masło, cukier, mąkę i cukier waniliowy, a dostępne są dwa sklepy Aha i Praha.
\[~\]
Poniższe tabele pokazują ilość woreczków, które oboje planują kupić, oraz ich ceny w sklepach Aha i Praha.
\[ \begin{array}{|c|c|c|c|c|}
\hline  
 &\begin{gathered}\text{Masło}\\\text{(kg)}\end{gathered}
 &\begin{gathered}\text{Cukier}\\\text{(kg)}\end{gathered}
 &\begin{gathered}\text{Mąka}\\\text{(kg)}\end{gathered}
 &\begin{gathered}\text{Cukier waniliowy}\\\text{(ilość sztuk)}\end{gathered}\\\hline
\text{mama Adama}& 0{,}5&  2&1&8 \\\hline
\text{mama Piotra}& 0{,}25 &2&2&5  \\\hline
\end{array}\]

\[ \begin{array}{|c|c|c|} \hline  &\mathrm{Aha} & \mathrm{Praha}\\\hline
\text{Masło (za }1\,\mathrm{kg)}& 119{,}6\,\mathrm{CZK}&  159{,}6\,\mathrm{CZK}  \\\hline 
\text{Cukier (za }1\,\mathrm{kg)} & 12{,}5\,\mathrm{CZK}& 9{,}9\,\mathrm{CZK} \\\hline 
\text{Mąka (za }1\,\mathrm{kg)} & 10{,}9\,\mathrm{CZK} & 9{,}5\,\mathrm{CZK} \\\hline
\text{Cukier waniliowy za  sztukę} & 5{,}9\,\mathrm{CZK} & 5{,}1\,\mathrm{CZK} \\\hline
\end{array}\]
Co wynika z kolejnego iloczynu dwóch macierzy?
\[ 
\left (\MATRIX{ 
0{,}5& 2 & 1 & 8\cr 
0{,}25& 2 & 2 & 5\cr }
\right )
\cdot
\left (\MATRIX{ 
119{,}6& 159{,}6 \cr 
12{,}5& 9{,}9 \cr 
10{,}9& 9{,}5 \cr 
5{,}9& 5{,}1 \cr }
\right )
=
\left (\MATRIX{ 
142{,}9& 149{,}9\cr 
106{,}2& 104{,}2 \cr }
\right )
\]}
{Mamie Adama wygodniej robić zakupy w Aha, a mamie Piotra wygodniej robić zakupy w Praha.}{Mamie Adama wygodniej robić zakupy w Praha, a mamie Piotra wygodniej robić zakupy w Aha.}{Zakupy w Aha są wygodniejsze dla obu mam.}{Zakupy w Praha są wygodniejsze dla obu mam.}{}{}}
\LANGcs{
\Question{
Adamova maminka (AM) a Petrova maminka (PM) se předhánějí v tom, která umí nakupovat levněji. Na letní víkend chtějí obě nakoupit máslo, cukr, mouku a vanilkový cukr. K dispozici mají dva obchody Aha a Praha. 
\[~\]
Následující tabulky vyjadřují, kolik kusů jednotlivých surovin mají maminky v plánu nakoupit a ceny v obchodech Aha a Praha.
\[ \begin{array}{|c|c|c|c|c|}
\hline  
 &\begin{gathered}\text{Máslo}\\\text{(kg)}\end{gathered}
 &\begin{gathered}\text{Cukr}\\\text{(kg)}\end{gathered}
 &\begin{gathered}\text{Mouka}\\\text{(kg)}\end{gathered}
 &\begin{gathered}\text{Vanilkový cukr}\\\text{(počet sáčků)}\end{gathered}\\\hline
\text{AM}& 0{,}5&  2&1&8 \\\hline
\text{PM}& 0{,}25 &2&2&5  \\\hline
\end{array}\]

\[ \begin{array}{|c|c|c|} \hline  &\mathrm{Aha} & \mathrm{Praha}\\\hline
\text{Máslo (za }1\,\mathrm{kg)}& 119{,}6\,\mathrm{CZK}&  159{,}6\,\mathrm{CZK}  \\\hline 
\text{Cukr (za }1\,\mathrm{kg)} & 12{,}5\,\mathrm{CZK}& 9{,}9\,\mathrm{CZK} \\\hline 
\text{Mouka (za }1\,\mathrm{kg)} & 10{,}9\,\mathrm{CZK} & 9{,}5\,\mathrm{CZK} \\\hline
\text{Vanilkový cukr (za }\mathrm{kus)} & 5{,}9\,\mathrm{CZK} & 5{,}1\,\mathrm{CZK} \\\hline
\end{array}\]
Co vyplývá z následujícího součinu matic?
\[ 
\left (\MATRIX{ 
0{,}5& 2 & 1 & 8\cr 
0{,}25& 2 & 2 & 5\cr }
\right )
\cdot
\left (\MATRIX{ 
119{,}6& 159{,}6 \cr 
12{,}5& 9{,}9 \cr 
10{,}9& 9{,}5 \cr 
5{,}9& 5{,}1 \cr }
\right )
=
\left (\MATRIX{ 
142{,}9& 149{,}9\cr 
106{,}2& 104{,}2 \cr }
\right )
\]}
{Pro Adamovu maminku je výhodnější nákup v obchodě Aha, pro Petrovu maminku je výhodnější nákup v obchodě Praha.}{Pro Adamovu maminku je výhodnější nákup v obchodě Praha, pro Petrovu maminku je výhodnější nákup v obchodě Aha.}{Pro obě maminky je výhodnější nákup v obchodě Aha.}{Pro obě maminky je výhodnější nákup v obchodě Praha.}{}{}}
\LANGsk{
\Question{
Mamičky Adama a Petra sa predbiehajú v tom, ktorá vie nakupovať lacnejšie. Na letný víkend chcú obe nakúpiť maslo, cukor, múku a vanilkový cukor. K dispozícii majú dva obchody Aha a Praha.
\[~\]
Nasledujúce tabuľky vyjadrujú, koľko kusov jednotlivých surovín majú mamičky v pláne nakúpiť a ceny v obchodoch Aha a Praha.
\[ \begin{array}{|c|c|c|c|c|}
\hline
 &\begin{gathered}\text{Maslo}\\\text{(kg)}\end{gathered}
 &\begin{gathered}\text{Cukor}\\\text{(kg)}\end{gathered}
 &\begin{gathered}\text{Múka}\\\text{(kg)}\end{gathered}
 &\begin{gathered}\text{Vanilkový cukor}\\\text{(počet vreciek)}\end{gathered}\\\hline
\text{Adamova mamička}& 0{,}5& 2&1&8 \\\hline
\text{Petrova mamička}& 0{,}25 &2&2&5 \\\hline
\end{array}\]

\[ \begin{array}{|c|c|c|} \hline  &\mathrm{Aha} & \mathrm{Praha}\\\hline
\text{Maslo (za }1\,\mathrm{kg)}& 119{,}6\,\mathrm{CZK}&  159{,}6\,\mathrm{CZK}  \\\hline 
\text{Cukor (za }1\,\mathrm{kg)} & 12{,}5\,\mathrm{CZK}& 9{,}9\,\mathrm{CZK} \\\hline 
\text{Múka (za }1\,\mathrm{kg)} & 10{,}9\,\mathrm{CZK} & 9{,}5\,\mathrm{CZK} \\\hline
\text{Vanilkový cukor (za }\mathrm{vrecko)} & 5{,}9\,\mathrm{CZK} & 5{,}1\,\mathrm{CZK} \\\hline

\end{array}\]
Čo vyplýva z nasledujúceho súčinu matíc?
\[
\left (\MATRIX{
0{,}5& 2 & 1 & 8\cr
0{,}25& 2 & 2 & 5\cr }
\right )
\cdot
\left (\MATRIX{
119{,}6& 159{,}6 \cr
12{,}5& 9{,}9 \cr
10{,}9& 9{,}5 \cr
5{,}9& 5{,}1 \cr }
\right )
=
\left (\MATRIX{
142{,}9& 149{,}9\cr
106{,}2& 104{,}2 \cr }
\right )
\]}
{Pre Adamovu mamu je výhodnejšie nakupovať v Aha, pre Petrovu mamu zase v Prahe.}{Pre Adamovu mamu je výhodnejšie nakupovať v Prahe, pre Petrovu mamu zase v Aha.}{Nakupovanie v Aha je výhodnejšie pre obe mamičky.}{Nakupovanie v Prahe je výhodnejšie pre obe mamičky.}{}{}}
\LANGes{
\Question{
La madre de Adam (MdA) y la madre de Peter (MdP) compiten para ver quién puede comprar alimentos más baratos. Para un fin de semana de verano, ambas quieren comprar mantequilla, azúcar, harina y esencia de vainilla y hay dos tiendas disponibles: Aha y Praha.
\[~\]
Las siguientes tablas muestran la cantidad de artículos que ambas planean comprar y sus precios en coronas checas en las tiendas de Aha y Praha.
\[ \begin{array}{|c|c|c|c|c|}
\hline  
 &\begin{gathered}\text{Mantequilla}\\\text{(kg)}\end{gathered}
 &\begin{gathered}\text{Azúcar}\\\text{(kg)}\end{gathered}
 &\begin{gathered}\text{Harina}\\\text{(kg)}\end{gathered}
 &\begin{gathered}\text{Esencia de vainilla}\\\text{(número de piezas)}\end{gathered}\\\hline
\text{MdA}& 0.5&  2&1&8 \\\hline
\text{MdP}& 0.25 &2&2&5  \\\hline
\end{array}\]

\[ \begin{array}{|c|c|c|} \hline  &\mathrm{Aha} & \mathrm{Praha}\\\hline
\text{Mantequilla (per }1\,\mathrm{kg)}& 119.6\,\mathrm{CZK}&  159.6\,\mathrm{CZK}  \\\hline 
\text{Azúcar (per }1\,\mathrm{kg)} & 12.5\,\mathrm{CZK}& 9.9\,\mathrm{CZK} \\\hline 
\text{Harina (per }1\,\mathrm{kg)} & 10.9\,\mathrm{CZK} & 9.5\,\mathrm{CZK} \\\hline
\text{Esencia de vainilla (per }\mathrm{pieza)} & 5.9\,\mathrm{CZK} & 5.1\,\mathrm{CZK} \\\hline
\end{array}\]
¿Qué se deduce del siguiente producto de matrices?
\[ 
\left (\MATRIX{ 
0.5& 2 & 1 & 8\cr 
0.25& 2 & 2 & 5\cr }
\right )
\cdot
\left (\MATRIX{ 
119.6& 159.6 \cr 
12.5& 9.9 \cr 
10.9& 9.5 \cr 
5.9& 5.1 \cr }
\right )
=
\left (\MATRIX{ 
142.9& 149.9\cr 
106.2& 104.2 \cr }
\right )
\]}
{Para la madre de Adam, es más conveniente comprar en Aha, mientras que para la madre de Peter, es más conveniente comprar en Praha.}{Para la madre de Adam, es más conveniente comprar en Praha, mientras que para la madre de Peter, es más conveniente comprar en Aha.}{Comprar en Aha es más conveniente para las dos madres.}{Comprar en Praha es más conveniente para las dos madres.}{}{}}
\End

\Data\ID{2000019305}
\Updated{2024-05-10 10:05:06}
\Level{A}
\SubArea{Matrices and determinants}
\LANGen{
\Question{
Mr. Wise buys fuel at three various gas stations – MOL, Shell and EuroOil. He always buys \(25\) liters of diesel, cafe latte and \(2\) nougat croissants.
\[~\]
Prices in CZK:
\[ \begin{array}{|c|c|c|c|} \hline  &\text{Diesel (}1~\text{liter)} &1~\text{Cafe Latte}&1~\text{Croissant}\\\hline
\text{MOL}& 31.20&  57&16.90\\\hline 
\text{Shell}& 27.20 &52 &20  \\\hline 
\text{EuroOil}& 29.60 &49 &18.20 \\\hline
\end{array}\]
Consider the next product of the two matrices:
\[ 
\left (\MATRIX{ 
31.20& 57 & 16.90 \cr
27.20& 52 & 20 \cr
29.60& 49 & 18.20 \cr }
\right )
\cdot
\left (\MATRIX{ 
25 \cr 
1 \cr 
2 \cr 
 }
\right )
\]
Choose the statement which is NOT true:}
{The most optimal prices are at the MOL gas station.}{The most optimal prices are at the Shell gas station.}{The least optimal prices are at the MOL gas station.}{}{}{}}
\LANGpl{
\Question{
Pan Wise kupuje paliwo na trzech różnych stacjach benzynowych – MOL, Shell i EuroOil. Zawsze kupuje \(25\) litrów oleju napędowego, cafe latte i \(2\) croissanty nugatowe.
\[~\]
Ceny w CZK:
\[ \begin{array}{|c|c|c|c|} \hline  &\text{Olej napędowy (}1~\text{litr)} &1~\text{Cafe Latte}&1~\text{croissant}\\\hline
\text{MOL}& 31{,}20&  57&16{,}90\\\hline 
\text{Shell}& 27{,}20 &52 &20  \\\hline 
\text{EuroOil}& 29{,}60 &49 &18{,}20 \\\hline
\end{array}\]
Rozważ następny iloczyn dwóch macierzy:
\[ 
\left (\MATRIX{ 
31{,}20& 57 & 16{,}90 \cr
27{,}20& 52 & 20 \cr
29{,}60& 49 & 18{,}20 \cr }
\right )
\cdot
\left (\MATRIX{ 
25 \cr 
1 \cr 
2 \cr 
 }
\right )
\]
Wybierz stwierdzenie, które NIE jest prawdziwe:}
{Najbardziej optymalne ceny są na stacji benzynowej MOL.}{Najbardziej optymalne ceny są na stacji paliw Shell.}{Najmniej optymalne ceny są na stacji benzynowej MOL.}{}{}{}}
\LANGcs{
\Question{
Pan Moudrý jezdí pro palivo na tři různé čerpací stanice MOL, Shell a EuroOil. Pokaždé koupí  \(25\) litrů nafty, cafe latte a \(2\) nugátové croissanty.
\[~\]
Ceny v CZK:
\[ \begin{array}{|c|c|c|c|} \hline  &\text{Nafta (}1~\text{litr)} &1~\text{Cafe Latte}&1~\text{Croissant}\\\hline
\text{MOL}& 31{,}20&  57&16{,}90\\\hline 
\text{Shell}& 27{,}20 &52 &20  \\\hline 
\text{EuroOil}& 29{,}60 &49 &18{,}20 \\\hline
\end{array}\]
Uvažujme následující součin dvou matic:
\[ 
\left (\MATRIX{ 
31{,}20& 57 & 16{,}90 \cr
27{,}20& 52 & 20 \cr
29{,}60& 49 & 18{,}20 \cr }
\right )
\cdot
\left (\MATRIX{ 
25 \cr 
1 \cr 
2 \cr 
 }
\right )
\]
Vyberte tvrzení, které NENÍ pravdivé.}
{Nejvýhodnější je nákup na čerpací stanici  MOL.}{Nejvýhodnější je nákup na čerpací stanici Shell.}{Nejméně výhodný je nákup na čerpací stanici  MOL.}{}{}{}}
\LANGsk{
\Question{
Pán Múdry jazdí pre palivo na tri rôzne čerpacie stanice MOL, Shell a EuroOil. Zakaždým kúpia \(25\) litrov nafty, cafe latte a \(2\) nugátové croissanty.
\[~\]
Ceny v EUR:
\[ \begin{array}{|c|c|c|c|} \hline &\text{Nafta (}1~\text{liter)} &1~\text{Cafe Latte}&1~\text{ Croissant}\\\hline
\text{MOL}& 31{,}20& 57&16{,}90\\\hline
\text{Shell}& 27{,}20 &52 &20 \\\hline
\text{EuroOil}& 29{,}60 &49 &18{,}20 \\\hline
\end{array}\]
Zvažujme nasledujúci súčin dvoch matíc:
\[
\left (\MATRIX{
31{,}20& 57 & 16{,}90 \cr
27{,}20& 52 & 20 \cr
29{,}60& 49 & 18{,}20 \cr }
\right )
\cdot
\left (\MATRIX{
25 \cr
1 \cr
2 \cr
  }
\right )
\]
Vyberte tvrdenie, ktoré NIE JE pravdivé.}
{Najoptimálnejšie ceny sú na čerpacej stanici MOL.}{Najoptimálnejšie ceny sú na čerpacej stanici Shell.}{Najmenej optimálne ceny sú na čerpacej stanici MOL.}{}{}{}}
\LANGes{
\Question{
El Sr. Wise compra combustible en tres gasolineras distintas – MOL, Shell y EuroOil. Siempre compra \(25\) litros de diesel, un café con leche y \(2\) cruasanes de turrón.
\[~\]
Precios en CZK:
\[ \begin{array}{|c|c|c|c|} \hline  &\text{Diesel (}1~\text{litro)} &1~\text{Café con leche}&1~\text{Cruasán}\\\hline
\text{MOL}& 31.20&  57&16.90\\\hline 
\text{Shell}& 27.20 &52 &20  \\\hline 
\text{EuroOil}& 29.60 &49 &18.20 \\\hline
\end{array}\]
Considera el siguiente producto de dos matrices:
\[ 
\left (\MATRIX{ 
31.20& 57 & 16.90 \cr
27.20& 52 & 20 \cr
29.60& 49 & 18.20 \cr }
\right )
\cdot
\left (\MATRIX{ 
25 \cr 
1 \cr 
2 \cr 
 }
\right )
\]
Elige la afirmación que NO es verdadera:}
{Los precios más óptimos están en la gasolinera MOL.}{Los precios más óptimos están en la gasolinera Shell.}{Los precios menos óptimos están en la gasolinera MOL.}{}{}{}}
\End

\Data\ID{2010006701}
\Updated{2024-03-03 09:03:26}
\Level{A}
\SubArea{Matrices and determinants}
\LANGen{
\Question{
Identify a true statement related to the following matrix
\(A\).
\[ 
A = \left (\MATRIX{ 
2& 4 & -3& 7\cr 
 9 & -5 & -1 & 8 \cr 
11& 0 & 8& 12 \cr 
-7 & -8 & 1& 13 \cr
9& 10 & -6& 2
 } \right )
\]}
{\(A\) is
a \(5\times 4\) matrix
and \(a_{(2,\, 3)} = -1\).}{\(A\) is
a \(5\times 4\) matrix
and \(a_{(2,\, 3)} = 0\).}{\(A\) is
a \(4\times 5\) matrix
and \(a_{(2,\, 3)} = 0\).}{\(A\) is
a \(4\times 5\) matrix
and \(a_{(2,\, 3)} = -1\).}{}{}}
\LANGpl{
\Question{
Wskaż prawdziwe stwierdzenie dotyczące macierzy 
\(A\).
\[ 
A = \left (\MATRIX{ 
2& 4 & -3& 7\cr 
 9 & -5 & -1 & 8 \cr 
11& 0 & 8& 12 \cr 
-7 & -8 & 1& 13 \cr
9& 10 & -6& 2
 } \right )
\]}
{\(A\) jest macierzą \(5\times 4\) 
i \(a_{(2,\, 3)} = -1\).}{\(A\) jest macierzą \(5\times 4\) i  \(a_{(2,\, 3)} = 0\).}{\(A\) jest macierzą \(4\times 5\) i  \(a_{(2,\, 3)} = 0\).}{\(A\) jest macierzą \(4\times 5\) i \(a_{(2,\, 3)} = -1\).}{}{}}
\LANGcs{
\Question{
Je dána matice \(A\). Vyberte správné tvrzení.
\[ 
A = \left (\MATRIX{ 
2& 4 & -3& 7\cr 
 9 & -5 & -1 & 8 \cr 
11& 0 & 8& 12 \cr 
-7 & -8 & 1& 13 \cr
9& 10 & -6& 2
 } \right )
\]}
{Matice \(A\) je typu 
 \(5\times 4\) 
a prvek \(a_{(2,\, 3)} = -1\).}{Matice \(A\) je typu
\(5\times 4\) 
a prvek  \(a_{(2,\, 3)} = 0\).}{Matice \(A\) je typu
\(4\times 5\) a prvek \(a_{(2,\, 3)} = 0\).}{Matice \(A\) je typu
 \(4\times 5\) 
a prvek \(a_{(2,\, 3)} = -1\).}{}{}}
\LANGsk{
\Question{
Je daná matica \(A\). Vyberte správne tvrdenie.
\[ 
A = \left (\MATRIX{ 
2& 4 & -3& 7\cr 
 9 & -5 & -1 & 8 \cr 
11& 0 & 8& 12 \cr 
-7 & -8 & 1& 13 \cr
9& 10 & -6& 2
 } \right )
\]}
{Matica \(A\) je typu
\(5\times 4\) 
a prvok \(a_{(2,\, 3)} = -1\).}{Matica \(A\) je typu 
\(5\times 4\) 
a prvok \(a_{(2,\, 3)} = 0\).}{Matice \(A\) je typu
\(4\times 5\) 
a prvok \(a_{(2,\, 3)} = 0\).}{Matica \(A\) je typu
\(4\times 5\) 
a prvok \(a_{(2,\, 3)} = -1\).}{}{}}
\LANGes{
\Question{
Identifica la proposición lógica verdadera relacionada con la siguiente matriz
\(A\).
\[ 
A = \left (\MATRIX{ 
2& 4 & -3& 7\cr 
 9 & -5 & -1 & 8 \cr 
11& 0 & 8& 12 \cr 
-7 & -8 & 1& 13 \cr
9& 10 & -6& 2
 } \right )
\]}
{\(A\) es \(5\times 4\) matriz
y \(a_{(2,\, 3)} = -1\).}{\(A\) es \(5\times 4\) matriz
y \(a_{(2,\, 3)} = 0\).}{\(A\) es \(4\times 5\) matriz
y \(a_{(2,\, 3)} = 0\).}{\(A\) es \(4\times 5\) matriz
y \(a_{(2,\, 3)} = -1\).}{}{}}
\End

\Data\ID{9000019903}
\Updated{2024-03-03 09:03:26}
\Level{A}
\SubArea{Matrices and determinants}
\LANGen{
\Question{
Identify a true statement related to the following matrix
\(A\).
\[ 
A = \left (\MATRIX{ 
-2& 3 & 10& 5 & -5\cr 
 6 & 11 & -7 & 2 & -3
\cr 
-7& 15& -6& 2 & 4\cr 
-8 & 1 & 13 & -5 & 0
 } \right )
\]}
{\(A\) is
a \(4\times 5\) matrix
and \(a_{(3,\, 2)} = 15\).}{\(A\) is
a \(4\times 5\) matrix
and \(a_{(2,\, 3)} = 15\).}{\(A\) is
a \(5\times 4\) matrix
and \(a_{(3,\, 2)} = -7\).}{\(A\) is
a \(5\times 4\) matrix
and \(a_{(3,\, 2)} = 15\).}{}{}}
\LANGpl{
\Question{
Które z poniższych zdań odnoszących się do macierzy
\(A\) jest zgodne z prawdą?
\[ 
A = \left (\MATRIX{ 
-2& 3 & 10& 5 & -5\cr 
 6 & 11 & -7 & 2 & -3
\cr 
-7& 15& -6& 2 & 4\cr 
-8 & 1 & 13 & -5 & 0
 } \right )
\]}
{\(A\) jest 
macierzą \(4\times 5\) 
i \(a_{(3,\, 2)} = 15\).}{\(A\) jest macierzą \(4\times 5\) 
i \(a_{(2,\, 3)} = 15\).}{\(A\) jest macierzą \(5\times 4\) 
i \(a_{(3,\, 2)} = -7\).}{\(A\) jest macierzą \(5\times 4\) 
i \(a_{(3,\, 2)} = 15\).}{}{}}
\LANGcs{
\Question{
Je dána matice \(A\).
Vyberte správné tvrzení.
\[ 
A = \left (\MATRIX{ 
-2& 3 & 10& 5 & -5\cr 
 6 & 11 & -7 & 2 & -3
\cr 
-7& 15& -6& 2 & 4\cr 
-8 & 1 & 13 & -5 & 0
 } \right )
\]}
{Matice je typu \((4,\, 5)\) 
a prvek \(a_{(3,\, 2)} = 15\).}{Matice je typu \((4,\, 5)\) 
a prvek \(a_{(2,\, 3)} = 15\).}{Matice je typu \((5,\, 4)\) 
a prvek \(a_{(3,\, 2)} = -7\).}{Matice je typu \((5,\, 4)\) 
a prvek \(a_{(3,\, 2)} = 15\).}{}{}}
\LANGsk{
\Question{
Je daná matica
\(A\). Vyberte správne tvrdenie.
\[ 
A = \left (\MATRIX{ 
-2& 3 & 10& 5 & -5\cr 
 6 & 11 & -7 & 2 & -3
\cr 
-7& 15& -6& 2 & 4\cr 
-8 & 1 & 13 & -5 & 0
 } \right )
\]}
{Matica \(A\) je typu
 \(4\times 5\) 
a \(a_{(3,\, 2)} = 15\).}{Matica \(A\) je typu
 \(4\times 5\) 
a \(a_{(2,\, 3)} = 15\).}{Matica \(A\) je typu
 \(5\times 4\) 
a \(a_{(3,\, 2)} = -7\).}{Matica \(A\) je typu
 \(5\times 4\) 
a \(a_{(3,\, 2)} = 15\).}{}{}}
\LANGes{
\Question{
Identifica la proposición lógica verdadera relacionada con la siguiente matriz
\(A\).
\[ 
A = \left (\MATRIX{ 
-2& 3 & 10& 5 & -5\cr 
 6 & 11 & -7 & 2 & -3
\cr 
-7& 15& -6& 2 & 4\cr 
-8 & 1 & 13 & -5 & 0
 } \right )
\]}
{\(A\) es \(4\times 5\) matriz
y \(a_{(3,\, 2)} = 15\).}{\(A\) es \(4\times 5\) matriz
y \(a_{(2,\, 3)} = 15\).}{\(A\) es \(5\times 4\) matriz
y \(a_{(3,\, 2)} = -7\).}{\(A\) es \(5\times 4\) matriz
y \(a_{(3,\, 2)} = 15\).}{}{}}
\End

\Data\ID{2000017101}
\Updated{2024-03-03 09:03:26}
\Level{B}
\SubArea{Matrices and determinants}
\LANGen{
\Question{
Determine all real numbers \(b\) such that there exists the inverse matrix to the matrix:
\[ 
\left (\MATRIX{ 
4& 3 & b\cr 
 2 & 1 & 2
\cr 
b& b & -1 } \right )
\]}
{all real numbers}{all non-negative real numbers}{all positive real numbers}{Such a number does not exists.}{}{}}
\LANGpl{
\Question{
Wyznacz wszystkie liczby rzeczywiste \(b\) tak aby zaistniała macierz odwrotna do podanej macierzy:
\[ 
\left (\MATRIX{ 
4& 3 & b\cr 
 2 & 1 & 2
\cr 
b& b & -1 } \right )
\]}
{wszystkie liczby rzeczywiste}{wszystkie nieujemne liczby rzeczywiste}{wszystkie dodatnie liczby rzeczywiste}{Taka liczba nie istnieje.}{}{}}
\LANGcs{
\Question{
Určete, pro která reálná čísla \(b\) existuje inverzní matice k matici:
\[ 
\left (\MATRIX{ 
4& 3 & b\cr 
 2 & 1 & 2
\cr 
b& b & -1 } \right )
\]}
{pro všechna reálná čísla}{pro všechna nezáporná reálná čísla}{pro všechna kladná reálná čísla}{Takové číslo neexistuje.}{}{}}
\LANGsk{
\Question{
Určte, pre ktoré reálne čísla \(b\) existuje inverzná matica k matici:
\[ 
\left (\MATRIX{ 
4& 3 & b\cr 
 2 & 1 & 2
\cr 
b& b & -1 } \right )
\]}
{pre všetky reálne čísla}{pre všetky nezáporné reálne čísla}{pre všetky kladné reálne čísla}{Také číslo neexistuje.}{}{}}
\LANGes{
\Question{
Determina todos los números reales \(b\) tales que exista la matriz inversa a la matriz:
\[ 
\left (\MATRIX{ 
4& 3 & b\cr 
 2 & 1 & 2
\cr 
b& b & -1 } \right )
\]}
{todos los números reales}{todos los números reales no negativos}{todos los números reales positivos}{El número así no existe.}{}{}}
\End

\Data\ID{2000017102}
\Updated{2024-03-03 09:03:26}
\Level{B}
\SubArea{Matrices and determinants}
\LANGen{
\Question{
Compute the inverse matrix to the matrix:
\[ 
\left (\MATRIX{ 
4& 3 & 0\cr 
 2 & 1 & 2
\cr 
0& 0 & -1 } \right )
\]}
{\[ 
\left (\MATRIX{ 
-\frac12& \frac32 & 3\cr 
 1 & -2 & -4
\cr 
0& 0 & -1 } \right )
\]}{\[ 
\left (\MATRIX{ 
\frac12& \frac32 & 3\cr 
 1 & -2 & -4
\cr 
0& 0 & 1 } \right )
\]}{\[ 
\left (\MATRIX{ 
-\frac12& \frac32 & -3\cr 
 -1 & -2 & -4
\cr 
0& 0 & -1 } \right )
\]}{\[ 
\left (\MATRIX{ 
-\frac12& \frac32 & 3\cr 
 1 & 2 & -4
\cr 
0& 0 & -1 } \right )
\]}{}{}}
\LANGpl{
\Question{
Oblicz macierz odwrotną do macierzy:
\[ 
\left (\MATRIX{ 
4& 3 & 0\cr 
 2 & 1 & 2
\cr 
0& 0 & -1 } \right )
\]}
{\[ 
\left (\MATRIX{ 
-\frac12& \frac32 & 3\cr 
 1 & -2 & -4
\cr 
0& 0 & -1 } \right )
\]}{\[ 
\left (\MATRIX{ 
\frac12& \frac32 & 3\cr 
 1 & -2 & -4
\cr 
0& 0 & 1 } \right )
\]}{\[ 
\left (\MATRIX{ 
-\frac12& \frac32 & -3\cr 
 -1 & -2 & -4
\cr 
0& 0 & -1 } \right )
\]}{\[ 
\left (\MATRIX{ 
-\frac12& \frac32 & 3\cr 
 1 & 2 & -4
\cr 
0& 0 & -1 } \right )
\]}{}{}}
\LANGcs{
\Question{
Vypočítejte matici inverzní k matici:
\[ 
\left (\MATRIX{ 
4& 3 & 0\cr 
 2 & 1 & 2
\cr 
0& 0 & -1 } \right )
\]}
{\[ 
\left (\MATRIX{ 
-\frac12& \frac32 & 3\cr 
 1 & -2 & -4
\cr 
0& 0 & -1 } \right )
\]}{\[ 
\left (\MATRIX{ 
\frac12& \frac32 & 3\cr 
 1 & -2 & -4
\cr 
0& 0 & 1 } \right )
\]}{\[ 
\left (\MATRIX{ 
-\frac12& \frac32 & -3\cr 
 -1 & -2 & -4
\cr 
0& 0 & -1 } \right )
\]}{\[ 
\left (\MATRIX{ 
-\frac12& \frac32 & 3\cr 
 1 & 2 & -4
\cr 
0& 0 & -1 } \right )
\]}{}{}}
\LANGsk{
\Question{
Vypočítajte maticu inverznú k matici:
\[ 
\left (\MATRIX{ 
4& 3 & 0\cr 
 2 & 1 & 2
\cr 
0& 0 & -1 } \right )
\]}
{\[ 
\left (\MATRIX{ 
-\frac12& \frac32 & 3\cr 
 1 & -2 & -4
\cr 
0& 0 & -1 } \right )
\]}{\[ 
\left (\MATRIX{ 
\frac12& \frac32 & 3\cr 
 1 & -2 & -4
\cr 
0& 0 & 1 } \right )
\]}{\[ 
\left (\MATRIX{ 
-\frac12& \frac32 & -3\cr 
 -1 & -2 & -4
\cr 
0& 0 & -1 } \right )
\]}{\[ 
\left (\MATRIX{ 
-\frac12& \frac32 & 3\cr 
 1 & 2 & -4
\cr 
0& 0 & -1 } \right )
\]}{}{}}
\LANGes{
\Question{
Calcula la matriz inversa a la matriz:
\[ 
\left (\MATRIX{ 
4& 3 & 0\cr 
 2 & 1 & 2
\cr 
0& 0 & -1 } \right )
\]}
{\[ 
\left (\MATRIX{ 
-\frac12& \frac32 & 3\cr 
 1 & -2 & -4
\cr 
0& 0 & -1 } \right )
\]}{\[ 
\left (\MATRIX{ 
\frac12& \frac32 & 3\cr 
 1 & -2 & -4
\cr 
0& 0 & 1 } \right )
\]}{\[ 
\left (\MATRIX{ 
-\frac12& \frac32 & -3\cr 
 -1 & -2 & -4
\cr 
0& 0 & -1 } \right )
\]}{\[ 
\left (\MATRIX{ 
-\frac12& \frac32 & 3\cr 
 1 & 2 & -4
\cr 
0& 0 & -1 } \right )
\]}{}{}}
\End

\Data\ID{2000017103}
\Updated{2024-03-03 09:03:26}
\Level{B}
\SubArea{Matrices and determinants}
\LANGen{
\Question{
Compute the inverse matrix to the matrix:
\[ 
\left (\MATRIX{ 
1& 0 & 0\cr 
 0 & 1 & 1
\cr 
-1& 0 & 1} \right )
\]}
{\[ 
\left (\MATRIX{ 
1& 0 & 0\cr 
 -1 & 1 & -1
\cr 
1& 0 & 1 } \right )
\]}{\[ 
\left (\MATRIX{ 
1& 0 & 0\cr 
 1 & 1 & 1
\cr 
1& 0 & 1 } \right )
\]}{\[ 
\left (\MATRIX{ 
1& 0 & 0\cr 
 1 & 1 & 1
\cr 
-1& 0 & -1 } \right )
\]}{\[ 
\left (\MATRIX{ 
-1& 0 & 0\cr 
 -1 & 1 & -1
\cr 
1& 0 & -1 } \right )
\]}{}{}}
\LANGpl{
\Question{
Oblicz macierz odwrotną do macierzy:
\[ 
\left (\MATRIX{ 
1& 0 & 0\cr 
 0 & 1 & 1
\cr 
-1& 0 & 1} \right )
\]}
{\[ 
\left (\MATRIX{ 
1& 0 & 0\cr 
 -1 & 1 & -1
\cr 
1& 0 & 1 } \right )
\]}{\[ 
\left (\MATRIX{ 
1& 0 & 0\cr 
 1 & 1 & 1
\cr 
1& 0 & 1 } \right )
\]}{\[ 
\left (\MATRIX{ 
1& 0 & 0\cr 
 1 & 1 & 1
\cr 
-1& 0 & -1 } \right )
\]}{\[ 
\left (\MATRIX{ 
-1& 0 & 0\cr 
 -1 & 1 & -1
\cr 
1& 0 & -1 } \right )
\]}{}{}}
\LANGcs{
\Question{
Vypočítejte matici inverzní k matici:
\[ 
\left (\MATRIX{ 
1& 0 & 0\cr 
 0 & 1 & 1
\cr 
-1& 0 & 1} \right )
\]}
{\[ 
\left (\MATRIX{ 
1& 0 & 0\cr 
 -1 & 1 & -1
\cr 
1& 0 & 1 } \right )
\]}{\[ 
\left (\MATRIX{ 
1& 0 & 0\cr 
 1 & 1 & 1
\cr 
1& 0 & 1 } \right )
\]}{\[ 
\left (\MATRIX{ 
1& 0 & 0\cr 
 1 & 1 & 1
\cr 
-1& 0 & -1 } \right )
\]}{\[ 
\left (\MATRIX{ 
-1& 0 & 0\cr 
 -1 & 1 & -1
\cr 
1& 0 & -1 } \right )
\]}{}{}}
\LANGsk{
\Question{
Vypočítajte maticu inverznú k matici:
\[ 
\left (\MATRIX{ 
1& 0 & 0\cr 
 0 & 1 & 1
\cr 
-1& 0 & 1} \right )
\]}
{\[ 
\left (\MATRIX{ 
1& 0 & 0\cr 
 -1 & 1 & -1
\cr 
1& 0 & 1 } \right )
\]}{\[ 
\left (\MATRIX{ 
1& 0 & 0\cr 
 1 & 1 & 1
\cr 
1& 0 & 1 } \right )
\]}{\[ 
\left (\MATRIX{ 
1& 0 & 0\cr 
 1 & 1 & 1
\cr 
-1& 0 & -1 } \right )
\]}{\[ 
\left (\MATRIX{ 
-1& 0 & 0\cr 
 -1 & 1 & -1
\cr 
1& 0 & -1 } \right )
\]}{}{}}
\LANGes{
\Question{
Calcula la matriz inversa a la matriz:
\[ 
\left (\MATRIX{ 
1& 0 & 0\cr 
 0 & 1 & 1
\cr 
-1& 0 & 1} \right )
\]}
{\[ 
\left (\MATRIX{ 
1& 0 & 0\cr 
 -1 & 1 & -1
\cr 
1& 0 & 1 } \right )
\]}{\[ 
\left (\MATRIX{ 
1& 0 & 0\cr 
 1 & 1 & 1
\cr 
1& 0 & 1 } \right )
\]}{\[ 
\left (\MATRIX{ 
1& 0 & 0\cr 
 1 & 1 & 1
\cr 
-1& 0 & -1 } \right )
\]}{\[ 
\left (\MATRIX{ 
-1& 0 & 0\cr 
 -1 & 1 & -1
\cr 
1& 0 & -1 } \right )
\]}{}{}}
\End

\Data\ID{2000017104}
\Updated{2024-03-03 09:03:26}
\Level{B}
\SubArea{Matrices and determinants}
\LANGen{
\Question{
Find the number \(a\) so that the matrices given below are inverse to each other. 
\[
\left (\MATRIX{ 
1& 2 \cr 
 2 & 2 \cr} \right ),
~
\left (\MATRIX{ 
-1& 1 \cr 
 1 & a \cr} \right )
\]}
{\(a=-\frac12\)}{\(a=\frac12\)}{\(a=0\)}{\(a\) cannot be found.}{}{}}
\LANGpl{
\Question{
Znajdź liczbę \(a\) tak, aby macierze podane poniżej były odwrotne względem siebie.
\[
\left (\MATRIX{ 
1& 2 \cr 
 2 & 2 \cr} \right ),
~
\left (\MATRIX{ 
-1& 1 \cr 
 1 & a \cr} \right )
\]}
{\(a=-\frac12\)}{\(a=\frac12\)}{\(a=0\)}{liczba \(a\) nie istnieje.}{}{}}
\LANGcs{
\Question{
Určete, pro které číslo \(a\) budou uvedené matice navzájem inverzními maticemi. 
\[
\left (\MATRIX{ 
1& 2 \cr 
 2 & 2 \cr} \right ),
~
\left (\MATRIX{ 
-1& 1 \cr 
 1 & a \cr} \right )
\]}
{\(a=-\frac12\)}{\(a=\frac12\)}{\(a=0\)}{\(a\) nelze určit.}{}{}}
\LANGsk{
\Question{
Určte, pre ktoré číslo \(a\) budú uvedené matice navzájom inverznými maticami. 
\[
\left (\MATRIX{ 
1& 2 \cr 
 2 & 2 \cr} \right ),
~
\left (\MATRIX{ 
-1& 1 \cr 
 1 & a \cr} \right )
\]}
{\(a=-\frac12\)}{\(a=\frac12\)}{\(a=0\)}{\(a\) nie je možné určiť.}{}{}}
\LANGes{
\Question{
Halla el número \(a\) para que las matrices dadas sean inversas.
\[
\left (\MATRIX{ 
1& 2 \cr 
 2 & 2 \cr} \right ),
~
\left (\MATRIX{ 
-1& 1 \cr 
 1 & a \cr} \right )
\]}
{\(a=-\frac12\)}{\(a=\frac12\)}{\(a=0\)}{\(a\) no se puede hallar.}{}{}}
\End

\Data\ID{2000017105}
\Updated{2024-03-03 09:03:26}
\Level{B}
\SubArea{Matrices and determinants}
\LANGen{
\Question{
Let \(x\) be a real number. Determine the inverse matrix to the matrix:
\[
\left (\MATRIX{ 
\cos x& -\sin x \cr 
 \sin x & \cos x \cr} \right )
\]}
{\[
\left (\MATRIX{ 
\cos x& \sin x \cr 
 -\sin x & \cos x \cr} \right )
\]}{\[
\left (\MATRIX{ 
1& 0 \cr 
 0 & 1 \cr} \right )
\]}{\[
\left (\MATRIX{ 
\frac{\cos x}{\sin 2x}& \frac{\sin x}{\sin 2x}\cr 
 -\frac{\sin x}{\sin 2x}& \frac{\cos x}{\sin 2x} \cr} \right )
\]}{The inverse matrix does not exist.}{}{}}
\LANGpl{
\Question{
Niech \(x\) będzie liczbą rzeczywistą. Wyznacz macierz odwrotną do macierzy:
\[
\left (\MATRIX{ 
\cos x& -\sin x \cr 
 \sin x & \cos x \cr} \right )
\]}
{\[
\left (\MATRIX{ 
\cos x& \sin x \cr 
 -\sin x & \cos x \cr} \right )
\]}{\[
\left (\MATRIX{ 
1& 0 \cr 
 0 & 1 \cr} \right )
\]}{\[
\left (\MATRIX{ 
\frac{\cos x}{\sin 2x}& \frac{\sin x}{\sin 2x}\cr 
 -\frac{\sin x}{\sin 2x}& \frac{\cos x}{\sin 2x} \cr} \right )
\]}{Macierz odwrotna nie istnieje.}{}{}}
\LANGcs{
\Question{
Nechť \(x\) je reálné číslo. Určete matici inverzní k matici:
\[
\left (\MATRIX{ 
\cos x& -\sin x \cr 
 \sin x & \cos x \cr} \right )
\]}
{\[
\left (\MATRIX{ 
\cos x& \sin x \cr 
 -\sin x & \cos x \cr} \right )
\]}{\[
\left (\MATRIX{ 
1& 0 \cr 
 0 & 1 \cr} \right )
\]}{\[
\left (\MATRIX{ 
\frac{\cos x}{\sin 2x}& \frac{\sin x}{\sin 2x}\cr 
 -\frac{\sin x}{\sin 2x}& \frac{\cos x}{\sin 2x} \cr} \right )
\]}{Inverzní matice neexistuje.}{}{}}
\LANGsk{
\Question{
Nech \(x\) je reálne číslo. Určte maticu inverznú k matici:
\[
\left (\MATRIX{ 
\cos x& -\sin x \cr 
 \sin x & \cos x \cr} \right )
\]}
{\[
\left (\MATRIX{ 
\cos x& \sin x \cr 
 -\sin x & \cos x \cr} \right )
\]}{\[
\left (\MATRIX{ 
1& 0 \cr 
 0 & 1 \cr} \right )
\]}{\[
\left (\MATRIX{ 
\frac{\cos x}{\sin 2x}& \frac{\sin x}{\sin 2x}\cr 
 -\frac{\sin x}{\sin 2x}& \frac{\cos x}{\sin 2x} \cr} \right )
\]}{Inverzná matica neexistuje.}{}{}}
\LANGes{
\Question{
Suponiendo que \(x\) es un número real, halla la matriz inversa a la matriz:
\[
\left (\MATRIX{ 
\cos x& -\sin x \cr 
 \sin x & \cos x \cr} \right )
\]}
{\[
\left (\MATRIX{ 
\cos x& \sin x \cr 
 -\sin x & \cos x \cr} \right )
\]}{\[
\left (\MATRIX{ 
1& 0 \cr 
 0 & 1 \cr} \right )
\]}{\[
\left (\MATRIX{ 
\frac{\cos x}{\sin 2x}& \frac{\sin x}{\sin 2x}\cr 
 -\frac{\sin x}{\sin 2x}& \frac{\cos x}{\sin 2x} \cr} \right )
\]}{La matriz inversa no existe.}{}{}}
\End

\Data\ID{2000017106}
\Updated{2024-03-03 09:03:26}
\Level{B}
\SubArea{Matrices and determinants}
\LANGen{
\Question{
For what values of the real parameters \(a\), \(b\), \(c\), \(d\), \(e\) and \(f\) will the matrices given below be inverse to each other?
\[ 
\left (\MATRIX{ 
1& 0 & 0\cr 
 0 & 1 & 1
\cr 
a& b & c} \right ),  
 \left (\MATRIX{ 
d& e & f\cr 
 -1 & 1 & -1
\cr 
1& 0 & 1} \right )
\]}
{\(a=-1\), \(b=0\), \(c=1\), \(d=1\), \(e=0\), \(f=0\)}{\(a=1\), \(b=0\), \(c=1\), \(d=-1\), \(e=0\), \(f=1\)}{\(a=-1\), \(b=0\), \(c=-1\), \(d=1\), \(e=0\), \(f=0\)}{\(a=1\), \(b=0\), \(c=-1\), \(d=-1\), \(e=1\), \(f=0\)}{}{}}
\LANGpl{
\Question{
Dla jakich wartości parametrów rzeczywistych \(a\), \(b\), \(c\), \(d\), \(e\) i \(f\) macierze podane poniżej będą odwrotne względem siebie?
\[ 
\left (\MATRIX{ 
1& 0 & 0\cr 
 0 & 1 & 1
\cr 
a& b & c} \right ), 
 \left (\MATRIX{ 
d& e & f\cr 
 -1 & 1 & -1
\cr 
1& 0 & 1} \right )
\]}
{\(a=-1\), \(b=0\), \(c=1\), \(d=1\), \(e=0\), \(f=0\)}{\(a=1\), \(b=0\), \(c=1\), \(d=-1\), \(e=0\), \(f=1\)}{\(a=-1\), \(b=0\), \(c=-1\), \(d=1\), \(e=0\), \(f=0\)}{\(a=1\), \(b=0\), \(c=-1\), \(d=-1\), \(e=1\), \(f=0\)}{}{}}
\LANGcs{
\Question{
Pro jaké hodnoty reálných parametrů \(a\), \(b\), \(c\), \(d\), \(e\) a \(f\) budou uvedené matice navzájem inverzní?
\[ 
\left (\MATRIX{ 
1& 0 & 0\cr 
 0 & 1 & 1
\cr 
a& b & c} \right ),  
 \left (\MATRIX{ 
d& e & f\cr 
 -1 & 1 & -1
\cr 
1& 0 & 1} \right )
\]}
{\(a=-1\), \(b=0\), \(c=1\), \(d=1\), \(e=0\), \(f=0\)}{\(a=1\), \(b=0\), \(c=1\), \(d=-1\), \(e=0\), \(f=1\)}{\(a=-1\), \(b=0\), \(c=-1\), \(d=1\), \(e=0\), \(f=0\)}{\(a=1\), \(b=0\), \(c=-1\), \(d=-1\), \(e=1\), \(f=0\)}{}{}}
\LANGsk{
\Question{
Pre aké hodnoty reálnych parametrov a, b, c, d, e a f budú uvedené matice navzájom inverzné?
\[ 
\left (\MATRIX{ 
1& 0 & 0\cr 
 0 & 1 & 1
\cr 
a& b & c} \right ),  
 \left (\MATRIX{ 
d& e & f\cr 
 -1 & 1 & -1
\cr 
1& 0 & 1} \right )
\]}
{\(a=-1\), \(b=0\), \(c=1\), \(d=1\), \(e=0\), \(f=0\)}{\(a=1\), \(b=0\), \(c=1\), \(d=-1\), \(e=0\), \(f=1\)}{\(a=-1\), \(b=0\), \(c=-1\), \(d=1\), \(e=0\), \(f=0\)}{\(a=1\), \(b=0\), \(c=-1\), \(d=-1\), \(e=1\), \(f=0\)}{}{}}
\LANGes{
\Question{
Halla los valores de los parámetros reales \(a\), \(b\), \(c\), \(d\), \(e\) y \(f\) para que las matrices dadas sean inversas.
\[ 
\left (\MATRIX{ 
1& 0 & 0\cr 
 0 & 1 & 1
\cr 
a& b & c} \right ),  
 \left (\MATRIX{ 
d& e & f\cr 
 -1 & 1 & -1
\cr 
1& 0 & 1} \right )
\]}
{\(a=-1\), \(b=0\), \(c=1\), \(d=1\), \(e=0\), \(f=0\)}{\(a=1\), \(b=0\), \(c=1\), \(d=-1\), \(e=0\), \(f=1\)}{\(a=-1\), \(b=0\), \(c=-1\), \(d=1\), \(e=0\), \(f=0\)}{\(a=1\), \(b=0\), \(c=-1\), \(d=-1\), \(e=1\), \(f=0\)}{}{}}
\End

\Data\ID{2000017107}
\Updated{2024-03-03 09:03:26}
\Level{B}
\SubArea{Matrices and determinants}
\LANGen{
\Question{
Compute the inverse matrix to the matrix:
\[ 
\left (\MATRIX{ 
\frac17&  -\frac3{14}\cr 
 \frac27 &  \frac1{14}} \right )
\]}
{\[ 
\left (\MATRIX{ 
1&  3\cr 
-4& 2}\right )
\]}{\[ 
\left (\MATRIX{ 
1&  3\cr 
4& -2}\right )
\]}{\[ 
\left (\MATRIX{ 
1&  3\cr 
2& -4}\right )
\]}{\[ 
\left (\MATRIX{ 
2&  -3\cr 
4& 1}\right )
\]}{}{}}
\LANGpl{
\Question{
Oblicz macierz odwrotną do macierzy:
\[ 
\left (\MATRIX{ 
\frac17&  -\frac3{14}\cr 
 \frac27 &  \frac1{14}} \right )
\]}
{\[ 
\left (\MATRIX{ 
1&  3\cr 
-4& 2}\right )
\]}{\[ 
\left (\MATRIX{ 
1&  3\cr 
4& -2}\right )
\]}{\[ 
\left (\MATRIX{ 
1&  3\cr 
2& -4}\right )
\]}{\[ 
\left (\MATRIX{ 
2&  -3\cr 
4& 1}\right )
\]}{}{}}
\LANGcs{
\Question{
Vypočítejte inverzní matici k matici:
\[ 
\left (\MATRIX{ 
\frac17&  -\frac3{14}\cr 
 \frac27 &  \frac1{14}} \right )
\]}
{\[ 
\left (\MATRIX{ 
1&  3\cr 
-4& 2}\right )
\]}{\[ 
\left (\MATRIX{ 
1&  3\cr 
4& -2}\right )
\]}{\[ 
\left (\MATRIX{ 
1&  3\cr 
2& -4}\right )
\]}{\[ 
\left (\MATRIX{ 
2&  -3\cr 
4& 1}\right )
\]}{}{}}
\LANGsk{
\Question{
Vypočítajte inverznú maticu k matici:
\[ 
\left (\MATRIX{ 
\frac17&  -\frac3{14}\cr 
 \frac27 &  \frac1{14}} \right )
\]}
{\[ 
\left (\MATRIX{ 
1&  3\cr 
-4& 2}\right )
\]}{\[ 
\left (\MATRIX{ 
1&  3\cr 
4& -2}\right )
\]}{\[ 
\left (\MATRIX{ 
1&  3\cr 
2& -4}\right )
\]}{\[ 
\left (\MATRIX{ 
2&  -3\cr 
4& 1}\right )
\]}{}{}}
\LANGes{
\Question{
Calcula la matriz inversa a la matriz:
\[ 
\left (\MATRIX{ 
\frac17&  -\frac3{14}\cr 
 \frac27 &  \frac1{14}} \right )
\]}
{\[ 
\left (\MATRIX{ 
1&  3\cr 
-4& 2}\right )
\]}{\[ 
\left (\MATRIX{ 
1&  3\cr 
4& -2}\right )
\]}{\[ 
\left (\MATRIX{ 
1&  3\cr 
2& -4}\right )
\]}{\[ 
\left (\MATRIX{ 
2&  -3\cr 
4& 1}\right )
\]}{}{}}
\End

\Data\ID{2000017108}
\Updated{2024-03-03 09:03:26}
\Level{B}
\SubArea{Matrices and determinants}
\LANGen{
\Question{
For what values of the real parameters \(a\) and \(b\) will the matrices given below be inverse to each other?
\[
\left (\MATRIX{ 
a& 7 \cr 
 3 & 1 \cr} \right ),
~
\left (\MATRIX{ 
-\frac1{16}& \frac7{16} \cr 
 b & -\frac5{16} \cr} \right )
\]}
{\(a=5\), \(b=\frac3{16}\)}{\(a=-5\), \(b=\frac3{16}\)}{\(a=5\), \(b=-\frac3{16}\)}{\(a=-5\), \(b=-\frac3{16}\)}{}{}}
\LANGpl{
\Question{
Dla jakich wartości parametrów rzeczywistych \(a\) i \(b\) macierze podane poniżej będą odwrotne względem siebie?
\[
\left (\MATRIX{ 
a& 7 \cr 
 3 & 1 \cr} \right ),
~
\left (\MATRIX{ 
-\frac1{16}& \frac7{16} \cr 
 b & -\frac5{16} \cr} \right )
\]}
{\(a=5\), \(b=\frac3{16}\)}{\(a=-5\), \(b=\frac3{16}\)}{\(a=5\), \(b=-\frac3{16}\)}{\(a=-5\), \(b=-\frac3{16}\)}{}{}}
\LANGcs{
\Question{
Pro jaké hodnoty reálných parametrů \(a\) a \(b\) budou níže uvedené matice navzájem inverzní?
\[
\left (\MATRIX{ 
a& 7 \cr 
 3 & 1 \cr} \right ),
~
\left (\MATRIX{ 
-\frac1{16}& \frac7{16} \cr 
 b & -\frac5{16} \cr} \right )
\]}
{\(a=5\), \(b=\frac3{16}\)}{\(a=-5\), \(b=\frac3{16}\)}{\(a=5\), \(b=-\frac3{16}\)}{\(a=-5\), \(b=-\frac3{16}\)}{}{}}
\LANGsk{
\Question{
Pre aké hodnoty reálnych parametrov \(a\) a \(b\) budú nižšie uvedené matice navzájom inverzné?
\[
\left (\MATRIX{ 
a& 7 \cr 
 3 & 1 \cr} \right ),
~
\left (\MATRIX{ 
-\frac1{16}& \frac7{16} \cr 
 b & -\frac5{16} \cr} \right )
\]}
{\(a=5\), \(b=\frac3{16}\)}{\(a=-5\), \(b=\frac3{16}\)}{\(a=5\), \(b=-\frac3{16}\)}{\(a=-5\), \(b=-\frac3{16}\)}{}{}}
\LANGes{
\Question{
Halla los valores de los parámetros reales \(a\) y \(b\) para que las matrices dadas sean inversas.
\[
\left (\MATRIX{ 
a& 7 \cr 
 3 & 1 \cr} \right ),
~
\left (\MATRIX{ 
-\frac1{16}& \frac7{16} \cr 
 b & -\frac5{16} \cr} \right )
\]}
{\(a=5\), \(b=\frac3{16}\)}{\(a=-5\), \(b=\frac3{16}\)}{\(a=5\), \(b=-\frac3{16}\)}{\(a=-5\), \(b=-\frac3{16}\)}{}{}}
\End

\Data\ID{2000017206}
\Updated{2024-03-03 09:03:25}
\Level{B}
\SubArea{Matrices and determinants}
\LANGen{
\Question{
For what values of \(a\), \(b\) and \(c\) is the matrix
\[
\left (\MATRIX{ 
a-b+c& 0 & a-2b\cr 
a+c-2b & 2^{a-2b} & 0\cr 
2b-a & 0 & c-a+3b
 } \right )
\]
an identity matrix?}
{\(a=2\), \(b=1\), \(c=0\)}{\(a=1\), \(b=2\), \(c=1\)}{\(a=2\), \(b=1\), \(c=3\)}{\(a=1\), \(b=1\), \(c=0\)}{}{}}
\LANGpl{
\Question{
Dla jakich wartości \(a\), \(b\) i \(c\) macierz
\[
\left (\MATRIX{ 
a-b+c& 0 & a-2b\cr 
a+c-2b & 2^{a-2b} & 0\cr 
2b-a & 0 & c-a+3b
 } \right )
\]
to macierz jednostkowa?}
{\(a=2\), \(b=1\), \(c=0\)}{\(a=1\), \(b=2\), \(c=1\)}{\(a=2\), \(b=1\), \(c=3\)}{\(a=1\), \(b=1\), \(c=0\)}{}{}}
\LANGcs{
\Question{
Pro jaké hodnoty \(a\), \(b\) a \(c\) je matice
\[
\left (\MATRIX{ 
a-b+c& 0 & a-2b\cr 
a+c-2b & 2^{a-2b} & 0\cr 
2b-a & 0 & c-a+3b
 } \right )
\]
jednotkovou maticí?}
{\(a=2\), \(b=1\), \(c=0\)}{\(a=1\), \(b=2\), \(c=1\)}{\(a=2\), \(b=1\), \(c=3\)}{\(a=1\), \(b=1\), \(c=0\)}{}{}}
\LANGsk{
\Question{
Pre aké hodnoty \(a\), \(b\) a \(c\) je matica
\[
\left (\MATRIX{ 
a-b+c& 0 & a-2b\cr 
a+c-2b & 2^{a-2b} & 0\cr 
2b-a & 0 & c-a+3b
 } \right )
\]
jednotkovou maticou?}
{\(a=2\), \(b=1\), \(c=0\)}{\(a=1\), \(b=2\), \(c=1\)}{\(a=2\), \(b=1\), \(c=3\)}{\(a=1\), \(b=1\), \(c=0\)}{}{}}
\LANGes{
\Question{
Halla los valores de \(a\), \(b\) y \(c\) para que la matriz sea una matriz identidad.
\[
\left (\MATRIX{ 
a-b+c& 0 & a-2b\cr 
a+c-2b & 2^{a-2b} & 0\cr 
2b-a & 0 & c-a+3b
 } \right )
\]}
{\(a=2\), \(b=1\), \(c=0\)}{\(a=1\), \(b=2\), \(c=1\)}{\(a=2\), \(b=1\), \(c=3\)}{\(a=1\), \(b=1\), \(c=0\)}{}{}}
\End

\Data\ID{2000017208}
\Updated{2024-03-03 09:03:25}
\Level{B}
\SubArea{Matrices and determinants}
\LANGen{
\Question{
Let \(x \in \mathbb{R}\). What is the sum of all the entries on the main diagonal of the following matrix?
\[
\left (\MATRIX{ 
\sin^2 x& 0& 1\cr 
5 & \cos^2 x & 4\cr 
\sin^2 x & 3 & -1
 } \right )
\]}
{\(0\)}{\(1\)}{\(2\)}{\(4\)}{}{}}
\LANGpl{
\Question{
Niech \(x \in \mathbb{R}\). Jaka jest suma wszystkich wartości na głównej przekątnej poniższej macierzy?
\[
\left (\MATRIX{ 
\sin^2 x& 0& 1\cr 
5 & \cos^2 x & 4\cr 
\sin^2 x & 3 & -1
 } \right )
\]}
{\(0\)}{\(1\)}{\(2\)}{\(4\)}{}{}}
\LANGcs{
\Question{
Nechť \(x \in \mathbb{R}\). Jaká je hodnota součtu prvků na hlavní diagonále následující matice?
\[
\left (\MATRIX{ 
\sin^2 x& 0& 1\cr 
5 & \cos^2 x & 4\cr 
\sin^2 x & 3 & -1
 } \right )
\]}
{\(0\)}{\(1\)}{\(2\)}{\(4\)}{}{}}
\LANGsk{
\Question{
Nech \(x \in \mathbb{R}\). Aká je hodnota súčtu prvkov na hlavnej diagonále nasledujúcej matice?
\[
\left (\MATRIX{ 
\sin^2 x& 0& 1\cr 
5 & \cos^2 x & 4\cr 
\sin^2 x & 3 & -1
 } \right )
\]}
{\(0\)}{\(1\)}{\(2\)}{\(4\)}{}{}}
\LANGes{
\Question{
Suponiendo que \(x \in \mathbb{R}\), halla la suma de todas las entradas en la diagonal principal de la siguiente matriz:
\[
\left (\MATRIX{ 
\sin^2 x& 0& 1\cr 
5 & \cos^2 x & 4\cr 
\sin^2 x & 3 & -1
 } \right )
\]}
{\(0\)}{\(1\)}{\(2\)}{\(4\)}{}{}}
\End

\Data\ID{2000017209}
\Updated{2024-03-03 09:03:25}
\Level{B}
\SubArea{Matrices and determinants}
\LANGen{
\Question{
For what real value of \(a\) is the following matrix an upper triangular matrix?
\[
\left (\MATRIX{ 
2^{a-3}& 3+a& a-2\cr 
3-a & (\sqrt3)^{a-1} & 1+a\cr 
\sin (\pi a) & a-\sqrt9 & 2+\sqrt3
 } \right )
\]}
{\(a=3\)}{\(a=2\)}{\(a=-3\)}{\(a=-2\)}{}{}}
\LANGpl{
\Question{
Dla jakiej wartości \(a\) podana macierz to górna macierz trójkątna?
\[
\left (\MATRIX{ 
2^{a-3}& 3+a& a-2\cr 
3-a & (\sqrt3)^{a-1} & 1+a\cr 
\sin (\pi a) & a-\sqrt9 & 2+\sqrt3
 } \right )
\]}
{\(a=3\)}{\(a=2\)}{\(a=-3\)}{\(a=-2\)}{}{}}
\LANGcs{
\Question{
Pro jaké reálné číslo \(a\) je následující matice horní trojúhelníkovou maticí?
\[
\left (\MATRIX{ 
2^{a-3}& 3+a& a-2\cr 
3-a & (\sqrt3)^{a-1} & 1+a\cr 
\sin (\pi a) & a-\sqrt9 & 2+\sqrt3
 } \right )
\]}
{\(a=3\)}{\(a=2\)}{\(a=-3\)}{\(a=-2\)}{}{}}
\LANGsk{
\Question{
Pre aké reálne číslo \(a\) je nasledujúca matica hornou trojuholníkovou maticou?
\[
\left (\MATRIX{ 
2^{a-3}& 3+a& a-2\cr 
3-a & (\sqrt3)^{a-1} & 1+a\cr 
\sin (\pi a) & a-\sqrt9 & 2+\sqrt3
 } \right )
\]}
{\(a=3\)}{\(a=2\)}{\(a=-3\)}{\(a=-2\)}{}{}}
\LANGes{
\Question{
Halla el valor real de \(a\) para que la matriz sea una matriz triangular superior.
\[
\left (\MATRIX{ 
2^{a-3}& 3+a& a-2\cr 
3-a & (\sqrt3)^{a-1} & 1+a\cr 
\sin (\pi a) & a-\sqrt9 & 2+\sqrt3
 } \right )
\]}
{\(a=3\)}{\(a=2\)}{\(a=-3\)}{\(a=-2\)}{}{}}
\End

\Data\ID{2000017210}
\Updated{2024-03-03 09:03:25}
\Level{B}
\SubArea{Matrices and determinants}
\LANGen{
\Question{
In the following matrix, what is the position of the entry with the largest value?
\[
\left (\MATRIX{ 
1+2& \sin \frac{\pi}4& \mathrm{tg}\,\frac{3\pi}4\cr 
3+\cos \frac{\pi}2 & 2^{\sqrt3} & 5+\mathrm{tg} (-\pi)\cr 
\sqrt{23}& 3-4 & \sin 2\pi -2
 } \right )
\]}
{It lays above the main diagonal.}{It lays on the main diagonal.}{It lays under the main diagonal.}{It lays on the counterdiagonal.}{}{}}
\LANGpl{
\Question{
Gdzie znajduje się największa wartość podanej macierzy?
\[
\left (\MATRIX{ 
1+2& \sin \frac{\pi}4& \mathrm{tg}\,\frac{3\pi}4\cr 
3+\cos \frac{\pi}2 & 2^{\sqrt3} & 5+\mathrm{tg} (-\pi)\cr 
\sqrt{23}& 3-4 & \sin 2\pi -2
 } \right )
\]}
{Powyżej głównej przekątnej.}{Na głównej przekątnej.}{Pod główną przekątna.}{Na przeciwprzekątnej.}{}{}}
\LANGcs{
\Question{
Kde je v následující matici prvek s největší hodnotou? 
\[
\left (\MATRIX{ 
1+2& \sin \frac{\pi}4& \mathrm{tg}\,\frac{3\pi}4\cr 
3+\cos \frac{\pi}2 & 2^{\sqrt3} & 5+\mathrm{tg} (-\pi)\cr 
\sqrt{23}& 3-4 & \sin 2\pi -2
 } \right )
\]}
{Prvek leží nad hlavní diagonálou.}{Prvek leží na hlavní diagonále.}{Prvek leží pod hlavní diagonálou.}{Prvek leží na vedlejší diagonále.}{}{}}
\LANGsk{
\Question{
Na akej pozícii je v nasledujúcej matici prvok s najväčšou hodnotou? 
\[
\left (\MATRIX{ 
1+2& \sin \frac{\pi}4& \mathrm{tg}\,\frac{3\pi}4\cr 
3+\cos \frac{\pi}2 & 2^{\sqrt3} & 5+\mathrm{tg} (-\pi)\cr 
\sqrt{23}& 3-4 & \sin 2\pi -2
 } \right )
\]}
{Prvok leží nad hlavnou diagonálou.}{Prvok leží na hlavnej diagonále.}{Prvok leží pod hlavnou diagonálou.}{Prvok leží na vedľajšej diagonále.}{}{}}
\LANGes{
\Question{
Dada la matriz:
\[
\left (\MATRIX{ 
1+2& \sin \frac{\pi}4& \mathrm{tg}\,\frac{3\pi}4\cr 
3+\cos \frac{\pi}2 & 2^{\sqrt3} & 5+\mathrm{tg} (-\pi)\cr 
\sqrt{23}& 3-4 & \sin 2\pi -2
 } \right )
\]
Determina la posición  de la entrada con el mayor valor.}
{Se encuentra pro encima de la diagonal principal.}{Se encuentra en la diagonal principal.}{Se encuentra debajo de da diagonal principal.}{Se encuentra en la contradiagonal.}{}{}}
\End

\Data\ID{2000017410}
\Updated{2024-03-03 09:03:25}
\Level{B}
\SubArea{Matrices and determinants}
\LANGen{
\Question{
Find the matrix \(M\) for which the following equality is true:
\[ 
3 \cdot \left (\MATRIX{ 
4& -1\cr 
 2 & 5 \cr 
} \right ) - 2\cdot M =
\left (\MATRIX{ 
14& -7 \cr 
4 & 7 \cr 
} \right )
\]}
{\(
\left (\MATRIX{ 
-1& 2\cr 
1 & 4 \cr 
} \right )
\)}{\(
\left (\MATRIX{ 
1& 2\cr 
1 & 4 \cr 
} \right )
\)}{\(
\left (\MATRIX{ 
-1& 2\cr 
-1 & 4 \cr 
} \right )
\)}{\(
\left (\MATRIX{ 
-1& -2\cr 
1 & 4 \cr 
} \right )
\)}{}{}}
\LANGpl{
\Question{
Znajdż macierz \(M\), dla której podana równość jest prawdziwa:
\[ 
3 \cdot \left (\MATRIX{ 
4& -1\cr 
 2 & 5 \cr 
} \right ) - 2\cdot M =
\left (\MATRIX{ 
14& -7 \cr 
4 & 7 \cr 
} \right )
\]}
{\(
\left (\MATRIX{ 
-1& 2\cr 
1 & 4 \cr 
} \right )
\)}{\(
\left (\MATRIX{ 
1& 2\cr 
1 & 4 \cr 
} \right )
\)}{\(
\left (\MATRIX{ 
-1& 2\cr 
-1 & 4 \cr 
} \right )
\)}{\(
\left (\MATRIX{ 
-1& -2\cr 
1 & 4 \cr 
} \right )
\)}{}{}}
\LANGcs{
\Question{
Určete, pro jakou matici \(M\) platí následující rovnost:
\[ 
3 \cdot \left (\MATRIX{ 
4& -1\cr 
 2 & 5 \cr 
} \right ) - 2\cdot M =
\left (\MATRIX{ 
14& -7 \cr 
4 & 7 \cr 
} \right )
\]}
{\(
\left (\MATRIX{ 
-1& 2\cr 
1 & 4 \cr 
} \right )
\)}{\(
\left (\MATRIX{ 
1& 2\cr 
1 & 4 \cr 
} \right )
\)}{\(
\left (\MATRIX{ 
-1& 2\cr 
-1 & 4 \cr 
} \right )
\)}{\(
\left (\MATRIX{ 
-1& -2\cr 
1 & 4 \cr 
} \right )
\)}{}{}}
\LANGsk{
\Question{
Nájdi maticu \(M\) pre ktorú platí rovnosť:
\[ 
3 \cdot \left (\MATRIX{ 
4& -1\cr 
 2 & 5 \cr 
} \right ) - 2\cdot M =
\left (\MATRIX{ 
14& -7 \cr 
4 & 7 \cr 
} \right )
\]}
{\(
\left (\MATRIX{ 
-1& 2\cr 
1 & 4 \cr 
} \right )
\)}{\(
\left (\MATRIX{ 
1& 2\cr 
1 & 4 \cr 
} \right )
\)}{\(
\left (\MATRIX{ 
-1& 2\cr 
-1 & 4 \cr 
} \right )
\)}{\(
\left (\MATRIX{ 
-1& -2\cr 
1 & 4 \cr 
} \right )
\)}{}{}}
\LANGes{
\Question{
Determina la matriz \(M\) para cual es verdadera la siguiente ecuación:
\[ 
3 \cdot \left (\MATRIX{ 
4& -1\cr 
 2 & 5 \cr 
} \right ) - 2\cdot M =
\left (\MATRIX{ 
14& -7 \cr 
4 & 7 \cr 
} \right )
\]}
{\(
\left (\MATRIX{ 
-1& 2\cr 
1 & 4 \cr 
} \right )
\)}{\(
\left (\MATRIX{ 
1& 2\cr 
1 & 4 \cr 
} \right )
\)}{\(
\left (\MATRIX{ 
-1& 2\cr 
-1 & 4 \cr 
} \right )
\)}{\(
\left (\MATRIX{ 
-1& -2\cr 
1 & 4 \cr 
} \right )
\)}{}{}}
\End

\Data\ID{2000017411}
\Updated{2024-03-03 09:03:25}
\Level{B}
\SubArea{Matrices and determinants}
\LANGen{
\Question{
For which matrix \(M\) the following equality holds?
\[ 
2 \cdot \left (\MATRIX{ 
5& -4\cr 
 7 & -3 \cr 
} \right ) - 4\cdot M =
\left (\MATRIX{ 
2& -20 \cr 
18 & -22 \cr 
} \right )
\]}
{\(
\left (\MATRIX{ 
2& 3\cr 
-1 & 4 \cr 
} \right )
\)}{\(
\left (\MATRIX{ 
2& 3\cr 
1 & 4 \cr 
} \right )
\)}{\(
\left (\MATRIX{ 
2& -3\cr 
-1 & 4 \cr 
} \right )
\)}{\(
\left (\MATRIX{ 
-2& 3\cr 
-1 & 4 \cr 
} \right )
\)}{}{}}
\LANGpl{
\Question{
Znajdź macierz  \(M\), dla której zachodzi następująca równość.
\[ 
2 \cdot \left (\MATRIX{ 
5& -4\cr 
 7 & -3 \cr 
} \right ) - 4\cdot M =
\left (\MATRIX{ 
2& -20 \cr 
18 & -22 \cr 
} \right )
\]}
{\(
\left (\MATRIX{ 
2& 3\cr 
-1 & 4 \cr 
} \right )
\)}{\(
\left (\MATRIX{ 
2& 3\cr 
1 & 4 \cr 
} \right )
\)}{\(
\left (\MATRIX{ 
2& -3\cr 
-1 & 4 \cr 
} \right )
\)}{\(
\left (\MATRIX{ 
-2& 3\cr 
-1 & 4 \cr 
} \right )
\)}{}{}}
\LANGcs{
\Question{
Pro kterou matici \(M\) platí následující rovnost?
\[ 
2 \cdot \left (\MATRIX{ 
5& -4\cr 
 7 & -3 \cr 
} \right ) - 4\cdot M =
\left (\MATRIX{ 
2& -20 \cr 
18 & -22 \cr 
} \right )
\]}
{\(
\left (\MATRIX{ 
2& 3\cr 
-1 & 4 \cr 
} \right )
\)}{\(
\left (\MATRIX{ 
2& 3\cr 
1 & 4 \cr 
} \right )
\)}{\(
\left (\MATRIX{ 
2& -3\cr 
-1 & 4 \cr 
} \right )
\)}{\(
\left (\MATRIX{ 
-2& 3\cr 
-1 & 4 \cr 
} \right )
\)}{}{}}
\LANGsk{
\Question{
Pre ktorú maticu \(M\) platí rovnosť?
\[ 
2 \cdot \left (\MATRIX{ 
5& -4\cr 
 7 & -3 \cr 
} \right ) - 4\cdot M =
\left (\MATRIX{ 
2& -20 \cr 
18 & -22 \cr 
} \right )
\]}
{\(
\left (\MATRIX{ 
2& 3\cr 
-1 & 4 \cr 
} \right )
\)}{\(
\left (\MATRIX{ 
2& 3\cr 
1 & 4 \cr 
} \right )
\)}{\(
\left (\MATRIX{ 
2& -3\cr 
-1 & 4 \cr 
} \right )
\)}{\(
\left (\MATRIX{ 
-2& 3\cr 
-1 & 4 \cr 
} \right )
\)}{}{}}
\LANGes{
\Question{
Determina la matriz \(M\) para cual es verdadera la siguiente ecuación:
\[ 
2 \cdot \left (\MATRIX{ 
5& -4\cr 
 7 & -3 \cr 
} \right ) - 4\cdot M =
\left (\MATRIX{ 
2& -20 \cr 
18 & -22 \cr 
} \right )
\]}
{\(
\left (\MATRIX{ 
2& 3\cr 
-1 & 4 \cr 
} \right )
\)}{\(
\left (\MATRIX{ 
2& 3\cr 
1 & 4 \cr 
} \right )
\)}{\(
\left (\MATRIX{ 
2& -3\cr 
-1 & 4 \cr 
} \right )
\)}{\(
\left (\MATRIX{ 
-2& 3\cr 
-1 & 4 \cr 
} \right )
\)}{}{}}
\End

\Data\ID{2000017412}
\Updated{2024-03-03 09:03:25}
\Level{B}
\SubArea{Matrices and determinants}
\LANGen{
\Question{
For which matrices \(X\) and \(Y\) are both following equations true?
\[ 
3\cdot X - Y= 
\left (\MATRIX{ 
2& -2 & -5\cr 
 8 & 3 & 2\cr 
} \right ) 
\]
\[ 
X + 2 \cdot Y= 
\left (\MATRIX{ 
3& 4& -4\cr 
 5 & 1 & 3\cr 
} \right ) 
\]}
{\(X= 
\left (\MATRIX{ 
1& 0& -2\cr 
 3 & 1 & 1\cr 
} \right ) 
\),
\(Y= 
\left (\MATRIX{ 
1& 2& -1\cr 
 1 & 0 & 1\cr 
} \right ) 
\)}{\(X= 
\left (\MATRIX{ 
1& 0& 2\cr 
 3 & 1 & 1\cr 
} \right ) 
\),
\(Y= 
\left (\MATRIX{ 
1& 2& -1\cr 
 1 & 0 & 1\cr 
} \right ) 
\)}{\(X= 
\left (\MATRIX{ 
1& 0& -2\cr 
 3 & 1 & 1\cr 
} \right ) 
\),
\(Y= 
\left (\MATRIX{ 
1& 2& 1\cr 
 1 & 0 & 1\cr 
} \right ) 
\)}{\(X= 
\left (\MATRIX{ 
1& 0& -2\cr 
 3 & -1 & 1\cr 
} \right ) 
\),
\(Y= 
\left (\MATRIX{ 
1& 2& -1\cr 
 -1 & 0 & 1\cr 
} \right ) 
\)}{}{}}
\LANGpl{
\Question{
Dla których macierzy \(X\) i \(Y\) oba poniższe równania są prawdziwe?
\[ 
3\cdot X - Y= 
\left (\MATRIX{ 
2& -2 & -5\cr 
 8 & 3 & 2\cr 
} \right ) 
\]
\[ 
X + 2 \cdot Y= 
\left (\MATRIX{ 
3& 4& -4\cr 
 5 & 1 & 3\cr 
} \right ) 
\]}
{\(X= 
\left (\MATRIX{ 
1& 0& -2\cr 
 3 & 1 & 1\cr 
} \right ) 
\),
\(Y= 
\left (\MATRIX{ 
1& 2& -1\cr 
 1 & 0 & 1\cr 
} \right ) 
\)}{\(X= 
\left (\MATRIX{ 
1& 0& 2\cr 
 3 & 1 & 1\cr 
} \right ) 
\),
\(Y= 
\left (\MATRIX{ 
1& 2& -1\cr 
 1 & 0 & 1\cr 
} \right ) 
\)}{\(X= 
\left (\MATRIX{ 
1& 0& -2\cr 
 3 & 1 & 1\cr 
} \right ) 
\),
\(Y= 
\left (\MATRIX{ 
1& 2& 1\cr 
 1 & 0 & 1\cr 
} \right ) 
\)}{\(X= 
\left (\MATRIX{ 
1& 0& -2\cr 
 3 & -1 & 1\cr 
} \right ) 
\),
\(Y= 
\left (\MATRIX{ 
1& 2& -1\cr 
 -1 & 0 & 1\cr 
} \right ) 
\)}{}{}}
\LANGcs{
\Question{
Pro které matice \(X\) a \(Y\) platí současně obě následující rovnice?
\[ 
3\cdot X - Y= 
\left (\MATRIX{ 
2& -2 & -5\cr 
 8 & 3 & 2\cr 
} \right ) 
\]
\[ 
X + 2 \cdot Y= 
\left (\MATRIX{ 
3& 4& -4\cr 
 5 & 1 & 3\cr 
} \right ) 
\]}
{\(X= 
\left (\MATRIX{ 
1& 0& -2\cr 
 3 & 1 & 1\cr 
} \right ) 
\),
\(Y= 
\left (\MATRIX{ 
1& 2& -1\cr 
 1 & 0 & 1\cr 
} \right ) 
\)}{\(X= 
\left (\MATRIX{ 
1& 0& 2\cr 
 3 & 1 & 1\cr 
} \right ) 
\),
\(Y= 
\left (\MATRIX{ 
1& 2& -1\cr 
 1 & 0 & 1\cr 
} \right ) 
\)}{\(X= 
\left (\MATRIX{ 
1& 0& -2\cr 
 3 & 1 & 1\cr 
} \right ) 
\),
\(Y= 
\left (\MATRIX{ 
1& 2& 1\cr 
 1 & 0 & 1\cr 
} \right ) 
\)}{\(X= 
\left (\MATRIX{ 
1& 0& -2\cr 
 3 & -1 & 1\cr 
} \right ) 
\),
\(Y= 
\left (\MATRIX{ 
1& 2& -1\cr 
 -1 & 0 & 1\cr 
} \right ) 
\)}{}{}}
\LANGsk{
\Question{
Pre ktoré matice \(X\) a \(Y\) platí súčasne rovnosť pre obe matice?
\[ 
3\cdot X - Y= 
\left (\MATRIX{ 
2& -2 & -5\cr 
 8 & 3 & 2\cr 
} \right ) 
\]
\[ 
X + 2 \cdot Y= 
\left (\MATRIX{ 
3& 4& -4\cr 
 5 & 1 & 3\cr 
} \right ) 
\]}
{\(X= 
\left (\MATRIX{ 
1& 0& -2\cr 
 3 & 1 & 1\cr 
} \right ) 
\),
\(Y= 
\left (\MATRIX{ 
1& 2& -1\cr 
 1 & 0 & 1\cr 
} \right ) 
\)}{\(X= 
\left (\MATRIX{ 
1& 0& 2\cr 
 3 & 1 & 1\cr 
} \right ) 
\),
\(Y= 
\left (\MATRIX{ 
1& 2& -1\cr 
 1 & 0 & 1\cr 
} \right ) 
\)}{\(X= 
\left (\MATRIX{ 
1& 0& -2\cr 
 3 & 1 & 1\cr 
} \right ) 
\),
\(Y= 
\left (\MATRIX{ 
1& 2& 1\cr 
 1 & 0 & 1\cr 
} \right ) 
\)}{\(X= 
\left (\MATRIX{ 
1& 0& -2\cr 
 3 & -1 & 1\cr 
} \right ) 
\),
\(Y= 
\left (\MATRIX{ 
1& 2& -1\cr 
 -1 & 0 & 1\cr 
} \right ) 
\)}{}{}}
\LANGes{
\Question{
Determina las matrices \(X\) y \(Y\) para cuales es verdadera la siguiente ecuación:
\[ 
3\cdot X - Y= 
\left (\MATRIX{ 
2& -2 & -5\cr 
 8 & 3 & 2\cr 
} \right ) 
\]
\[ 
X + 2 \cdot Y= 
\left (\MATRIX{ 
3& 4& -4\cr 
 5 & 1 & 3\cr 
} \right ) 
\]}
{\(X= 
\left (\MATRIX{ 
1& 0& -2\cr 
 3 & 1 & 1\cr 
} \right ) 
\),
\(Y= 
\left (\MATRIX{ 
1& 2& -1\cr 
 1 & 0 & 1\cr 
} \right ) 
\)}{\(X= 
\left (\MATRIX{ 
1& 0& 2\cr 
 3 & 1 & 1\cr 
} \right ) 
\),
\(Y= 
\left (\MATRIX{ 
1& 2& -1\cr 
 1 & 0 & 1\cr 
} \right ) 
\)}{\(X= 
\left (\MATRIX{ 
1& 0& -2\cr 
 3 & 1 & 1\cr 
} \right ) 
\),
\(Y= 
\left (\MATRIX{ 
1& 2& 1\cr 
 1 & 0 & 1\cr 
} \right ) 
\)}{\(X= 
\left (\MATRIX{ 
1& 0& -2\cr 
 3 & -1 & 1\cr 
} \right ) 
\),
\(Y= 
\left (\MATRIX{ 
1& 2& -1\cr 
 -1 & 0 & 1\cr 
} \right ) 
\)}{}{}}
\End

\Data\ID{2000017413}
\Updated{2024-03-03 09:03:25}
\Level{B}
\SubArea{Matrices and determinants}
\LANGen{
\Question{
Consider the matrices
\[ 
A=\left (\MATRIX{ 
1& 0 & 0\cr 
 -1 & 1 & 1 \cr 
2& 1 & 0} \right ),~
B=\left (\MATRIX{ 
-1& 1 & 0\cr 
0 & 1 & 1 \cr 
1 &1& -1 } \right ).
\]
For which matrix \(X\) holds \(A \cdot B - X= B^2\)?}
{\(
\left (\MATRIX{ 
-2& 1 & -1\cr 
 1 & -1& 0\cr 
0 & 2& -1 } \right )
\)}{\(
\left (\MATRIX{ 
-2& 1 & -1\cr 
 1 & 1& 0\cr 
0 & 2& -1 } \right )
\)}{\(
\left (\MATRIX{ 
-2& 1 & 1\cr 
 1 & -1& 0\cr 
0 & 2& -1 } \right )
\)}{\(
\left (\MATRIX{ 
2& 1 & -1\cr 
 1 & -1& 0\cr 
0 & 2& -1 } \right )
\)}{}{}}
\LANGpl{
\Question{
Dane są macierze
\[ 
A=\left (\MATRIX{ 
1& 0 & 0\cr 
 -1 & 1 & 1 \cr 
2& 1 & 0} \right ),~
B=\left (\MATRIX{ 
-1& 1 & 0\cr 
0 & 1 & 1 \cr 
1 &1& -1 } \right ).
\]
Wyznacz macierz \(X\), dla której spełnione jest równanie \(A \cdot B - X= B^2\).}
{\(
\left (\MATRIX{ 
-2& 1 & -1\cr 
 1 & -1& 0\cr 
0 & 2& -1 } \right )
\)}{\(
\left (\MATRIX{ 
-2& 1 & -1\cr 
 1 & 1& 0\cr 
0 & 2& -1 } \right )
\)}{\(
\left (\MATRIX{ 
-2& 1 & 1\cr 
 1 & -1& 0\cr 
0 & 2& -1 } \right )
\)}{\(
\left (\MATRIX{ 
2& 1 & -1\cr 
 1 & -1& 0\cr 
0 & 2& -1 } \right )
\)}{}{}}
\LANGcs{
\Question{
Jsou dány matice
\[ 
A=\left (\MATRIX{ 
1& 0 & 0\cr 
 -1 & 1 & 1 \cr 
2& 1 & 0} \right ),~
B=\left (\MATRIX{ 
-1& 1 & 0\cr 
0 & 1 & 1 \cr 
1 &1& -1 } \right ).
\]
Pro kterou matici \(X\)  platí \(A \cdot B - X= B^2\)?}
{\(
\left (\MATRIX{ 
-2& 1 & -1\cr 
 1 & -1& 0\cr 
0 & 2& -1 } \right )
\)}{\(
\left (\MATRIX{ 
-2& 1 & -1\cr 
 1 & 1& 0\cr 
0 & 2& -1 } \right )
\)}{\(
\left (\MATRIX{ 
-2& 1 & 1\cr 
 1 & -1& 0\cr 
0 & 2& -1 } \right )
\)}{\(
\left (\MATRIX{ 
2& 1 & -1\cr 
 1 & -1& 0\cr 
0 & 2& -1 } \right )
\)}{}{}}
\LANGsk{
\Question{
Sú dané matice
\[ 
A=\left (\MATRIX{ 
1& 0 & 0\cr 
 -1 & 1 & 1 \cr 
2& 1 & 0} \right ),~
B=\left (\MATRIX{ 
-1& 1 & 0\cr 
0 & 1 & 1 \cr 
1 &1& -1 } \right ).
\]
Pre ktorú maticu \(X\) platí \(A \cdot B - X= B^2\)?}
{\(
\left (\MATRIX{ 
-2& 1 & -1\cr 
 1 & -1& 0\cr 
0 & 2& -1 } \right )
\)}{\(
\left (\MATRIX{ 
-2& 1 & -1\cr 
 1 & 1& 0\cr 
0 & 2& -1 } \right )
\)}{\(
\left (\MATRIX{ 
-2& 1 & 1\cr 
 1 & -1& 0\cr 
0 & 2& -1 } \right )
\)}{\(
\left (\MATRIX{ 
2& 1 & -1\cr 
 1 & -1& 0\cr 
0 & 2& -1 } \right )
\)}{}{}}
\LANGes{
\Question{
Dadas las matrices:
\[ 
A=\left (\MATRIX{ 
1& 0 & 0\cr 
 -1 & 1 & 1 \cr 
2& 1 & 0} \right ),~
B=\left (\MATRIX{ 
-1& 1 & 0\cr 
0 & 1 & 1 \cr 
1 &1& -1 } \right).
\]
¿Para qué matriz \(X\) vale \(A \cdot B - X= B^2\)?}
{\(
\left (\MATRIX{ 
-2& 1 & -1\cr 
 1 & -1& 0\cr 
0 & 2& -1 } \right )
\)}{\(
\left (\MATRIX{ 
-2& 1 & -1\cr 
 1 & 1& 0\cr 
0 & 2& -1 } \right )
\)}{\(
\left (\MATRIX{ 
-2& 1 & 1\cr 
 1 & -1& 0\cr 
0 & 2& -1 } \right )
\)}{\(
\left (\MATRIX{ 
2& 1 & -1\cr 
 1 & -1& 0\cr 
0 & 2& -1 } \right )
\)}{}{}}
\End

\Data\ID{2000018304}
\Updated{2024-03-03 09:03:23}
\Level{B}
\SubArea{Matrices and determinants}
\LANGen{
\Question{
Which matrices X and Y make valid both equalities given below?
\[ 
2X+Y
=
\left (\MATRIX{ 
1 &4\cr 
2 & 0\cr 
} \right ) 
\]
\[ 
X-Y
=
\left (\MATRIX{ 
1 &-1\cr 
1 & 0\cr 
} \right ) 
\]}
{\(
X
=
\left (\MATRIX{ 
\frac23 &1\cr 
1 & 0\cr 
} \right ) 
\) and
\(
Y
=
\left (\MATRIX{ 
-\frac13 &2\cr 
0& 0\cr 
} \right ) 
\)}{\(
X
=
\left (\MATRIX{ 
\frac23 &1\cr 
1 & 0\cr 
} \right ) 
\) and
\(
Y
=
\left (\MATRIX{ 
-\frac13 &4\cr 
0& 0\cr 
} \right ) 
\)}{\(
X
=
\left (\MATRIX{ 
\frac23 &1\cr 
1 & 0\cr 
} \right ) 
\) and
\(
Y
=
\left (\MATRIX{ 
\frac13 &2\cr 
0& 0\cr 
} \right ) 
\)}{\(
X
=
\left (\MATRIX{ 
\frac23 &1\cr 
1 & 1\cr 
} \right ) 
\) and
\(
Y
=
\left (\MATRIX{ 
-\frac13 &4\cr 
0& 0\cr 
} \right ) 
\)}{}{}}
\LANGpl{
\Question{
Które macierze X i Y spełniają poniższe równania?
\[ 
2X+Y
=
\left (\MATRIX{ 
1 &4\cr 
2 & 0\cr 
} \right ) 
\]
\[ 
X-Y
=
\left (\MATRIX{ 
1 &-1\cr 
1 & 0\cr 
} \right ) 
\]}
{\(
X
=
\left (\MATRIX{ 
\frac23 &1\cr 
1 & 0\cr 
} \right ) 
\) i
\(
Y
=
\left (\MATRIX{ 
-\frac13 &2\cr 
0& 0\cr 
} \right ) 
\)}{\(
X
=
\left (\MATRIX{ 
\frac23 &1\cr 
1 & 0\cr 
} \right ) 
\) i
\(
Y
=
\left (\MATRIX{ 
-\frac13 &4\cr 
0& 0\cr 
} \right ) 
\)}{\(
X
=
\left (\MATRIX{ 
\frac23 &1\cr 
1 & 0\cr 
} \right ) 
\) i
\(
Y
=
\left (\MATRIX{ 
\frac13 &2\cr 
0& 0\cr 
} \right ) 
\)}{\(
X
=
\left (\MATRIX{ 
\frac23 &1\cr 
1 & 1\cr 
} \right ) 
\) i
\(
Y
=
\left (\MATRIX{ 
-\frac13 &4\cr 
0& 0\cr 
} \right ) 
\)}{}{}}
\LANGcs{
\Question{
Které matice X a Y vyhovují níže uvedeným rovnostem?
\[ 
2X+Y
=
\left (\MATRIX{ 
1 &4\cr 
2 & 0\cr 
} \right ) 
\]
\[ 
X-Y
=
\left (\MATRIX{ 
1 &-1\cr 
1 & 0\cr 
} \right ) 
\]}
{\(
X
=
\left (\MATRIX{ 
\frac23 &1\cr 
1 & 0\cr 
} \right ) 
\) a
\(
Y
=
\left (\MATRIX{ 
-\frac13 &2\cr 
0& 0\cr 
} \right ) 
\)}{\(
X
=
\left (\MATRIX{ 
\frac23 &1\cr 
1 & 0\cr 
} \right ) 
\) a
\(
Y
=
\left (\MATRIX{ 
-\frac13 &4\cr 
0& 0\cr 
} \right ) 
\)}{\(
X
=
\left (\MATRIX{ 
\frac23 &1\cr 
1 & 0\cr 
} \right ) 
\) a
\(
Y
=
\left (\MATRIX{ 
\frac13 &2\cr 
0& 0\cr 
} \right ) 
\)}{\(
X
=
\left (\MATRIX{ 
\frac23 &1\cr 
1 & 1\cr 
} \right ) 
\) a
\(
Y
=
\left (\MATRIX{ 
-\frac13 &4\cr 
0& 0\cr 
} \right ) 
\)}{}{}}
\LANGsk{
\Question{
Ktoré matice X a Y vyhovujú nižšie uvedeným rovnostiam?
\[ 
2X+Y
=
\left (\MATRIX{ 
1 &4\cr 
2 & 0\cr 
} \right ) 
\]
\[ 
X-Y
=
\left (\MATRIX{ 
1 &-1\cr 
1 & 0\cr 
} \right ) 
\]}
{\(
X
=
\left (\MATRIX{ 
\frac23 &1\cr 
1 & 0\cr 
} \right ) 
\) a
\(
Y
=
\left (\MATRIX{ 
-\frac13 &2\cr 
0& 0\cr 
} \right ) 
\)}{\(
X
=
\left (\MATRIX{ 
\frac23 &1\cr 
1 & 0\cr 
} \right ) 
\) a
\(
Y
=
\left (\MATRIX{ 
-\frac13 &4\cr 
0& 0\cr 
} \right ) 
\)}{\(
X
=
\left (\MATRIX{ 
\frac23 &1\cr 
1 & 0\cr 
} \right ) 
\) a
\(
Y
=
\left (\MATRIX{ 
\frac13 &2\cr 
0& 0\cr 
} \right ) 
\)}{\(
X
=
\left (\MATRIX{ 
\frac23 &1\cr 
1 & 1\cr 
} \right ) 
\) a
\(
Y
=
\left (\MATRIX{ 
-\frac13 &4\cr 
0& 0\cr 
} \right ) 
\)}{}{}}
\LANGes{
\Question{
¿Qué matrices X y Y son válidas para las siguientes ecuaciones?
\[ 
2X+Y
=
\left (\MATRIX{ 
1 &4\cr 
2 & 0\cr 
} \right ) 
\]
\[ 
X-Y
=
\left (\MATRIX{ 
1 &-1\cr 
1 & 0\cr 
} \right ) 
\]}
{\(
X
=
\left (\MATRIX{ 
\frac23 &1\cr 
1 & 0\cr 
} \right ) 
\) y
\(
Y
=
\left (\MATRIX{ 
-\frac13 &2\cr 
0& 0\cr 
} \right ) 
\)}{\(
X
=
\left (\MATRIX{ 
\frac23 &1\cr 
1 & 0\cr 
} \right ) 
\) y
\(
Y
=
\left (\MATRIX{ 
-\frac13 &4\cr 
0& 0\cr 
} \right ) 
\)}{\(
X
=
\left (\MATRIX{ 
\frac23 &1\cr 
1 & 0\cr 
} \right ) 
\) y
\(
Y
=
\left (\MATRIX{ 
\frac13 &2\cr 
0& 0\cr 
} \right ) 
\)}{\(
X
=
\left (\MATRIX{ 
\frac23 &1\cr 
1 & 1\cr 
} \right ) 
\) y
\(
Y
=
\left (\MATRIX{ 
-\frac13 &4\cr 
0& 0\cr 
} \right ) 
\)}{}{}}
\End

\Data\ID{2000018901}
\Updated{2024-03-03 09:03:23}
\Level{B}
\SubArea{Matrices and determinants}
\LANGen{
\Question{
Determine the rank of the following matrix.
\[ 
\left (\MATRIX{ 
3& 8 \cr 
 1 & 5
\cr 
5 & 18  } \right )
\]}
{\(2\)}{\(1\)}{\(3\)}{\(4\)}{}{}}
\LANGpl{
\Question{
Oblicz rangę następującej macierzy.
\[ 
\left (\MATRIX{ 
3& 8 \cr 
 1 & 5
\cr 
5 & 18  } \right )
\]}
{\(2\)}{\(1\)}{\(3\)}{\(4\)}{}{}}
\LANGcs{
\Question{
Vypočtěte hodnost následující matice.
\[ 
\left (\MATRIX{ 
3& 8 \cr 
 1 & 5
\cr 
5 & 18  } \right )
\]}
{\(2\)}{\(1\)}{\(3\)}{\(4\)}{}{}}
\LANGsk{
\Question{
Vypočítajte hodnosť nasledujúcej matice.
\[ 
\left (\MATRIX{ 
3& 8 \cr 
 1 & 5
\cr 
5 & 18  } \right )
\]}
{\(2\)}{\(1\)}{\(3\)}{\(4\)}{}{}}
\LANGes{
\Question{
Halla el rango de la siguiente matriz.
\[ 
\left (\MATRIX{ 
3& 8 \cr 
 1 & 5
\cr 
5 & 18  } \right )
\]}
{\(2\)}{\(1\)}{\(3\)}{\(4\)}{}{}}
\End

\Data\ID{2000018902}
\Updated{2024-03-03 09:03:23}
\Level{B}
\SubArea{Matrices and determinants}
\LANGen{
\Question{
Determine the value of \(p\) so that the rank of the following matrix differs from \(3\).
\[ 
\left (\MATRIX{ 
2& 5 & 1\cr 
 4 & 10 & p-1
\cr 
1 & 3 & 8} \right )
\]}
{\(3\)}{\(0\)}{\(1\)}{\(2\)}{}{}}
\LANGpl{
\Question{
Oblicz wartość \(p\) tak, aby ranga poniższej macierzy różniła się od \(3\).
\[ 
\left (\MATRIX{ 
2& 5 & 1\cr 
 4 & 10 & p-1
\cr 
1 & 3 & 8} \right )
\]}
{\(3\)}{\(0\)}{\(1\)}{\(2\)}{}{}}
\LANGcs{
\Question{
Vypočtěte hodnotu \(p\),  pro kterou nemá následující matice hodnost  \(3\):
\[ 
\left (\MATRIX{ 
2& 5 & 1\cr 
 4 & 10 & p-1
\cr 
1 & 3 & 8} \right )
\]}
{\(3\)}{\(0\)}{\(1\)}{\(2\)}{}{}}
\LANGsk{
\Question{
Vypočítajte hodnotu \(p\), pre ktorú nemá nasledujúca matica hodnosť \(3\).
\[ 
\left (\MATRIX{ 
2& 5 & 1\cr 
 4 & 10 & p-1
\cr 
1 & 3 & 8} \right )
\]}
{\(3\)}{\(0\)}{\(1\)}{\(2\)}{}{}}
\LANGes{
\Question{
Halla el valor de \(p\) de modo que el rango de la siguiente matriz difiera de \(3\).
\[ 
\left (\MATRIX{ 
2& 5 & 1\cr 
 4 & 10 & p-1
\cr 
1 & 3 & 8} \right )
\]}
{\(3\)}{\(0\)}{\(1\)}{\(2\)}{}{}}
\End

\Data\ID{2000018903}
\Updated{2024-03-03 09:03:23}
\Level{B}
\SubArea{Matrices and determinants}
\LANGen{
\Question{
Determine values of \(k\) and \(l\) so that the following matrix has rank \(1\).
\[ 
\left (\MATRIX{ 
4& 16 & k-25\cr 
 -7 & l-30 & 35
\cr 
3 & 12 & -15} \right )
\]}
{\(k=5\), \(l=2\)}{\(k=2\), \(l=5\)}{\(k=-5\), \(l=2\)}{\(k=-2\), \(l=-5\)}{}{}}
\LANGpl{
\Question{
Podaj wartości dla \(k\) oraz \(l\)  tak, aby poniższa macierz miała rangę \(1\).
\[ 
\left (\MATRIX{ 
4& 16 & k-25\cr 
 -7 & l-30 & 35
\cr 
3 & 12 & -15} \right )
\]}
{\(k=5\), \(l=2\)}{\(k=2\), \(l=5\)}{\(k=-5\), \(l=2\)}{\(k=-2\), \(l=-5\)}{}{}}
\LANGcs{
\Question{
Určete hodnoty  \(k\) a \(l\) tak, aby následující matice měla hodnost \(1\).
\[ 
\left (\MATRIX{ 
4& 16 & k-25\cr 
 -7 & l-30 & 35
\cr 
3 & 12 & -15} \right )
\]}
{\(k=5\), \(l=2\)}{\(k=2\), \(l=5\)}{\(k=-5\), \(l=2\)}{\(k=-2\), \(l=-5\)}{}{}}
\LANGsk{
\Question{
Určte hodnoty  \(k\) a \(l\) tak, aby nasledujúca matica mala hodnosť \(1\).
\[ 
\left (\MATRIX{ 
4& 16 & k-25\cr 
 -7 & l-30 & 35
\cr 
3 & 12 & -15} \right )
\]}
{\(k=5\), \(l=2\)}{\(k=2\), \(l=5\)}{\(k=-5\), \(l=2\)}{\(k=-2\), \(l=-5\)}{}{}}
\LANGes{
\Question{
Halla los valores de \(k\) y \(l\) para que la siguiente matriz tenga el rango \(1\).
\[ 
\left (\MATRIX{ 
4& 16 & k-25\cr 
 -7 & l-30 & 35
\cr 
3 & 12 & -15} \right )
\]}
{\(k=5\), \(l=2\)}{\(k=2\), \(l=5\)}{\(k=-5\), \(l=2\)}{\(k=-2\), \(l=-5\)}{}{}}
\End

\Data\ID{2000018904}
\Updated{2024-03-03 09:03:23}
\Level{B}
\SubArea{Matrices and determinants}
\LANGen{
\Question{
Specify \(p\) so that the matrix \(A\) has rank \(2\).
\[ 
A=\left(\MATRIX{ 
1& p+1\cr 
 3 & 6+2p \cr 
-1 & -8} \right)
\]}
{\(A\) has rank \(2\) for every real number \(p\).}{\(p=7\)}{\(p=9\)}{\(A\) has not rank \(2\) for any \(p\).}{}{}}
\LANGpl{
\Question{
Określ \(p\), tak aby macierz \(A\) miała rangę \(2\).
\[ 
A=\left(\MATRIX{ 
1& p+1\cr 
 3 & 6+2p \cr 
-1 & -8} \right)
\]}
{\(A\) ma rangę \(2\) dla każdej liczby rzeczywistej \(p\).}{\(p=7\)}{\(p=9\)}{\(A\) nie ma rangi  \(2\) dla \(p\).}{}{}}
\LANGcs{
\Question{
Najděte \(p\) tak, aby matice \(A\) měla hodnost \(2\).
\[ 
A=\left(\MATRIX{ 
1& p+1\cr 
 3 & 6+2p \cr 
-1 & -8} \right)
\]}
{\(A\) má hodnost \(2\) pro každé \(p\in\mathbb{R}\).}{\(p=7\)}{\(p=9\)}{\(A\) nemá hodnost \(2\) pro žádné \(p\in\mathbb{R}\).}{}{}}
\LANGsk{
\Question{
Určte \(p\) tak, aby matica \(A\) mala hodnosť \(2\).
\[
A=\left(\MATRIX{
1& p+1\cr
  3& 6+2p \cr
-1 & -8} \right)
\]}
{\(A\) má hodnosť \(2\) pre každé reálne číslo \(p\).}{\(p=7\)}{\(p=9\)}{\(A\) nemá hodnosť \(2\) pre žiadne \(p\).}{}{}}
\LANGes{
\Question{
Halla \(p\) para que la matriz \(A\) tenga el rango \(2\).
\[ 
A=\left(\MATRIX{ 
1& p+1\cr 
 3 & 6+2p \cr 
-1 & -8} \right)
\]}
{\(A\) tiene el rango \(2\) para cada número real \(p\).}{\(p=7\)}{\(p=9\)}{\(A\) no tiene el rango \(2\) para cualquier \(p\).}{}{}}
\End

\Data\ID{2000018905}
\Updated{2024-03-03 09:03:23}
\Level{B}
\SubArea{Matrices and determinants}
\LANGen{
\Question{
What is the rank of the following matrix?
\[ 
\left (\MATRIX{ 
2& 6& 10\cr 
 3 & 9& 15\cr 
7 & 0 & 1 \cr
10 & 9 & 16} \right )
\]}
{\(2\)}{\(1\)}{\(3\)}{\(4\)}{}{}}
\LANGpl{
\Question{
Jaka jest ranga poniższej macierzy?
\[ 
\left (\MATRIX{ 
2& 6& 10\cr 
 3 & 9& 15\cr 
7 & 0 & 1 \cr
10 & 9 & 16} \right )
\]}
{\(2\)}{\(1\)}{\(3\)}{\(4\)}{}{}}
\LANGcs{
\Question{
Jaká je hodnost následující matice? 
\[ 
\left (\MATRIX{ 
2& 6& 10\cr 
 3 & 9& 15\cr 
7 & 0 & 1 \cr
10 & 9 & 16} \right )
\]}
{\(2\)}{\(1\)}{\(3\)}{\(4\)}{}{}}
\LANGsk{
\Question{
Aká je hodnosť nasledujúcej matice?
\[ 
\left (\MATRIX{ 
2& 6& 10\cr 
 3 & 9& 15\cr 
7 & 0 & 1 \cr
10 & 9 & 16} \right )
\]}
{\(2\)}{\(1\)}{\(3\)}{\(4\)}{}{}}
\LANGes{
\Question{
Halla el rango de la siguiente matriz:
\[ 
\left (\MATRIX{ 
2& 6& 10\cr 
 3 & 9& 15\cr 
7 & 0 & 1 \cr
10 & 9 & 16} \right )
\]}
{\(2\)}{\(1\)}{\(3\)}{\(4\)}{}{}}
\End

\Data\ID{2000018906}
\Updated{2024-03-03 09:03:23}
\Level{B}
\SubArea{Matrices and determinants}
\LANGen{
\Question{
Specify how the rank of the matrix \(A\) changes depending on the value of \(t\), where
\[ 
A=\left (\MATRIX{ 
3& -2& 1&-4\cr 
-6& 4& -2&8\cr 
0& t& 0&t} \right ).
\]}
{If \(t=0\), the rank is \(1\), otherwise the rank is \(2\).}{If \(t=0\), the rank is \(1\), otherwise the rank is \(3\).}{If \(t=0\), the rank is \(2\), otherwise the rank is \(1\).}{If \(t=2\), the rank is \(3\), otherwise the rank is \(1\).}{}{}}
\LANGpl{
\Question{
Określ, jak zmienia się ranga macierzy \(A\) w zależności od wartości \(t\), gdzie
\[ 
A=\left (\MATRIX{ 
3& -2& 1&-4\cr 
-6& 4& -2&8\cr 
0& t& 0&t} \right ).
\]}
{Jeśli \(t=0\), to ranga wynosi \(1\), w przeciwnym razie wynosi \(2\).}{Jeśli \(t=0\), to ranga wynosi \(1\), w przeciwnym razie wynosi \(3\).}{Jeśli \(t=0\), to ranga wynosi \(2\), w przeciwnym razie wynosi \(1\).}{Jeśli \(t=2\), to ranga wynosi \(3\), w przeciwnym razie wynosi \(1\).}{}{}}
\LANGcs{
\Question{
Popište, jak se bude měnit hodnost matice \(A\) v závislosti na \(t\), jestliže
\[ 
A=\left (\MATRIX{ 
3& -2& 1&-4\cr 
-6& 4& -2&8\cr 
0& t& 0&t} \right ).
\]}
{Pro \(t=0\) je hodnost  \(1\), pro jiné hodnoty \(t\) je hodnost \(2\).}{Pro \(t=0\) je hodnost  \(1\), pro jiné hodnoty \(t\) je hodnost \(3\).}{Pro \(t=0\) je hodnost  \(2\), pro jiné hodnoty \(t\) je hodnost \(1\).}{Pro \(t=2\) je hodnost  \(3\), pro jiné hodnoty \(t\) je hodnost \(1\).}{}{}}
\LANGsk{
\Question{
Popíšte, ako sa bude meniť hodnosť matice \(A\) v závislosti na \(t\), ak
\[ 
A=\left (\MATRIX{ 
3& -2& 1&-4\cr 
-6& 4& -2&8\cr 
0& t& 0&t} \right ).
\]}
{Pre \(t=0\), je hodnosť \(1\), pre iné hodnoty \(t\) je hodnosť \(2\).}{Pre \(t=0\), je hodnosť \(1\), pre iné hodnoty \(t\) je hodnosť \(3\).}{Pre \(t=0\), je hodnosť \(2\), pre iné hodnoty \(t\) je hodnosť \(1\).}{Pre \(t=2\), je hodnosť \(3\), pre iné hodnoty \(t\) je hodnosť \(1\).}{}{}}
\LANGes{
\Question{
Especifica cómo cambia el rango de la matriz \(A\) dependiendo del valor de \(t\), donde
\[ 
A=\left (\MATRIX{ 
3& -2& 1&-4\cr 
-6& 4& -2&8\cr 
0& t& 0&t} \right ).
\]}
{Si \(t=0\), el rango es \(1\), por otra parte el rango es \(2\).}{Si \(t=0\), el rango es \(1\), por otra parte el rango es \(3\).}{Si \(t=0\), el rango es \(2\), por otra parte el rango es \(1\).}{Si \(t=2\), el rango es \(3\), por otra parte el rango es \(1\).}{}{}}
\End

\Data\ID{9000019901}
\Updated{2024-03-03 09:03:26}
\Level{B}
\SubArea{Matrices and determinants}
\LANGen{
\Question{
Find the rank of the following matrix
\(A\).
\[ 
A = \left (\MATRIX{ 
-3& 3 & 3\cr 
 2 & -2 & 2
\cr 
-1& 1 & 1 } \right )
\]}
{\(\mathop{\mathrm{rank}}(A) = 2\)}{\(\mathop{\mathrm{rank}}(A) = 3\)}{\(\mathop{\mathrm{rank}}(A) = 1\)}{\(\mathop{\mathrm{rank}}(A) = 0\)}{}{}}
\LANGpl{
\Question{
Wyznacz rząd \( r(A) \) macierzy 
\(A\).
\[ 
A = \left (\MATRIX{ 
-3& 3 & 3\cr 
 2 & -2 & 2
\cr 
-1& 1 & 1 } \right )
\]}
{\(\mathop{\mathrm{r}}(A) = 2\)}{\(\mathop{\mathrm{r}}(A) = 3\)}{\(\mathop{\mathrm{r}}(A) = 1\)}{\(\mathop{\mathrm{r}}(A) = 0\)}{}{}}
\LANGcs{
\Question{
Určete hodnost matice \(A\).
\[ 
A = \left (\MATRIX{ 
-3& 3 & 3\cr 
 2 & -2 & 2
\cr 
-1& 1 & 1 } \right )
\]}
{\(h(A) = 2\)}{\(h(A) = 3\)}{\(h(A) = 1\)}{\(h(A) = 0\)}{}{}}
\LANGsk{
\Question{
Určte hodnosť matice
\(A\).
\[ 
A = \left (\MATRIX{ 
-3& 3 & 3\cr 
 2 & -2 & 2
\cr 
-1& 1 & 1 } \right )
\]}
{\(h(A) = 2\)}{\(h(A) = 3\)}{\(h(A) = 1\)}{\(h(A) = 0\)}{}{}}
\LANGes{
\Question{
Halla el rango de la siguiente matriz
\(A\).
\[ 
A = \left (\MATRIX{ 
-3& 3 & 3\cr 
 2 & -2 & 2
\cr 
-1& 1 & 1 } \right )
\]}
{\(\mathop{\mathrm{rango}}(A) = 2\)}{\(\mathop{\mathrm{rango}}(A) = 3\)}{\(\mathop{\mathrm{rango}}(A) = 1\)}{\(\mathop{\mathrm{rango}}(A) = 0\)}{}{}}
\End

\Data\ID{9000019902}
\Updated{2024-03-03 09:03:26}
\Level{B}
\SubArea{Matrices and determinants}
\LANGen{
\Question{
Find the rank of the following matrix
\(B\).
\[ 
B = \left (\MATRIX{ 
3& 3& -2& 0\cr 
0& 1 & 2 & 1
\cr 
0& 0& -3& 0 } \right )
\]}
{\(\mathop{\mathrm{rank}}(B) = 3\)}{\(\mathop{\mathrm{rank}}(B) = 2\)}{\(\mathop{\mathrm{rank}}(B) = 1\)}{\(\mathop{\mathrm{rank}}(B) = 0\)}{}{}}
\LANGpl{
\Question{
Wyznacz rząd \( r(B) \) macierzy
\(B\).
\[ 
B = \left (\MATRIX{ 
3& 3& -2& 0\cr 
0& 1 & 2 & 1
\cr 
0& 0& -3& 0 } \right )
\]}
{\(\mathop{\mathrm{r}}(B) = 3\)}{\(\mathop{\mathrm{r}}(B) = 2\)}{\(\mathop{\mathrm{r}}(B) = 1\)}{\(\mathop{\mathrm{r}}(B) = 0\)}{}{}}
\LANGcs{
\Question{
Určete hodnost matice \(B\).
\[ 
B = \left (\MATRIX{ 
3& 3& -2& 0\cr 
0& 1 & 2 & 1
\cr 
0& 0& -3& 0 } \right )
\]}
{\(h(B) = 3\)}{\(h(B) = 2\)}{\(h(B) = 1\)}{\(h(B) = 0\)}{}{}}
\LANGsk{
\Question{
Určte hodnosť matice 
\(B\).
\[ 
B = \left (\MATRIX{ 
3& 3& -2& 0\cr 
0& 1 & 2 & 1
\cr 
0& 0& -3& 0 } \right )
\]}
{\(h(B) = 3\)}{\(h(B) = 2\)}{\(h(B) = 1\)}{\(h(B) = 0\)}{}{}}
\LANGes{
\Question{
Halla el rango de la siguiente matriz
\(B\).
\[ 
B = \left (\MATRIX{ 
3& 3& -2& 0\cr 
0& 1 & 2 & 1
\cr 
0& 0& -3& 0 } \right )
\]}
{\(\mathop{\mathrm{rango}}(B) = 3\)}{\(\mathop{\mathrm{rango}}(B) = 2\)}{\(\mathop{\mathrm{rango}}(B) = 1\)}{\(\mathop{\mathrm{rango}}(B) = 0\)}{}{}}
\End

\Data\ID{2000017601}
\Updated{2024-03-03 09:03:25}
\Level{C}
\SubArea{Matrices and determinants}
\LANGen{
\Question{
Calculate the determinant of the following matrix:
\[ 
\left (\MATRIX{ 
3  & 0 & -1\cr 
2  & 4& 0\cr 
1  & 1 & 5\cr 
} \right )
\]}
{\( 62\)}{\( 54\)}{\( 66\)}{The result does not belong to the set \(\{54;62;66\}\).}{}{}}
\LANGpl{
\Question{
Oblicz wyznacznik następującej macierzy:
\[ 
\left (\MATRIX{ 
3  & 0 & -1\cr 
2  & 4& 0\cr 
1  & 1 & 5\cr 
} \right )
\]}
{\( 62\)}{\( 54\)}{\( 66\)}{Wynik nie należy do zbioru \(\{54;62;66\}\).}{}{}}
\LANGcs{
\Question{
Vypočítejte determinant následující matice:
\[ 
\left (\MATRIX{ 
3  & 0 & -1\cr 
2  & 4& 0\cr 
1  & 1 & 5\cr 
} \right )
\]}
{\( 62\)}{\( 54\)}{\( 66\)}{Výsledek nepatří do množiny \(\{54;62;66\}\).}{}{}}
\LANGsk{
\Question{
Vypočítajte determinant nasledujúcej matice:
\[ 
\left (\MATRIX{ 
3  & 0 & -1\cr 
2  & 4& 0\cr 
1  & 1 & 5\cr 
} \right )
\]}
{\( 62\)}{\( 54\)}{\( 66\)}{Výsledok nepatrí do množiny \(\{54;62;66\}\).}{}{}}
\LANGes{
\Question{
Calcula el determinante de la siguiente matriz:
\[ 
\left (\MATRIX{ 
3  & 0 & -1\cr 
2  & 4& 0\cr 
1  & 1 & 5\cr 
} \right )
\]}
{\( 62\)}{\( 54\)}{\( 66\)}{El resultado no pertenece al conjunto \(\{54;62;66\}\).}{}{}}
\End

\Data\ID{2000017602}
\Updated{2024-03-03 09:03:25}
\Level{C}
\SubArea{Matrices and determinants}
\LANGen{
\Question{
Calculate the determinant of the following matrix:
\[ 
\left (\MATRIX{ 
\frac12  & 0 & -1\cr 
2  & 1& 0\cr 
1  & \frac12 & 2\cr 
} \right )
\]}
{\( 1\)}{\( 3\)}{\( -1\)}{The result does not belong to the set \(\{-1;1;3\}\).}{}{}}
\LANGpl{
\Question{
Oblicz wyznacznik następującej macierzy:
\[ 
\left (\MATRIX{ 
\frac12  & 0 & -1\cr 
2  & 1& 0\cr 
1  & \frac12 & 2\cr 
} \right )
\]}
{\( 1\)}{\( 3\)}{\( -1\)}{Wynik nie należy do zbioru \(\{-1;1;3\}\).}{}{}}
\LANGcs{
\Question{
Vypočítejte determinant následující matice:
\[ 
\left (\MATRIX{ 
\frac12  & 0 & -1\cr 
2  & 1& 0\cr 
1  & \frac12 & 2\cr 
} \right )
\]}
{\( 1\)}{\( 3\)}{\( -1\)}{Řešení nepatří do množiny \(\{-1;1;3\}\).}{}{}}
\LANGsk{
\Question{
Vypočítajte determinant nasledujúcej matice:
\[ 
\left (\MATRIX{ 
\frac12  & 0 & -1\cr 
2  & 1& 0\cr 
1  & \frac12 & 2\cr 
} \right )
\]}
{\( 1\)}{\( 3\)}{\( -1\)}{Výsledok nepatrí do množiny \(\{-1;1;3\}\).}{}{}}
\LANGes{
\Question{
Calcula el determinante de la siguiente matriz:
\[ 
\left (\MATRIX{ 
\frac12  & 0 & -1\cr 
2  & 1& 0\cr 
1  & \frac12 & 2\cr 
} \right )
\]}
{\( 1\)}{\( 3\)}{\( -1\)}{El resultado no pertenece al conjunto  \(\{-1;1;3\}\).}{}{}}
\End

\Data\ID{2000017603}
\Updated{2024-03-03 09:03:25}
\Level{C}
\SubArea{Matrices and determinants}
\LANGen{
\Question{
For what value of \(a\) is the determinant of the following matrix equal to \(14\)?
\[ 
\left (\MATRIX{ 
a  & 1 & -1\cr 
0  & 1& 3\cr 
1  & 0 & 2\cr 
} \right )
\]}
{\( a=5\)}{\( a=6\)}{\( a=9\)}{The value of \(a\) does not belong to the set \(\{5;6;9\}\).}{}{}}
\LANGpl{
\Question{
Dla jakiej wartości \(a\) wyznacznik poniższej macierzy jest równy \(14\)?
\[ 
\left (\MATRIX{ 
a  & 1 & -1\cr 
0  & 1& 3\cr 
1  & 0 & 2\cr 
} \right )
\]}
{\( a=5\)}{\( a=6\)}{\( a=9\)}{Wartość \(a\) nie należy do zbioru \(\{5;6;9\}\).}{}{}}
\LANGcs{
\Question{
Pro jakou hodnotu \(a\) je determinant následující matice roven \(14\)?
\[ 
\left (\MATRIX{ 
a  & 1 & -1\cr 
0  & 1& 3\cr 
1  & 0 & 2\cr 
} \right )
\]}
{\( a=5\)}{\( a=6\)}{\( a=9\)}{Hodnota \(a\) nepatří do množiny \(\{5;6;9\}\).}{}{}}
\LANGsk{
\Question{
Pre akú hodnotu \(a\) je determinant nasledujúcej matice rovný \(14\)?
\[ 
\left (\MATRIX{ 
a  & 1 & -1\cr 
0  & 1& 3\cr 
1  & 0 & 2\cr 
} \right )
\]}
{\( a=5\)}{\( a=6\)}{\( a=9\)}{Hodnota \(a\) nepatrí do množiny \(\{5;6;9\}\).}{}{}}
\LANGes{
\Question{
¿Para qué valor de \(a\) el determinante de la siguiente matriz equivale a \(14\)?
\[ 
\left (\MATRIX{ 
a  & 1 & -1\cr 
0  & 1& 3\cr 
1  & 0 & 2\cr 
} \right )
\]}
{\( a=5\)}{\( a=6\)}{\( a=9\)}{El valor de \(a\) no pertenece al conjunto \(\{5;6;9\}\).}{}{}}
\End

\Data\ID{2000017604}
\Updated{2024-03-03 09:03:25}
\Level{C}
\SubArea{Matrices and determinants}
\LANGen{
\Question{
Which of the given matrices \(A\), \(B\), \(C\) and \(D\) has a different determinant than the others?
\[ 
A=\left (\MATRIX{ 
6 & 11 \cr 
2 & 2\cr 
} \right )
\]
\[ 
B=\left (\MATRIX{ 
1 & 3 \cr 
5 & 2\cr 
} \right )
\]
\[ 
C=\left (\MATRIX{ 
5 & -2 \cr 
10 & -6\cr 
} \right )
\]
\[ 
D=\left (\MATRIX{ 
10 & 0 \cr 
-7 & -1\cr 
} \right )
\]}
{\( B\)}{\( D\)}{All given matrices have the same determinant.}{Each matrix has a different determinant.}{}{}}
\LANGpl{
\Question{
Która z podanych macierzy \(A\), \(B\), \(C\) i \(D\) ma inny wyznacznik niż pozostałe?
\[ 
A=\left (\MATRIX{ 
6 & 11 \cr 
2 & 2\cr 
} \right )
\]
\[ 
B=\left (\MATRIX{ 
1 & 3 \cr 
5 & 2\cr 
} \right )
\]
\[ 
C=\left (\MATRIX{ 
5 & -2 \cr 
10 & -6\cr 
} \right )
\]
\[ 
D=\left (\MATRIX{ 
10 & 0 \cr 
-7 & -1\cr 
} \right )
\]}
{\( B\)}{\( D\)}{Wszystkie podane macierze mają ten sam wyznacznik.}{Każda macierz ma inny wyznacznik.}{}{}}
\LANGcs{
\Question{
Která z uvedených matic \(A\), \(B\), \(C\) a \(D\) má jiný determinant než ostatní?
\[ 
A=\left (\MATRIX{ 
6 & 11 \cr 
2 & 2\cr 
} \right )
\]
\[ 
B=\left (\MATRIX{ 
1 & 3 \cr 
5 & 2\cr 
} \right )
\]
\[ 
C=\left (\MATRIX{ 
5 & -2 \cr 
10 & -6\cr 
} \right )
\]
\[ 
D=\left (\MATRIX{ 
10 & 0 \cr 
-7 & -1\cr 
} \right )
\]}
{\( B\)}{\( D\)}{Všechny uvedené matice mají stejný determinant.}{Každá matice má jiný determinant.}{}{}}
\LANGsk{
\Question{
Ktorá z daných matíc \(A\), \(B\), \(C\) a \(D\) má iný determinant?
\[ 
A=\left (\MATRIX{ 
6 & 11 \cr 
2 & 2\cr 
} \right )
\]
\[ 
B=\left (\MATRIX{ 
1 & 3 \cr 
5 & 2\cr 
} \right )
\]
\[ 
C=\left (\MATRIX{ 
5 & -2 \cr 
10 & -6\cr 
} \right )
\]
\[ 
D=\left (\MATRIX{ 
10 & 0 \cr 
-7 & -1\cr 
} \right )
\]}
{\( B\)}{\( D\)}{Všetky matice majú rovnaký determinant.}{Každá matica má iný determinant.}{}{}}
\LANGes{
\Question{
¿Cuál de las siguientes matrices \(A\), \(B\), \(C\) y \(D\) tiene el determinante diferente de los demás?
\[ 
A=\left (\MATRIX{ 
6 & 11 \cr 
2 & 2\cr 
} \right )
\]
\[ 
B=\left (\MATRIX{ 
1 & 3 \cr 
5 & 2\cr 
} \right )
\]
\[ 
C=\left (\MATRIX{ 
5 & -2 \cr 
10 & -6\cr 
} \right )
\]
\[ 
D=\left (\MATRIX{ 
10 & 0 \cr 
-7 & -1\cr 
} \right )
\]}
{\( B\)}{\( D\)}{Todas las matrices tienen el mismo determinante.}{Cada matriz tiene un determinante diferente.}{}{}}
\End

\Data\ID{2000017605}
\Updated{2024-03-03 09:03:25}
\Level{C}
\SubArea{Matrices and determinants}
\LANGen{
\Question{
Let \(a\) and \(b\) be real numbers. What is the determinant of the following matrix equal to?
\[ 
\left (\MATRIX{ 
a +b & 1 & 0\cr 
1  & a-b& 1\cr 
a & 0 & -1\cr 
} \right )
\]}
{\( b^2-a^2+a+1\)}{\( a^2-b^2+a+1\)}{\( a^2-b^2+a-1\)}{\( b^2-a^2-a+1\)}{}{}}
\LANGpl{
\Question{
\(a\) i \(b\) to liczby rzeczywiste. Jaki jest wyznacznik poniższej macierzy?
\[ 
\left (\MATRIX{ 
a +b & 1 & 0\cr 
1  & a-b& 1\cr 
a & 0 & -1\cr 
} \right )
\]}
{\( b^2-a^2+a+1\)}{\( a^2-b^2+a+1\)}{\( a^2-b^2+a-1\)}{\( b^2-a^2-a+1\)}{}{}}
\LANGcs{
\Question{
Nechť \(a\) a \(b\) jsou reálná čísla. Čemu se rovná determinant následující matice?
\[ 
\left (\MATRIX{ 
a +b & 1 & 0\cr 
1  & a-b& 1\cr 
a & 0 & -1\cr 
} \right )
\]}
{\( b^2-a^2+a+1\)}{\( a^2-b^2+a+1\)}{\( a^2-b^2+a-1\)}{\( b^2-a^2-a+1\)}{}{}}
\LANGsk{
\Question{
Nech \(a\) a \(b\) sú reálne čísla. Čomu sa rovná determinant nasledujúcej matice?
\[ 
\left (\MATRIX{ 
a +b & 1 & 0\cr 
1  & a-b& 1\cr 
a & 0 & -1\cr 
} \right )
\]}
{\( b^2-a^2+a+1\)}{\( a^2-b^2+a+1\)}{\( a^2-b^2+a-1\)}{\( b^2-a^2-a+1\)}{}{}}
\LANGes{
\Question{
Suponiendo que \(a\) y \(b\) son números reales. ¿A qué es igual el determinante de la siguiente matriz?
\[ 
\left (\MATRIX{ 
a +b & 1 & 0\cr 
1  & a-b& 1\cr 
a & 0 & -1\cr 
} \right )
\]}
{\( b^2-a^2+a+1\)}{\( a^2-b^2+a+1\)}{\( a^2-b^2+a-1\)}{\( b^2-a^2-a+1\)}{}{}}
\End

\Data\ID{2000017606}
\Updated{2024-03-03 09:03:25}
\Level{C}
\SubArea{Matrices and determinants}
\LANGen{
\Question{
Determine the set of all real numbers \(b\) for which the determinant of the following matrix equals to \(5\).
\[ 
\left (\MATRIX{ 
4 & b & -1\cr 
3 &0& 2\cr 
b & 0 & -1\cr 
} \right )
\]}
{\( \left\{1;-\frac52\right\}\)}{\( \left\{-\frac52\right\}\)}{\( \left\{1\right\}\)}{\( \emptyset\)}{}{}}
\LANGpl{
\Question{
Wyznacz zbiór wszystkich liczb rzeczywistych \(b\), dla których wyznacznik poniższej macierzy jest równy \(5\). 
\[ 
\left (\MATRIX{ 
4 & b & -1\cr 
3 &0& 2\cr 
b & 0 & -1\cr 
} \right )
\]}
{\( \left\{1;-\frac52\right\}\)}{\( \left\{-\frac52\right\}\)}{\( \left\{1\right\}\)}{\( \emptyset\)}{}{}}
\LANGcs{
\Question{
Určete množinu všech reálných čísel \(b\), pro něž je determinant následující matice roven  \(5\).
\[ 
\left (\MATRIX{ 
4 & b & -1\cr 
3 &0& 2\cr 
b & 0 & -1\cr 
} \right )
\]}
{\( \left\{1;-\frac52\right\}\)}{\( \left\{-\frac52\right\}\)}{\( \left\{1\right\}\)}{\( \emptyset\)}{}{}}
\LANGsk{
\Question{
Určte množinu všetkých reálnych čísel \(b\), pre ktoré je determinant nasledujúcej matice rovný \(5\).
\[ 
\left (\MATRIX{ 
4 & b & -1\cr 
3 &0& 2\cr 
b & 0 & -1\cr 
} \right )
\]}
{\( \left\{1;-\frac52\right\}\)}{\( \left\{-\frac52\right\}\)}{\( \left\{1\right\}\)}{\( \emptyset\)}{}{}}
\LANGes{
\Question{
Determina el conjunto de todos los números reales \(b\) para los cuales el determinante de la siguiente matriz equivale a \(5\).
\[ 
\left (\MATRIX{ 
4 & b & -1\cr 
3 &0& 2\cr 
b & 0 & -1\cr 
} \right )
\]}
{\( \left\{1;-\frac52\right\}\)}{\( \left\{-\frac52\right\}\)}{\( \left\{1\right\}\)}{\( \emptyset\)}{}{}}
\End

\Data\ID{2000017607}
\Updated{2024-03-03 09:03:25}
\Level{C}
\SubArea{Matrices and determinants}
\LANGen{
\Question{
For which real number \(x\) is the determinant of the matrix equal to \(124\)?
\[ 
\left (\MATRIX{ 
2^2 & \sqrt{625} & \sin\left(-\frac{\pi}2\right)\cr 
\log_3 x &0& \cos 0\cr 
2 & 7^{-1} & -1\cr 
} \right )
\]}
{\( 27\)}{\( \frac1{27}\)}{\( 9\)}{\(\frac19\)}{}{}}
\LANGpl{
\Question{
Dla jakiej liczby rzeczywistej \(x\) wyznacznik macierzy jest równy \(124\)?
\[ 
\left (\MATRIX{ 
2^2 & \sqrt{625} & \sin\left(-\frac{\pi}2\right)\cr 
\log_3 x &0& \cos 0\cr 
2 & 7^{-1} & -1\cr 
} \right )
\]}
{\( 27\)}{\( \frac1{27}\)}{\( 9\)}{\(\frac19\)}{}{}}
\LANGcs{
\Question{
Pro jaké reálné číslo \(x\) je determinant matice roven \(124\)?
\[ 
\left (\MATRIX{ 
2^2 & \sqrt{625} & \sin\left(-\frac{\pi}2\right)\cr 
\log_3 x &0& \cos 0\cr 
2 & 7^{-1} & -1\cr 
} \right )
\]}
{\( 27\)}{\( \frac1{27}\)}{\( 9\)}{\(\frac19\)}{}{}}
\LANGsk{
\Question{
Pre ktoré reálne číslo \(x\) je determinant matice rovný \(124\)?
\[ 
\left (\MATRIX{ 
2^2 & \sqrt{625} & \sin\left(-\frac{\pi}2\right)\cr 
\log_3 x &0& \cos 0\cr 
2 & 7^{-1} & -1\cr 
} \right )
\]}
{\( 27\)}{\( \frac1{27}\)}{\( 9\)}{\(\frac19\)}{}{}}
\LANGes{
\Question{
¿Para qué número real \(x\) el determinante de la siguiente matriz equivale a \(124\)?
\[ 
\left (\MATRIX{ 
2^2 & \sqrt{625} & \sin\left(-\frac{\pi}2\right)\cr 
\log_3 x &0& \cos 0\cr 
2 & 7^{-1} & -1\cr 
} \right )
\]}
{\( 27\)}{\( \frac1{27}\)}{\( 9\)}{\(\frac19\)}{}{}}
\End

\Data\ID{2000017608}
\Updated{2024-03-03 09:03:25}
\Level{C}
\SubArea{Matrices and determinants}
\LANGen{
\Question{
Calculate the determinant of the following matrix:
\[ 
\left (\MATRIX{ 
7  & 6 & 1\cr 
3  & 2& 1\cr 
2  & 7& -1\cr 
} \right )
\]}
{\( -16\)}{\( 12\)}{\( 54\)}{\( -52\)}{}{}}
\LANGpl{
\Question{
Oblicz wyznacznik następującej macierzy:
\[ 
\left (\MATRIX{ 
7  & 6 & 1\cr 
3  & 2& 1\cr 
2  & 7& -1\cr 
} \right )
\]}
{\( -16\)}{\( 12\)}{\( 54\)}{\( -52\)}{}{}}
\LANGcs{
\Question{
Vypočítejte determinant následující matice:
\[ 
\left (\MATRIX{ 
7  & 6 & 1\cr 
3  & 2& 1\cr 
2  & 7& -1\cr 
} \right )
\]}
{\( -16\)}{\( 12\)}{\( 54\)}{\( -52\)}{}{}}
\LANGsk{
\Question{
Vypočítajte determinant nasledujúcej matice:
\[ 
\left (\MATRIX{ 
7  & 6 & 1\cr 
3  & 2& 1\cr 
2  & 7& -1\cr 
} \right )
\]}
{\( -16\)}{\( 12\)}{\( 54\)}{\( -52\)}{}{}}
\LANGes{
\Question{
Calcula el determinante de la siguiente matriz:
\[ 
\left (\MATRIX{ 
7  & 6 & 1\cr 
3  & 2& 1\cr 
2  & 7& -1\cr 
} \right )
\]}
{\( -16\)}{\( 12\)}{\( 54\)}{\( -52\)}{}{}}
\End

\Data\ID{2000017609}
\Updated{2024-03-03 09:03:25}
\Level{C}
\SubArea{Matrices and determinants}
\LANGen{
\Question{
Calculate the determinant of the following matrix:
\[ 
\left (\MATRIX{ 
1+\sqrt5  & 4-\sqrt3 \cr 
4+\sqrt3  & 1-\sqrt5\cr 
} \right )
\]}
{\( -17\)}{\( 9\)}{\( -23\)}{\( 17\)}{}{}}
\LANGpl{
\Question{
Oblicz wyznacznik następującej macierzy:
\[ 
\left (\MATRIX{ 
1+\sqrt5  & 4-\sqrt3 \cr 
4+\sqrt3  & 1-\sqrt5\cr 
} \right )
\]}
{\( -17\)}{\( 9\)}{\( -23\)}{\( 17\)}{}{}}
\LANGcs{
\Question{
Vypočítejte determinant následující matice:
\[ 
\left (\MATRIX{ 
1+\sqrt5  & 4-\sqrt3 \cr 
4+\sqrt3  & 1-\sqrt5\cr 
} \right )
\]}
{\( -17\)}{\( 9\)}{\( -23\)}{\( 17\)}{}{}}
\LANGsk{
\Question{
Vypočítate determinant nasledujúcej matice:
\[ 
\left (\MATRIX{ 
1+\sqrt5  & 4-\sqrt3 \cr 
4+\sqrt3  & 1-\sqrt5\cr 
} \right )
\]}
{\( -17\)}{\( 9\)}{\( -23\)}{\( 17\)}{}{}}
\LANGes{
\Question{
Calcula el determinante de la siguiente matriz:
\[ 
\left (\MATRIX{ 
1+\sqrt5  & 4-\sqrt3 \cr 
4+\sqrt3  & 1-\sqrt5\cr 
} \right )
\]}
{\( -17\)}{\( 9\)}{\( -23\)}{\( 17\)}{}{}}
\End

\Data\ID{2000017610}
\Updated{2024-03-03 09:03:23}
\Level{C}
\SubArea{Matrices and determinants}
\LANGen{
\Question{
For which positive real number \(a\) is the determinant of the matrix equal to \(36\)?
\[ 
\left (\MATRIX{ 
2a  & 0 & a\cr 
0  & 1& -1\cr 
a & 2a & 6\cr 
} \right )
\]}
{\( a=2\)}{\( a=4\)}{\( a=6\)}{\(a=0\)}{}{}}
\LANGpl{
\Question{
Dla jakiej dodatniej liczby rzeczywistej \(a\) wyznacznik macierzy jest równy \(36\)?
\[ 
\left (\MATRIX{ 
2a  & 0 & a\cr 
0  & 1& -1\cr 
a & 2a & 6\cr 
} \right )
\]}
{\( a=2\)}{\( a=4\)}{\( a=6\)}{\(a=0\)}{}{}}
\LANGcs{
\Question{
Pro které kladné reálné číslo  \(a\) je determinant následující matice roven \(36\)?
\[ 
\left (\MATRIX{ 
2a  & 0 & a\cr 
0  & 1& -1\cr 
a & 2a & 6\cr 
} \right )
\]}
{\( a=2\)}{\( a=4\)}{\( a=6\)}{\(a=0\)}{}{}}
\LANGsk{
\Question{
Pre ktoré kladné reálne číslo \(a\) je determinant matice rovný \(36\)?
\[ 
\left (\MATRIX{ 
2a  & 0 & a\cr 
0  & 1& -1\cr 
a & 2a & 6\cr 
} \right )
\]}
{\( a=2\)}{\( a=4\)}{\( a=6\)}{\(a=0\)}{}{}}
\LANGes{
\Question{
¿Para qué número real positivo \(a\) el determinante de la siguiente matriz equivale a \(36\)?
\[ 
\left (\MATRIX{ 
2a  & 0 & a\cr 
0  & 1& -1\cr 
a & 2a & 6\cr 
} \right )
\]}
{\( a=2\)}{\( a=4\)}{\( a=6\)}{\(a=0\)}{}{}}
\End

\Data\ID{2000018401}
\Updated{2024-05-10 11:05:26}
\Level{C}
\SubArea{Matrices and determinants}
\LANGen{
\Question{
Assuming that
\[
\left| 
\MATRIX{ 
1  & 1 & 1\cr 
a & b& c\cr 
x  & y & z\cr 
}
\right| =5
\]
compute the following determinant:
\[
\left| 
\MATRIX{ 
a & b & c\cr 
x & y& z\cr 
3  & 3 & 3\cr 
}
\right|. 
\]}
{\(15\)}{\(-15\)}{\(27\)}{other}{}{}}
\LANGpl{
\Question{
Zakładając, że
\[
\left| 
\MATRIX{ 
1  & 1 & 1\cr 
a & b& c\cr 
x  & y & z\cr 
}
\right| =5
\]
obliczyć następujący wyznacznik:
\[
\left| 
\MATRIX{ 
a & b & c\cr 
x & y& z\cr 
3  & 3 & 3\cr 
}
\right|. 
\]}
{\(15\)}{\(-15\)}{\(27\)}{inny}{}{}}
\LANGcs{
\Question{
Nechť
\[
\left| 
\MATRIX{ 
1  & 1 & 1\cr 
a & b& c\cr 
x  & y & z\cr 
}
\right| =5.
\]
Vypočítejte následující determinant.
\[
\left| 
\MATRIX{ 
a & b & c\cr 
x & y& z\cr 
3  & 3 & 3\cr 
}
\right| 
\]}
{\(15\)}{\(-15\)}{\(27\)}{jiný výsledek}{}{}}
\LANGsk{
\Question{
Predpokladajme, že
\[
\left| 
\MATRIX{ 
1  & 1 & 1\cr 
a & b& c\cr 
x  & y & z\cr 
}
\right| =5.
\]
Vypočítajte determinant.
\[
\left| 
\MATRIX{ 
a & b & c\cr 
x & y& z\cr 
3  & 3 & 3\cr 
}
\right| 
\]}
{\(15\)}{\(-15\)}{\(27\)}{iný výsledok}{}{}}
\LANGes{
\Question{
Suponiendo que
\[
\left| 
\MATRIX{ 
1  & 1 & 1\cr 
a & b& c\cr 
x  & y & z\cr 
}
\right| =5
\]
calcula el siguiente determinante:
\[
\left| 
\MATRIX{ 
a & b & c\cr 
x & y& z\cr 
3  & 3 & 3\cr 
}
\right|. 
\]}
{\(15\)}{\(-15\)}{\(27\)}{otro}{}{}}
\End

\Data\ID{2000018402}
\Updated{2024-03-03 09:03:23}
\Level{C}
\SubArea{Matrices and determinants}
\LANGen{
\Question{
Which two of the matrices, \(A\), \(B\), and \(C\) have the same determinant?
\[
A=\left(
\MATRIX{ 
1  & 3& 5\cr 
5 & 3& 2\cr 
1  & 5 & 3\cr 
}
\right),~
B=\left(
\MATRIX{ 
1  & 2& 5\cr 
5 & 3& 3\cr 
1  & 5 & 3\cr 
}
\right),~
C=\left(
\MATRIX{ 
1  & 5& 1\cr 
3 & 3& 5\cr 
5  & 2 & 3\cr 
}
\right) 
\]}
{\(A\) and \(C\)}{\(A\) and \(B\)}{\(B\) and \(C\)}{none}{}{}}
\LANGpl{
\Question{
Które dwie z podanych macierzy, \(A\), \(B\), i \(C\) mają ten sam wyznacznik?
\[
A=\left(
\MATRIX{ 
1  & 3& 5\cr 
5 & 3& 2\cr 
1  & 5 & 3\cr 
}
\right),~
B=\left(
\MATRIX{ 
1  & 2& 5\cr 
5 & 3& 3\cr 
1  & 5 & 3\cr 
}
\right),~
C=\left(
\MATRIX{ 
1  & 5& 1\cr 
3 & 3& 5\cr 
5  & 2 & 3\cr 
}
\right) 
\]}
{\(A\) i \(C\)}{\(A\) i \(B\)}{\(B\) i \(C\)}{żadne}{}{}}
\LANGcs{
\Question{
Které dvě z níže uvedených matic A, B, C mají stejný determinant?
\[
A=\left(
\MATRIX{ 
1  & 3& 5\cr 
5 & 3& 2\cr 
1  & 5 & 3\cr 
}
\right),~
B=\left(
\MATRIX{ 
1  & 2& 5\cr 
5 & 3& 3\cr 
1  & 5 & 3\cr 
}
\right),~
C=\left(
\MATRIX{ 
1  & 5& 1\cr 
3 & 3& 5\cr 
5  & 2 & 3\cr 
}
\right) 
\]}
{\(A\) a \(C\)}{\(A\) a \(B\)}{\(B\) a \(C\)}{žádné dvě}{}{}}
\LANGsk{
\Question{
Ktoré dve z nižšie uvedených matíc A, B, C majú rovnaký determinant?
\[
A=\left(
\MATRIX{ 
1  & 3& 5\cr 
5 & 3& 2\cr 
1  & 5 & 3\cr 
}
\right),~
B=\left(
\MATRIX{ 
1  & 2& 5\cr 
5 & 3& 3\cr 
1  & 5 & 3\cr 
}
\right),~
C=\left(
\MATRIX{ 
1  & 5& 1\cr 
3 & 3& 5\cr 
5  & 2 & 3\cr 
}
\right) 
\]}
{\(A\) a \(C\)}{\(A\) a \(B\)}{\(B\) a \(C\)}{none}{}{}}
\LANGes{
\Question{
¿Cuáles de las dos matrices \(A\), \(B\), y \(C\) tienen el mismo determinante?
\[
A=\left(
\MATRIX{ 
1  & 3& 5\cr 
5 & 3& 2\cr 
1  & 5 & 3\cr 
}
\right),~
B=\left(
\MATRIX{ 
1  & 2& 5\cr 
5 & 3& 3\cr 
1  & 5 & 3\cr 
}
\right),~
C=\left(
\MATRIX{ 
1  & 5& 1\cr 
3 & 3& 5\cr 
5  & 2 & 3\cr 
}
\right) 
\]}
{\(A\) y \(C\)}{\(A\) y \(B\)}{\(B\) y \(C\)}{ningunas}{}{}}
\End

\Data\ID{2000018403}
\Updated{2024-03-03 09:03:23}
\Level{C}
\SubArea{Matrices and determinants}
\LANGen{
\Question{
Assuming that
\[
\left|
\MATRIX{ 
1  & 1& 1\cr 
a & b& c\cr 
x & y & z\cr 
}
\right|=5
\]
compute the following determinant.
\[
\left|
\MATRIX{ 
5  & 5& 5\cr 
a+2 & b+2& c+2\cr 
\frac{x}3 & \frac{y}3 & \frac{z}3\cr 
}
\right|
\]}
{\(\frac{25}3\)}{\(\frac{5}3\)}{\(7\)}{\(\frac{7}3\)}{}{}}
\LANGpl{
\Question{
Zakładając, że
\[
\left|
\MATRIX{ 
1  & 1& 1\cr 
a & b& c\cr 
x & y & z\cr 
}
\right|=5
\]
obliczyć następujący wyznacznik.
\[
\left|
\MATRIX{ 
5  & 5& 5\cr 
a+2 & b+2& c+2\cr 
\frac{x}3 & \frac{y}3 & \frac{z}3\cr 
}
\right|
\]}
{\(\frac{25}3\)}{\(\frac{5}3\)}{\(7\)}{\(\frac{7}3\)}{}{}}
\LANGcs{
\Question{
Nechť
\[
\left|
\MATRIX{ 
1  & 1& 1\cr 
a & b& c\cr 
x & y & z\cr 
}
\right|=5.
\]
Vypočítejte následující determinant.
\[
\left|
\MATRIX{ 
5  & 5& 5\cr 
a+2 & b+2& c+2\cr 
\frac{x}3 & \frac{y}3 & \frac{z}3\cr 
}
\right|
\]}
{\(\frac{25}3\)}{\(\frac{5}3\)}{\(7\)}{\(\frac{7}3\)}{}{}}
\LANGsk{
\Question{
Predpokladajme, že
\[
\left|
\MATRIX{ 
1  & 1& 1\cr 
a & b& c\cr 
x & y & z\cr 
}
\right|=5.
\]
Vypočítajte determinant.
\[
\left|
\MATRIX{ 
5  & 5& 5\cr 
a+2 & b+2& c+2\cr 
\frac{x}3 & \frac{y}3 & \frac{z}3\cr 
}
\right|
\]}
{\(\frac{25}3\)}{\(\frac{5}3\)}{\(7\)}{\(\frac{7}3\)}{}{}}
\LANGes{
\Question{
Suponiendo que
\[
\left|
\MATRIX{ 
1  & 1& 1\cr 
a & b& c\cr 
x & y & z\cr 
}
\right|=5
\]
calcula el siguiente determinante.
\[
\left|
\MATRIX{ 
5  & 5& 5\cr 
a+2 & b+2& c+2\cr 
\frac{x}3 & \frac{y}3 & \frac{z}3\cr 
}
\right|
\]}
{\(\frac{25}3\)}{\(\frac{5}3\)}{\(7\)}{\(\frac{7}3\)}{}{}}
\End

\Data\ID{2000018404}
\Updated{2024-03-03 09:03:23}
\Level{C}
\SubArea{Matrices and determinants}
\LANGen{
\Question{
Assuming that
\[
\left|
\MATRIX{ 
4  & 0& 2\cr 
a & b& c\cr 
2 & 4 & 6\cr 
}
\right|=36
\]
compute the following determinant.
\[
\left|
\MATRIX{ 
1  & 2& 3\cr 
a & b& c\cr 
4 & 0 & 2\cr 
}
\right|
\]}
{\(-18\)}{\(18\)}{\(72\)}{\(-72\)}{}{}}
\LANGpl{
\Question{
Zakładając, że
\[
\left|
\MATRIX{ 
4  & 0& 2\cr 
a & b& c\cr 
2 & 4 & 6\cr 
}
\right|=36
\]
obliczyć następujący wyznacznik.
\[
\left|
\MATRIX{ 
1  & 2& 3\cr 
a & b& c\cr 
4 & 0 & 2\cr 
}
\right|
\]}
{\(-18\)}{\(18\)}{\(72\)}{\(-72\)}{}{}}
\LANGcs{
\Question{
Nechť
\[
\left|
\MATRIX{ 
4  & 0& 2\cr 
a & b& c\cr 
2 & 4 & 6\cr 
}
\right|=36.
\]
Vypočítejte následující determinant.
\[
\left|
\MATRIX{ 
1  & 2& 3\cr 
a & b& c\cr 
4 & 0 & 2\cr 
}
\right|
\]}
{\(-18\)}{\(18\)}{\(72\)}{\(-72\)}{}{}}
\LANGsk{
\Question{
Predpokladajme, že
\[
\left|
\MATRIX{ 
4  & 0& 2\cr 
a & b& c\cr 
2 & 4 & 6\cr 
}
\right|=36.
\]
Vypočítajte determinant.
\[
\left|
\MATRIX{ 
1  & 2& 3\cr 
a & b& c\cr 
4 & 0 & 2\cr 
}
\right|
\]}
{\(-18\)}{\(18\)}{\(72\)}{\(-72\)}{}{}}
\LANGes{
\Question{
Suponiendo que
\[
\left|
\MATRIX{ 
4  & 0& 2\cr 
a & b& c\cr 
2 & 4 & 6\cr 
}
\right|=36
\]
calcula el siguiente determinante.
\[
\left|
\MATRIX{ 
1  & 2& 3\cr 
a & b& c\cr 
4 & 0 & 2\cr 
}
\right|
\]}
{\(-18\)}{\(18\)}{\(72\)}{\(-72\)}{}{}}
\End

\Data\ID{2000018405}
\Updated{2024-03-03 09:03:23}
\Level{C}
\SubArea{Matrices and determinants}
\LANGen{
\Question{
Assuming that
\[
\left|
\MATRIX{ 
1  & a& 5\cr 
1 & b& 6\cr 
1 & c & 7\cr 
}
\right|=-4
\]
compute the following determinant:
\[
\left|
\MATRIX{ 
a  & -2& 6\cr 
b & -2& 7\cr 
c & -2 & 8\cr 
}
\right|
\]}
{\(-8\)}{\(8\)}{\(-7\)}{\(9\)}{}{}}
\LANGpl{
\Question{
Zakładając, że 
\[
\left|
\MATRIX{ 
1  & a& 5\cr 
1 & b& 6\cr 
1 & c & 7\cr 
}
\right|=-4
\]
obliczyć następujący wyznacznik:
\[
\left|
\MATRIX{ 
a  & -2& 6\cr 
b & -2& 7\cr 
c & -2 & 8\cr 
}
\right|
\]}
{\(-8\)}{\(8\)}{\(-7\)}{\(9\)}{}{}}
\LANGcs{
\Question{
Nechť
\[
\left|
\MATRIX{ 
1  & a& 5\cr 
1 & b& 6\cr 
1 & c & 7\cr 
}
\right|=-4.
\]
Vypočítejte následující determinant:
\[
\left|
\MATRIX{ 
a  & -2& 6\cr 
b & -2& 7\cr 
c & -2 & 8\cr 
}
\right|
\]}
{\(-8\)}{\(8\)}{\(-7\)}{\(9\)}{}{}}
\LANGsk{
\Question{
Predpokladajme, že
\[
\left|
\MATRIX{ 
1  & a& 5\cr 
1 & b& 6\cr 
1 & c & 7\cr 
}
\right|=-4
\]
Vypočítajte determinant:
\[
\left|
\MATRIX{ 
a  & -2& 6\cr 
b & -2& 7\cr 
c & -2 & 8\cr 
}
\right|
\]}
{\(-8\)}{\(8\)}{\(-7\)}{\(9\)}{}{}}
\LANGes{
\Question{
Suponiendo que
\[
\left|
\MATRIX{ 
1  & a& 5\cr 
1 & b& 6\cr 
1 & c & 7\cr 
}
\right|=-4
\]
calcula el siguiente determinante:
\[
\left|
\MATRIX{ 
a  & -2& 6\cr 
b & -2& 7\cr 
c & -2 & 8\cr 
}
\right|
\]}
{\(-8\)}{\(8\)}{\(-7\)}{\(9\)}{}{}}
\End

\Data\ID{2000018406}
\Updated{2024-03-03 09:03:23}
\Level{C}
\SubArea{Matrices and determinants}
\LANGen{
\Question{
Which matrix from \(A\), \(B\), \(C\) and \(D\) does not have a determinant equal to zero?
\[\]
$A=\left(
\MATRIX{ 
1  & 2& 5\cr 
1 & 3& 6\cr 
1  & 4 & 7\cr 
}
\right),$ 
$B=\left(
\MATRIX{ 
1  & 2& 3\cr 
0 & 1& -1\cr 
2  & 4 & 6\cr 
}
\right),$ 
$C=\left(
\MATRIX{ 
1  & 1& 1\cr 
2 & 3& 4\cr 
15  & 16 & 17\cr 
}
\right),$ 
$D=\left(
\MATRIX{ 
1  & 2& 5\cr 
1 & -4& -6\cr 
1  & -4 & 7\cr 
}
\right)$}
{\(D\)}{\(A\)}{\(B\)}{\( C\)}{}{}}
\LANGpl{
\Question{
Która z podanych macierzy \(A\), \(B\), \(C\) i \(D\) nie posiada wyznacznika równego zeru?
\[\]
$A=\left(
\MATRIX{ 
1  & 2& 5\cr 
1 & 3& 6\cr 
1  & 4 & 7\cr 
}
\right),$ 
$B=\left(
\MATRIX{ 
1  & 2& 3\cr 
0 & 1& -1\cr 
2  & 4 & 6\cr 
}
\right),$ 
$C=\left(
\MATRIX{ 
1  & 1& 1\cr 
2 & 3& 4\cr 
15  & 16 & 17\cr 
}
\right),$ 
$D=\left(
\MATRIX{ 
1  & 2& 5\cr 
1 & -4& -6\cr 
1  & -4 & 7\cr 
}
\right)$}
{\(D\)}{\(A\)}{\(B\)}{\( C\)}{}{}}
\LANGcs{
\Question{
Která z uvedených matic \(A\), \(B\), \(C\) a \(D\) nemá nulový determinant?
\[\]
$A=\left(
\MATRIX{ 
1  & 2& 5\cr 
1 & 3& 6\cr 
1  & 4 & 7\cr 
}
\right),$ 
$B=\left(
\MATRIX{ 
1  & 2& 3\cr 
0 & 1& -1\cr 
2  & 4 & 6\cr 
}
\right),$ 
$C=\left(
\MATRIX{ 
1  & 1& 1\cr 
2 & 3& 4\cr 
15  & 16 & 17\cr 
}
\right),$ 
$D=\left(
\MATRIX{ 
1  & 2& 5\cr 
1 & -4& -6\cr 
1  & -4 & 7\cr 
}
\right)$}
{\(D\)}{\(A\)}{\(B\)}{\( C\)}{}{}}
\LANGsk{
\Question{
Ktoré z matíc \(A\), \(B\), \(C\) a \(D\) nemajú determinant rovný nule?
\[\]
$A=\left(
\MATRIX{ 
1  & 2& 5\cr 
1 & 3& 6\cr 
1  & 4 & 7\cr 
}
\right),$ 
$B=\left(
\MATRIX{ 
1  & 2& 3\cr 
0 & 1& -1\cr 
2  & 4 & 6\cr 
}
\right),$ 
$C=\left(
\MATRIX{ 
1  & 1& 1\cr 
2 & 3& 4\cr 
15  & 16 & 17\cr 
}
\right),$ 
$D=\left(
\MATRIX{ 
1  & 2& 5\cr 
1 & -4& -6\cr 
1  & -4 & 7\cr 
}
\right)$}
{\(D\)}{\(A\)}{\(B\)}{\( C\)}{}{}}
\LANGes{
\Question{
Elige la matriz \(A\), \(B\), \(C\) o \(D\) cuyo determinante no equivale a 0.
\[\]
$A=\left(
\MATRIX{ 
1  & 2& 5\cr 
1 & 3& 6\cr 
1  & 4 & 7\cr 
}
\right),$ 
$B=\left(
\MATRIX{ 
1  & 2& 3\cr 
0 & 1& -1\cr 
2  & 4 & 6\cr 
}
\right),$ 
$C=\left(
\MATRIX{ 
1  & 1& 1\cr 
2 & 3& 4\cr 
15  & 16 & 17\cr 
}
\right),$ 
$D=\left(
\MATRIX{ 
1  & 2& 5\cr 
1 & -4& -6\cr 
1  & -4 & 7\cr 
}
\right)$}
{\(D\)}{\(A\)}{\(B\)}{\( C\)}{}{}}
\End

\Data\ID{2000018407}
\Updated{2024-03-03 09:03:23}
\Level{C}
\SubArea{Matrices and determinants}
\LANGen{
\Question{
Compute the determinant of the following matrix:
\[
A=\left(
\MATRIX{ 
1  & 0& 0\cr 
2 & 3& 0\cr 
-1  & 4 & 11\cr 
}
\right)
\]}
{\(33\)}{\(0\)}{\(34\)}{\( 15\)}{}{}}
\LANGpl{
\Question{
Oblicz wyznacznik podanej macierzy:
\[
A=\left(
\MATRIX{ 
1  & 0& 0\cr 
2 & 3& 0\cr 
-1  & 4 & 11\cr 
}
\right)
\]}
{\(33\)}{\(0\)}{\(34\)}{\( 15\)}{}{}}
\LANGcs{
\Question{
Vypočítejte determinant následující matice
\[
A=\left(
\MATRIX{ 
1  & 0& 0\cr 
2 & 3& 0\cr 
-1  & 4 & 11\cr 
}
\right)
\]}
{\(33\)}{\(0\)}{\(34\)}{\( 15\)}{}{}}
\LANGsk{
\Question{
Vypočítaj determinant nasledujúcej matice:
\[
A=\left(
\MATRIX{ 
1  & 0& 0\cr 
2 & 3& 0\cr 
-1  & 4 & 11\cr 
}
\right)
\]}
{\(33\)}{\(0\)}{\(34\)}{\( 15\)}{}{}}
\LANGes{
\Question{
Calcula el determinante de la siguiente matriz:
\[
A=\left(
\MATRIX{ 
1  & 0& 0\cr 
2 & 3& 0\cr 
-1  & 4 & 11\cr 
}
\right)
\]}
{\(33\)}{\(0\)}{\(34\)}{\( 15\)}{}{}}
\End
